% Generated mechanically from the frozen simplicial.tex target.
% Do not edit this generated file; edit the bound locale source instead.

% OK, start here.
%

\setcounter{chapter}{13}
\chapter{单纯方法}



\phantomsection
\label{simplicial-section-phantom}


\section{引言}
\label{simplicial-section-introduction}

\noindent
这是对单纯方法的一份最低限度的介绍。我们只在后文需要某项内容时，才把它加入
这里。关于这些材料，一份一般性的参考文献或许是 \cite{SimpHom}。这些方法能够
完成何种工作的例子包括 Quillen 关于同伦代数的论文，见 \cite{Quillen}，以及
Artin 与 Mazur 关于 'Etale 同伦的论文，见 \cite{ArtinMazur}。

\section{有限全序集的范畴}
\label{simplicial-section-Delta}

\noindent
范畴 $\Delta$ 定义如下：
\begin{enumerate}
\item 其对象为 $[0], [1], [2], \ldots$，其中
$[n] = \{0, 1, 2, \ldots, n\}$；
\item 态射 $[n] \to [m]$ 是相应集合之间的非减映射
$\{0, 1, 2, \ldots, n\} \to \{0, 1, 2, \ldots, m\}$。
\end{enumerate}
这里，映射 $\varphi : [n] \to [m]$ 为{\it 非减的}，按定义是指：若
$i \geq j$，则 $\varphi(i) \geq \varphi(j)$。换言之，$\Delta$ 等价于以非空有限
全序集为对象、以非减映射为态射的“大”范畴。恰有 $n + 1$ 个态射
$[0] \to [n]$，而恰有 $1$ 个态射 $[n] \to [0]$。恰有
$(n + 1)(n + 2)/2$ 个态射 $[1] \to [n]$，而恰有 $n + 2$ 个态射
$[n] \to [1]$。其余情形以此类推。

\begin{definition}
\label{simplicial-definition-face-degeneracy}
对任意整数 $n\geq 1$ 及任意 $0\leq j \leq n$，以
{\it $\delta^n_j : [n-1] \to [n]$}
表示跳过 $j$ 的保序单射。对任意整数 $n\geq 0$ 及任意
$0\leq j \leq n$，以
{\it $\sigma^n_j : [n + 1]  \to [n]$}
表示满足
$(\sigma^n_j)^{-1}(\{j\}) = \{j, j + 1\}$
的保序满射。
\end{definition}

\begin{lemma}
\label{simplicial-lemma-face-degeneracy}
$\Delta$ 中的任一态射都可写成若干态射 $\delta^n_j$ 与 $\sigma^n_j$ 的复合。
\end{lemma}

\begin{proof}
设 $\varphi : [n] \to [m]$ 为 $\Delta$ 的态射。若
$j \not \in \Im(\varphi)$，则对某个态射
$\psi : [n] \to [m - 1]$，可将 $\varphi$ 写成
$\delta^m_j \circ \psi$。若 $\varphi(j) = \varphi(j + 1)$，则对某个态射
$\psi : [n - 1] \to [m]$，可将 $\varphi$ 写成
$\psi \circ \sigma^{n - 1}_j$。上述每次替换都会减小 $n + m$，故在某一步
$\varphi$ 将同时为单射和满射，因而为恒等态射，结论随即成立。
\end{proof}

\begin{lemma}
\label{simplicial-lemma-relations-face-degeneracy}
态射 $\delta^n_j$ 与 $\sigma^n_j$ 满足下列关系。
\begin{enumerate}
\item 若 $0 \leq i < j \leq n + 1$，则
$\delta^{n + 1}_j \circ \delta^n_i =
\delta^{n + 1}_i \circ \delta^n_{j - 1}$。
换言之，下图交换：
$$
\xymatrix{
& [n] \ar[rd]^{\delta^{n + 1}_j} & \\
[n - 1] \ar[ru]^{\delta^n_i} \ar[rd]_{\delta^n_{j - 1}} & &
[n + 1] \\
& [n] \ar[ru]_{\delta^{n + 1}_i} &
}
$$
\item 若 $0 \leq i < j \leq n - 1$，则
$\sigma^{n - 1}_j \circ \delta^n_i =
\delta^{n - 1}_i \circ \sigma^{n - 2}_{j - 1}$。
换言之，下图交换：
$$
\xymatrix{
& [n] \ar[rd]^{\sigma^{n - 1}_j} & \\
[n - 1] \ar[ru]^{\delta^n_i} \ar[rd]_{\sigma^{n - 2}_{j - 1}} & &
[n - 1] \\
& [n - 2] \ar[ru]_{\delta^{n - 1}_i} &
}
$$
\item 若 $0 \leq j \leq n - 1$，则
$\sigma^{n - 1}_j \circ \delta^n_j = \text{id}_{[n - 1]}$
且
$\sigma^{n - 1}_j \circ \delta^n_{j + 1} = \text{id}_{[n - 1]}$。
换言之，下图交换：
$$
\xymatrix{
& [n] \ar[rd]^{\sigma^{n - 1}_j} & \\
[n - 1]
\ar[ru]^{\delta^n_j}
\ar[rd]_{\delta^n_{j + 1}}
\ar[rr]^{\text{id}_{[n - 1]}} & & [n - 1] \\
& [n] \ar[ru]_{\sigma^{n - 1}_j} &
}
$$
\item 若 $0 < j + 1 < i \leq n$，则
$\sigma^{n - 1}_j \circ \delta^n_i =
\delta^{n - 1}_{i - 1} \circ \sigma^{n - 2}_j$。
换言之，下图交换：
$$
\xymatrix{
& [n] \ar[rd]^{\sigma^{n - 1}_j} & \\
[n - 1] \ar[ru]^{\delta^n_i} \ar[rd]_{\sigma^{n - 2}_j} & &
[n - 1] \\
& [n - 2] \ar[ru]_{\delta^{n - 1}_{i - 1}} &
}
$$
\item 若 $0 \leq i \leq j \leq n - 1$，则
$\sigma^{n - 1}_j \circ \sigma^n_i =
\sigma^{n - 1}_i \circ \sigma^n_{j + 1}$。
换言之，下图交换：
$$
\xymatrix{
& [n] \ar[rd]^{\sigma^{n - 1}_j} & \\
[n + 1] \ar[ru]^{\sigma^n_i} \ar[rd]_{\sigma^n_{j + 1}} & &
[n - 1] \\
& [n] \ar[ru]_{\sigma^{n - 1}_i} &
}
$$
\end{enumerate}
\end{lemma}

\begin{proof}
略。
\end{proof}

\begin{lemma}
\label{simplicial-lemma-face-degeneracy-category}
范畴 $\Delta$ 是具有对象 $[n]$（$n \geq 0$）及态射
$\delta^n_j$、$\sigma^n_j$ 的泛范畴，并满足：(a) 每个态射都是这些态射的
复合；(b) 引理 \ref{simplicial-lemma-relations-face-degeneracy} 所列关系成立；
(c) 这些态射之间的任一关系都是上述关系的推论。
\end{lemma}

\begin{proof}
略。
\end{proof}




\section{单纯对象}
\label{simplicial-section-simplicial-object}

\begin{definition}
\label{simplicial-definition-simplicial-object}
设 $\mathcal{C}$ 为范畴。
\begin{enumerate}
\item $\mathcal{C}$ 的一个{\it 单纯对象 $U$} 是从 $\Delta$ 到
$\mathcal{C}$ 的反变函子 $U$，即
$$
U : \Delta^{opp} \longrightarrow \mathcal{C}
$$
\item 若 $\mathcal{C}$ 为集合范畴，则称 $U$ 为{\it 单纯集}。
\item 若 $\mathcal{C}$ 为阿贝尔群范畴，则称 $U$ 为{\it 单纯阿贝尔群}。
\item 单纯对象之间的一个{\it 态射 $U \to U'$} 是函子之间的自然变换。
\item $\mathcal{C}$ 的{\it 单纯对象范畴}记为 $\text{Simp}(\mathcal{C})$。
\end{enumerate}
\end{definition}

\noindent
这意味着存在对象 $U([0]), U([1]), U([2]), \ldots$，并且对任意非减映射
$\varphi$，即 $\varphi : [m] \to [n]$，存在态射
$U(\varphi) : U([n]) \to U([m])$，满足
$U(\varphi \circ \psi) = U(\psi) \circ U(\varphi)$。

\medskip\noindent
特别地，存在唯一态射 $U([0]) \to U([n])$，并且与 $n + 1$ 个映射
$[0] \to [n]$ 相对应，恰有 $n + 1$ 个态射 $U([n]) \to U([0])$。
显然，为了清楚地讨论这些单纯对象，还需要更多记号。为此我们使用上面
第 \ref{simplicial-section-Delta} 节中挑出的那些态射。

\begin{lemma}
\label{simplicial-lemma-characterize-simplicial-object}
设 $\mathcal{C}$ 为范畴。
\begin{enumerate}
\item 给定 $\mathcal{C}$ 中的单纯对象 $U$，得到对象列
$U_n = U([n])$，以及态射
$d^n_j = U(\delta^n_j) : U_n \to U_{n-1}$ 和
$s^n_j = U(\sigma^n_j) : U_n \to U_{n + 1}$。这些态射满足引理
\ref{simplicial-lemma-relations-face-degeneracy} 所列关系的反向关系，即
\begin{enumerate}
\item 若 $0 \leq i < j \leq n + 1$，则
$d^n_i \circ d^{n + 1}_j = d^n_{j - 1} \circ d^{n + 1}_i$。
\item 若 $0 \leq i < j \leq n - 1$，则
$d^n_i \circ s^{n - 1}_j = s^{n - 2}_{j - 1} \circ d^{n - 1}_i$。
\item 若 $0 \leq j \leq n - 1$，则
$\text{id} = d^n_j \circ s^{n - 1}_j = d^n_{j + 1} \circ s^{n - 1}_j$。
\item 若 $0 < j + 1 < i \leq n$，则
$d^n_i \circ s^{n - 1}_j = s^{n - 2}_j \circ d^{n - 1}_{i - 1}$。
\item 若 $0 \leq i \leq j \leq n - 1$，则
$s^n_i \circ s^{n - 1}_j = s^n_{j + 1} \circ s^{n - 1}_i$。
\end{enumerate}
\item 反过来，给定对象列 $U_n$ 及满足 (1)(a)--(e) 的态射
$d^n_j$、$s^n_j$，则存在唯一的 $\mathcal{C}$ 中的单纯对象 $U$，使得
$U_n = U([n])$、$d^n_j = U(\delta^n_j)$ 且
$s^n_j = U(\sigma^n_j)$。
\item 单纯对象 $U$ 与 $U'$ 之间的态射由一族与态射 $d^n_j$ 和
$s^n_j$ 交换的态射 $U_n \to U'_n$ 给出。
\end{enumerate}
\end{lemma}

\begin{proof}
这由引理 \ref{simplicial-lemma-face-degeneracy-category} 得出。
\end{proof}

\begin{remark}
\label{simplicial-remark-relations}
稍微滥用记号，我们有时以 $d_i : U_n \to U_{n - 1}$ 代替
$d^n_i$，并对 $s_i : U_n \to U_{n + 1}$ 使用类似记法。
态射 $d^n_i$ 与 $s^n_i$ 之间的关系可表为：
\begin{enumerate}
\item 若 $i < j$，则 $d_i \circ d_j = d_{j - 1} \circ d_i$。
\item 若 $i < j$，则 $d_i \circ s_j = s_{j - 1} \circ d_i$。
\item 有 $\text{id} = d_j \circ s_j = d_{j + 1} \circ s_j$。
\item 若 $i > j + 1$，则 $d_i \circ s_j = s_j \circ d_{i - 1}$。
\item 若 $i \leq j$，则 $s_i \circ s_j = s_{j + 1} \circ s_i$。
\end{enumerate}
这表示，只要等式左右两侧的复合都有定义，相应等式就成立。
\end{remark}

\noindent
存在唯一态射 $s^0_0 = U(\sigma^0_0) : U_0 \to U_1$，以及两个态射
$d^1_0 = U(\delta^1_0)$ 和 $d^1_1 = U(\delta^1_1)$，它们都是
$U_1 \to U_0$ 的态射。存在两个态射
$s^1_0 = U(\sigma^1_0)$、$s^1_1 = U(\sigma^1_1)$，它们都是
$U_1 \to U_2$ 的态射。还存在三个态射
$d^2_0 = U(\delta^2_0)$、$d^2_1 = U(\delta^2_1)$、$d^2_2 = U(\delta^2_2)$，
它们都是 $U_3 \to U_2$ 的态射。依此类推。
% SOURCE-ERRATUM-CANDIDATE P02-E0102: frozen source says d^2_i: U_3 -> U_2, but its own definition gives d^2_i: U_2 -> U_1.

\medskip\noindent
在图示上，我们把 $U$ 看作：
$$
\xymatrix{
U_2
\ar@<2ex>[r]
\ar@<0ex>[r]
\ar@<-2ex>[r]
&
U_1
\ar@<1ex>[r]
\ar@<-1ex>[r]
\ar@<1ex>[l]
\ar@<-1ex>[l]
&
U_0
\ar@<0ex>[l]
}
$$
这里，指向右侧的箭头是 $d$-态射，指向左侧的箭头是 $s$-态射。

\begin{example}
\label{simplicial-example-constant-simplicial-object}
最简单的例子是取值为 $X \in \Ob(\mathcal{C})$ 的{\it 常值}单纯对象。
换言之，$U_n = X$，且所有映射均为 $\text{id}_X$。
\end{example}

\begin{example}
\label{simplicial-example-fibre-products-simplicial-object}
设 $Y\to X$ 为 $\mathcal{C}$ 的态射，并且所有纤维积
$Y \times_X Y \times_X \ldots \times_X Y$ 都存在。令 $U_n$ 等于
$(n + 1)$ 重纤维积，并使 $\varphi : [n] \to [m]$ 对应于如下映射
（用“坐标”表示）：
$(y_0, \ldots, y_m) \mapsto (y_{\varphi(0)}, \ldots, y_{\varphi(n)})$。
换言之，映射 $U_0 = Y \to U_1 = Y \times_X Y$ 是对角映射，而两个映射
$U_1 = Y \times_X Y \to U_0 = Y$ 是投影映射。
\end{example}

\noindent
从几何上说，上面的例 \ref{simplicial-example-fibre-products-simplicial-object}
十分重要。它表明，最好把映射 $d^n_j : U_n \to U_{n - 1}$ 看成投影映射
（忘掉第 $j$ 个分量），并把映射 $s^n_j : U_n \to U_{n + 1}$ 看成对角映射
（重复第 $j$ 个坐标）。后面各节还会回到这一点。

\begin{lemma}
\label{simplicial-lemma-si-injective}
设 $\mathcal{C}$ 为范畴，$U$ 为 $\mathcal{C}$ 的单纯对象。每个态射
$s^n_i : U_n \to U_{n + 1}$ 都有左逆。特别地，$s^n_i$ 是单态射。
\end{lemma}

\begin{proof}
这是因为 $d_i^{n + 1} \circ s^n_i = \text{id}_{U_n}$。
\end{proof}

\section{作为预层的单纯对象}
\label{simplicial-section-simplicial-presheaves}

\noindent
另一个观察是，可以把 $\mathcal{C}$ 的单纯对象看作 $\Delta$ 上取值于
$\mathcal{C}$ 的预层。见 Sites，定义 \ref{sites-definition-presheaf}。
事实上，若 $U$、$U'$ 是 $\mathcal{C}$ 的单纯对象，则
\begin{equation}
\label{simplicial-equation-simplicial-set-presheaf}
\Mor(U, U') = \Mor_{\textit{PSh}(\Delta)}(U, U').
\end{equation}
下面的一些材料可以用位点一章中更一般的构造来替代。不过，采用单纯对象所特有的
论证似乎能给出更清楚的图景。















\section{余单纯对象}
\label{simplicial-section-cosimplicial-object}

\noindent
范畴 $\mathcal{C}$ 的余单纯对象当然可以直接定义为对偶范畴
$\mathcal{C}^{opp}$ 的单纯对象。不过，人的思维通常并不这样运作，
因此我们在这里单独引入余单纯对象，并指出它们的一些简单性质。

\begin{definition}
\label{simplicial-definition-cosimplicial-object}
设 $\mathcal{C}$ 为范畴。
\begin{enumerate}
\item $\mathcal{C}$ 的一个{\it 余单纯对象 $U$} 是从 $\Delta$ 到
$\mathcal{C}$ 的协变函子 $U$，即
$$
U : \Delta \longrightarrow \mathcal{C}
$$
\item 若 $\mathcal{C}$ 是集合范畴，则称 $U$ 为{\it 余单纯集}。
\item 若 $\mathcal{C}$ 是阿贝尔群范畴，则称 $U$ 为{\it 余单纯阿贝尔群}。
\item {\it 余单纯对象之间的态射 $U \to U'$} 是函子之间的变换。
\item $\mathcal{C}$ 的{\it 余单纯对象范畴}记作
$\text{CoSimp}(\mathcal{C})$。
\end{enumerate}
\end{definition}

\noindent
这就是说，我们有对象 $U([0]), U([1]), U([2]), \ldots$；对于任意非降映射
$\varphi$，即 $\varphi : [m] \to [n]$，都有一个态射
$U(\varphi) : U([m]) \to U([n])$，并且满足
$U(\varphi \circ \psi) = U(\varphi) \circ U(\psi)$。

\medskip\noindent
特别地，有唯一的态射 $U([n]) \to U([0])$，并且恰有 $n + 1$ 个态射
$U([0]) \to U([n])$，它们对应于 $n + 1$ 个映射 $[0] \to [n]$。
显然，要恰当地讨论这些单纯对象，还需要更多记号。
% P02-E0103: frozen source says "these simplicial objects" in the cosimplicial section.
为此，我们考察上面第 \ref{simplicial-section-Delta} 节中挑出的那些态射。

\begin{lemma}
\label{simplicial-lemma-characterize-cosimplicial-object}
设 $\mathcal{C}$ 为范畴。
\begin{enumerate}
\item 给定 $\mathcal{C}$ 中的余单纯对象 $U$，我们得到一列对象
$U_n = U([n])$，并带有态射
$\delta^n_j = U(\delta^n_j) : U_{n - 1} \to U_n$ 和
$\sigma^n_j = U(\sigma^n_j) : U_{n + 1} \to U_n$。这些态射满足
引理 \ref{simplicial-lemma-relations-face-degeneracy} 中列出的关系。
\item 反之，给定一列对象 $U_n$ 和满足这些关系的态射
$\delta^n_j$、$\sigma^n_j$，存在唯一的 $\mathcal{C}$ 中的余单纯对象 $U$，
使得 $U_n = U([n])$、$\delta^n_j = U(\delta^n_j)$ 且
$\sigma^n_j = U(\sigma^n_j)$。
\item 余单纯对象 $U$ 与 $U'$ 之间的态射由一族态射
$U_n \to U'_n$ 给出，并且它们与态射 $\delta^n_j$ 和 $\sigma^n_j$ 可交换。
\end{enumerate}
\end{lemma}

\begin{proof}
这由引理 \ref{simplicial-lemma-face-degeneracy-category} 推出。
\end{proof}

\begin{remark}
\label{simplicial-remark-relations-cosimplicial}
我们有时滥用记号，以 $\delta_i : U_{n - 1} \to U_n$ 代替
$\delta^n_i$，并对 $\sigma_i : U_{n + 1} \to U_n$ 使用类似记法。
态射 $\delta^n_i$ 与 $\sigma^n_i$ 之间的关系可以表述如下：
\begin{enumerate}
\item 若 $i < j$，则
$\delta_j \circ \delta_i = \delta_i \circ \delta_{j - 1}$。
\item 若 $i < j$，则
$\sigma_j \circ \delta_i = \delta_i \circ \sigma_{j - 1}$。
\item 有
$\text{id} = \sigma_j \circ \delta_j = \sigma_j \circ \delta_{j + 1}$。
\item 若 $i > j + 1$，则
$\sigma_j \circ \delta_i = \delta_{i - 1} \circ \sigma_j$。
\item 若 $i \leq j$，则
$\sigma_j \circ \sigma_i = \sigma_i \circ \sigma_{j + 1}$。
\end{enumerate}
这里的意思是，只要等式两边的复合都有定义，相应等式就应成立。
\end{remark}

\noindent
我们得到唯一的态射 $\sigma^0_0 = U(\sigma^0_0) : U_1 \to U_0$，以及
两个态射 $\delta^1_0 = U(\delta^1_0)$ 和
$\delta^1_1 = U(\delta^1_1)$，它们都是态射 $U_0 \to U_1$。
有两个态射 $\sigma^1_0 = U(\sigma^1_0)$、
$\sigma^1_1 = U(\sigma^1_1)$，它们都是态射 $U_2 \to U_1$。
还有三个态射 $\delta^2_0 = U(\delta^2_0)$、
$\delta^2_1 = U(\delta^2_1)$、$\delta^2_2 = U(\delta^2_2)$，
它们都是态射 $U_2 \to U_3$。依此类推。

\medskip\noindent
在图示上，我们把 $U$ 看作：
$$
\xymatrix{
U_0
\ar@<1ex>[r]
\ar@<-1ex>[r]
&
U_1
\ar@<0ex>[l]
\ar@<2ex>[r]
\ar@<0ex>[r]
\ar@<-2ex>[r]
&
U_2
\ar@<1ex>[l]
\ar@<-1ex>[l]
}
$$
这里，指向右侧的箭头是 $\delta$-态射，指向左侧的箭头是
$\sigma$-态射。

\begin{example}
\label{simplicial-example-constant-cosimplicial-object}
最简单的例子是取值为 $X \in \Ob(\mathcal{C})$ 的{\it 常值}余单纯对象。
换言之，$U_n = X$，且所有映射均为 $\text{id}_X$。
\end{example}

\begin{example}
\label{simplicial-example-push-outs-simplicial-object}
设 $X\to Y$ 为 $C$ 的态射，并且所有推出
$Y\amalg_X Y \amalg_X \ldots \amalg_X Y$ 都存在。
% P02-E0104: frozen source uses plain C although the ambient category is mathcal C.
令 $U_n$ 等于 $(n + 1)$ 重推出，并使 $\varphi : [n] \to [m]$ 对应于
“坐标”上的映射
$$
(y \text{ 于 }i\text{次分量})
\mapsto
(y \text{ 于 }\varphi(i)\text{次分量})
$$
换言之，映射 $U_1 = Y \amalg_X Y \to U_0 = Y$ 在每个分量上都是恒等映射。
两个映射 $U_0 = Y \to U_1 = Y \amalg_X Y$ 是两个余投影。
\end{example}

\begin{example}
\label{simplicial-example-simplex-cosimplicial-set}
对每个 $n \geq 0$，以 $C[n]$ 表示余单纯集
$$
\Delta \longrightarrow \textit{Sets},\quad
[k] \longmapsto \Mor_{\Delta}([n], [k])
$$
这个例子与例 \ref{simplicial-example-simplex-simplicial-set} 对偶。
\end{example}

\begin{lemma}
\label{simplicial-lemma-di-injective}
设 $\mathcal{C}$ 为范畴，$U$ 为 $\mathcal{C}$ 的余单纯对象。每个态射
$\delta^n_i : U_{n - 1} \to U_n$ 都有左逆。特别地，$\delta^n_i$ 是单态射。
\end{lemma}

\begin{proof}
这是因为当 $j < n$ 时，
$\sigma_i^{n - 1} \circ \delta^n_i = \text{id}_{U_n}$。
% P02-E0105: frozen source has a type/index mismatch: id_{U_n} and j<n.
\end{proof}













\section{单纯对象的积}
\label{simplicial-section-products}

\noindent
当然，我们应当把单纯对象的积定义为单纯对象范畴中的积。
这可能导致一种容易混淆的情形：积虽然存在，却不能按下面的方式描述。
为避免这种情形，我们直接作如下定义。

\begin{definition}
\label{simplicial-definition-product}
设 $\mathcal{C}$ 为范畴，$U$ 和 $V$ 为 $\mathcal{C}$ 的单纯对象。
假设积 $U_n \times V_n$ 在 $\mathcal{C}$ 中存在。
$U$ 与 $V$ 的{\it 积}是如下定义的单纯对象 $U \times V$：
\begin{enumerate}
\item $(U \times V)_n = U_n \times V_n$，
\item $d^n_i = (d^n_i, d^n_i)$，且
\item $s^n_i = (s^n_i, s^n_i)$。
\end{enumerate}
换言之，$U \times V$ 是 $\Delta$ 上预层 $U$ 与 $V$ 的积。
\end{definition}

\begin{lemma}
\label{simplicial-lemma-product}
若 $U$ 和 $V$ 是范畴 $\mathcal{C}$ 中的单纯对象，且 $U \times V$ 存在，
则对 $\mathcal{C}$ 的任意第三个单纯对象 $W$，都有
$$
\Mor(W, U \times V) =
\Mor(W, U) \times
\Mor(W, V)
$$
\end{lemma}

\begin{proof}
略。
\end{proof}


\section{单纯对象的纤维积}
\label{simplicial-section-fibre-products}

\noindent
当然，我们应当把单纯对象的纤维积定义为单纯对象范畴中的纤维积。
这可能导致一种容易混淆的情形：纤维积虽然存在，却不能按下面的方式描述。
为避免这种情形，我们直接作如下定义。

\begin{definition}
\label{simplicial-definition-fibre-product}
设 $\mathcal{C}$ 为范畴，$U, V, W$ 为 $\mathcal{C}$ 的单纯对象，并设
$a : V \to U$、$b : W \to U$ 为态射。假设纤维积
$V_n \times_{U_n} W_n$ 在 $\mathcal{C}$ 中存在。
$V$ 与 $W$ 在 $U$ 上的{\it 纤维积}是如下定义的单纯对象
$V \times_U W$：
\begin{enumerate}
\item $(V \times_U W)_n = V_n \times_{U_n} W_n$，
\item $d^n_i = (d^n_i, d^n_i)$，且
\item $s^n_i = (s^n_i, s^n_i)$。
\end{enumerate}
换言之，$V \times_U W$ 是 $\Delta$ 上预层 $V$ 与 $W$
在预层 $U$ 上的纤维积。
\end{definition}

\begin{lemma}
\label{simplicial-lemma-fibre-product}
若 $U, V, W$ 是范畴 $\mathcal{C}$ 中的单纯对象，
$a : V \to U$、$b : W \to U$ 是态射，并且 $V \times_U W$ 存在，
则对 $\mathcal{C}$ 的任意第四个单纯对象 $T$，都有
$$
\Mor(T, V \times_U W) =
\Mor(T, V) \times_{\Mor(T, U)}
\Mor(T, W)
$$
\end{lemma}

\begin{proof}
略。
\end{proof}

\section{单纯对象的推出}
\label{simplicial-section-push-outs}

\noindent
当然，我们应当把单纯对象的推出定义为单纯对象范畴中的推出。
这可能导致一种容易混淆的情形：推出虽然存在，却不能按下面的方式描述。
为避免这种情形，我们直接作如下定义。

\begin{definition}
\label{simplicial-definition-push-out}
设 $\mathcal{C}$ 为范畴，$U, V, W$ 为 $\mathcal{C}$ 的单纯对象，并设
$a : U \to V$、$b : U \to W$ 为态射。假设推出
$V_n \amalg_{U_n} W_n$ 在 $\mathcal{C}$ 中存在。
$V$ 与 $W$ 在 $U$ 上的{\it 推出}是如下定义的单纯对象
$V\amalg_U W$：
\begin{enumerate}
\item $(V \amalg_U W)_n = V_n \amalg_{U_n} W_n$，
\item $d^n_i = (d^n_i, d^n_i)$，且
\item $s^n_i = (s^n_i, s^n_i)$。
\end{enumerate}
换言之，$V\amalg_U W$ 是 $\Delta$ 上预层 $V$ 与 $W$
在预层 $U$ 上的推出。
\end{definition}

\begin{lemma}
\label{simplicial-lemma-push-out}
若 $U, V, W$ 是范畴 $\mathcal{C}$ 中的单纯对象，
$a : U \to V$、$b : U \to W$ 是态射，并且 $V\amalg_U W$ 存在，
则对 $\mathcal{C}$ 的任意第四个单纯对象 $T$，都有
$$
\Mor(V\amalg_U W, T) =
\Mor(V, T) \times_{\Mor(U, T)}
\Mor(W, T)
$$
\end{lemma}

\begin{proof}
略。
\end{proof}











\section{余单纯对象的积}
\label{simplicial-section-products-cosimplicial}

\noindent
当然，我们应当把余单纯对象的积定义为余单纯对象范畴中的积。
这可能导致一种容易混淆的情形：积虽然存在，却不能按下面的方式描述。
为避免这种情形，我们直接作如下定义。

\begin{definition}
\label{simplicial-definition-product-cosimplicial-objects}
设 $\mathcal{C}$ 为范畴，$U$ 和 $V$ 为 $\mathcal{C}$ 的余单纯对象。
假设积 $U_n \times V_n$ 在 $\mathcal{C}$ 中存在。
$U$ 与 $V$ 的{\it 积}是如下定义的余单纯对象 $U \times V$：
\begin{enumerate}
\item $(U \times V)_n = U_n \times V_n$，
\item 对任意 $\varphi : [n] \to [m]$，映射
$(U \times V)(\varphi) : U_n \times V_n \to U_m \times V_m$
是积 $U(\varphi) \times V(\varphi)$。
\end{enumerate}
\end{definition}

\begin{lemma}
\label{simplicial-lemma-product-cosimplicial-objects}
若 $U$ 和 $V$ 是范畴 $\mathcal{C}$ 中的余单纯对象，且 $U \times V$ 存在，
则对 $\mathcal{C}$ 的任意第三个余单纯对象 $W$，都有
$$
\Mor(W, U \times V) =
\Mor(W, U) \times
\Mor(W, V)
$$
\end{lemma}

\begin{proof}
略。
\end{proof}

\section{余单纯对象的纤维积}
\label{simplicial-section-fibre-products-cosimplicial}

\noindent
当然，我们应当把余单纯对象的纤维积定义为余单纯对象范畴中的纤维积。
这可能导致一种容易混淆的情形：积虽然存在，却不能按下面的方式描述。
为避免这种情形，我们直接作如下定义。

\begin{definition}
\label{simplicial-definition-fibre-product-cosimplicial-objects}
设 $\mathcal{C}$ 为范畴，$U, V, W$ 为 $\mathcal{C}$ 的余单纯对象，并设
$a : V \to U$ 和 $b : W \to U$ 为态射。假设纤维积
$V_n \times_{U_n} W_n$ 在 $\mathcal{C}$ 中存在。
$V$ 与 $W$ 在 $U$ 上的{\it 纤维积}是如下定义的余单纯对象
$V \times_U W$：
\begin{enumerate}
\item $(V \times_U W)_n = V_n \times_{U_n} W_n$，
\item 对任意 $\varphi : [n] \to [m]$，映射
$(V \times_U W)(\varphi) : V_n \times_{U_n} W_n \to V_m \times_{U_m} W_m$
是积 $V(\varphi) \times_{U(\varphi)} W(\varphi)$。
\end{enumerate}
\end{definition}

\begin{lemma}
\label{simplicial-lemma-fibre-product-cosimplicial-objects}
若 $U, V, W$ 是范畴 $\mathcal{C}$ 中的余单纯对象，
$a : V \to U$、$b : W \to U$ 是态射，并且 $V \times_U W$ 存在，
则对 $\mathcal{C}$ 的任意第四个余单纯对象 $T$，都有
$$
\Mor(T, V \times_U W) =
\Mor(T, V) \times_{\Mor(T, U)}
\Mor(T, W)
$$
\end{lemma}

\begin{proof}
略。
\end{proof}













\section{单纯集}
\label{simplicial-section-simplicial-set}

\noindent
设 $U$ 为单纯集。可以把 $U_0$ 看作{\it $0$-单形}的集合，
把 $U_1$ 看作{\it $1$-单形}的集合，把 $U_2$ 看作
{\it $2$-单形}的集合，依此类推。

\medskip\noindent
我们把映射 $s^n_j : U_n \to U_{n + 1}$ 看作这样的映射：它把
$n$-单形 $A$ 送到退化的 $(n + 1)$-单形 $B$，其中该单形的
$(j, j + 1)$-边被收缩到 $A$ 的第 $j$ 个顶点。我们把映射
$d^n_j : U_n \to U_{n - 1}$ 看作把 $n$-单形 $A$ 送到它的一个面，
即略去第 $j$ 个顶点所得的面。这样一来，就可以从几何上理解
映射 $s^n_j$ 与 $d^n_j$ 之间的关系。

\begin{definition}
\label{simplicial-definition-terminology-simplicial-sets}
设 $U$ 为单纯集。称 $x$ 是 $U$ 的一个{\it $n$-单形}，是指
$x$ 是 $U_n$ 的元素。称 $y$ 是 $x$ 的第 $j$ 个{\it 面}，是指
$d^n_jx = y$。若 $z = s^n_jx$，则称 $z$ 是 $x$ 的第 $j$ 个
{\it 退化}。若一个单形是另一单形的退化，则称它为{\it 退化的}。
\end{definition}

\noindent
下面给出几个基本例子。

\begin{example}
\label{simplicial-example-simplex-simplicial-set}
对每个 $n \geq 0$，以 $\Delta[n]$ 表示单纯集
$$
\Delta^{opp} \longrightarrow \textit{Sets},\quad
[k] \longmapsto \Mor_{\Delta}([k], [n])
$$
下列断言留给读者验证。若 $m > n$，则 $\Delta[n]$ 的每个
$m$-单形都是退化的。$\Delta[n]$ 有唯一的非退化 $n$-单形，
即 $\text{id}_{[n]}$。
\end{example}

\begin{lemma}
\label{simplicial-lemma-simplex-map}
设 $U$ 为单纯集，$n \geq 0$ 为整数。存在典范双射
$$
\Mor(\Delta[n], U)
\longrightarrow
U_n
$$
它把态射 $\varphi$ 送到 $\varphi$ 在 $\Delta[n]$ 的唯一非退化
$n$-单形上的取值。
\end{lemma}

\begin{proof}
略。
\end{proof}

\begin{example}
\label{simplicial-example-simplex-category}
考虑 $\Delta$ 中 $[n]$ 上的对象所成的范畴 $\Delta/[n]$，见
Categories，例 \ref{categories-example-category-over-X}。
有一个函子 $p : \Delta/[n] \to \Delta$。$p$ 在 $[k]$ 上的纤维范畴
（见 Categories，第 \ref{categories-section-fibred-groupoids} 节）
以 $\Delta[n]$ 中 $k$-单形的集合 $\Delta[n]_k$ 为对象，并且只有恒等态射。
对 $\Delta$ 的每个态射 $\varphi : [k] \to [l]$ 以及 $[l]$ 上纤维范畴中的
每个对象 $\psi : [l] \to [n]$，在 $[k]$ 上恰有一个对象带有覆盖
$\varphi$ 的态射，即 $\psi \circ \varphi : [k] \to [n]$。
因此，$\Delta/[n]$ 是 $\Delta$ 上以集合为纤维的范畴。换言之，
可把 $\Delta/[n]$ 看成 $\Delta$ 上的集合预层。另见 Categories，
例 \ref{categories-example-fibred-category-from-functor-of-points}。
这个集合预层与单纯集 $\Delta[n]$ 一致。特别地，由公式
(\ref{simplicial-equation-simplicial-set-presheaf}) 和上面的引理
\ref{simplicial-lemma-simplex-map}，对任意单纯集 $U$，得到公式
$$
\Mor_{\textit{PSh}(\Delta)}(\Delta/[n], U) = U_n
$$
\end{example}

\begin{lemma}
\label{simplicial-lemma-product-degenerate}
设 $U$、$V$ 为单纯集，$a, b \geq 0$ 为整数。假设当 $n > a$ 时，
$U$ 的每个 $n$-单形都是退化的；并假设当 $n > b$ 时，$V$ 的每个
$n$-单形都是退化的。则当 $n > a + b$ 时，$U \times V$ 的每个
$n$-单形都是退化的。
\end{lemma}

\begin{proof}
设 $n > a + b$，并取
$(u, v) \in (U \times V)_n = U_n \times V_n$。由假设，存在
$\alpha : [n] \to [a]$、$u' \in U_a$、$\beta : [n] \to [b]$
和 $v' \in V_b$，使得 $u = U(\alpha)(u')$ 且
$v = V(\beta)(v')$。由于 $n > a + b$，存在
$0 \leq i \leq a + b$，使得 $\alpha(i) = \alpha(i + 1)$ 且
$\beta(i) = \beta(i + 1)$。立即可知，$(u, v)$ 属于
$s^{n - 1}_i$ 的像。
\end{proof}



\section{截断单纯对象与骨架函子}
\label{simplicial-section-skeleton}

\noindent
以 $\Delta_{\leq n}$ 表示 $\Delta$ 的全子范畴，其对象为
$[0], [1], [2], \ldots, [n]$。设 $\mathcal{C}$ 为范畴。

\begin{definition}
\label{simplicial-definition-truncated-simplicial-object}
$\mathcal{C}$ 的一个{\it $n$-截断单纯对象}是从
$\Delta_{\leq n}$ 到 $\mathcal{C}$ 的反变函子。
{\it $n$-截断单纯对象之间的态射}是函子之间的变换。
$\mathcal{C}$ 的 $n$-截断单纯对象范畴记作
$\text{Simp}_n(\mathcal{C})$。
\end{definition}

\noindent
给定 $\mathcal{C}$ 的单纯对象 $U$，截断 $\text{sk}_n U$ 是
$U$ 在子范畴 $\Delta_{\leq n}$ 上的限制。由此定义一个{\it 骨架函子}
$$
\text{sk}_n :
\text{Simp}(\mathcal{C}) \longrightarrow \text{Simp}_n(\mathcal{C})
$$
它从 $\mathcal{C}$ 的单纯对象范畴到 $\mathcal{C}$ 的
$n$-截断单纯对象范畴。为避免与文献中的其他函子混淆，见备注
\ref{simplicial-remark-sk-literature}。




\section{与单纯集的积}
\label{simplicial-section-product-with-simplicial-sets}

\noindent
设 $\mathcal{C}$ 为范畴，$U$ 为单纯集，$V$ 为 $\mathcal{C}$ 的单纯对象。
可以考虑如下协变函子：它把 $\mathcal{C}$ 的单纯对象 $W$ 送到集合
\begin{equation}
\label{simplicial-equation-functor-product-with-simplicial-set}
\left\{
(f_{n, u} : V_n \to W_n)_{n \geq 0, u \in U_n}
\text{ 使得 }
\begin{matrix}
\forall \varphi : [m] \to [n] \\
f_{m, U(\varphi)(u)} \circ V(\varphi) = W(\varphi) \circ f_{n, u}
\end{matrix}
\right\}
\end{equation}
若此函子具有形式 $\Mor_{\text{Simp}(\mathcal{C})}(Q, -)$，则可以把
$Q$ 看作 $U$ 与 $V$ 的积。我们不作这种形式化，而是直接作如下定义。

\begin{definition}
\label{simplicial-definition-product-with-simplicial-set}
设 $\mathcal{C}$ 为范畴，且 $\mathcal{C}$ 的任意两个对象的余积都存在。
设 $U$ 为单纯集，
$V$ 为 $\mathcal{C}$ 的单纯对象，并假设每个 $U_n$ 都是有限非空的。
在这种情形下，定义 $U$ 与 $V$ 的{\it 积 $U \times V$}为
$\mathcal{C}$ 的如下单纯对象，其第 $n$ 项为
$$
(U \times V)_n = \coprod\nolimits_{u\in U_n} V_n
$$
对 $\varphi : [m] \to [n]$，相应映射为
$$
\coprod\nolimits_{u\in U_n} V_n
\longrightarrow
\coprod\nolimits_{u'\in U_m} V_m
$$
它通过态射 $V(\varphi)$，把对应于 $u$ 的分量 $V_n$ 映到对应于
$u' = U(\varphi)(u)$ 的分量 $V_m$。
更宽泛地，若上面列出的所有余积都存在（不对 $\mathcal{C}$ 作任何假设），
则称{\it 积 $U \times V$ 存在}。
\end{definition}

\begin{lemma}
\label{simplicial-lemma-check-product-with-simplicial-set}
设 $\mathcal{C}$ 为范畴，且 $\mathcal{C}$ 的任意两个对象的余积都存在。
设 $U$ 为单纯集，
$V$ 为 $\mathcal{C}$ 的单纯对象，并假设每个 $U_n$ 都是有限非空的。
函子 $W \mapsto \Mor_{\text{Simp}(\mathcal{C})}(U \times V, W)$
与把 $W$ 送到公式
(\ref{simplicial-equation-functor-product-with-simplicial-set}) 中集合的函子典范同构。
\end{lemma}

\begin{proof}
略。
\end{proof}

\begin{lemma}
\label{simplicial-lemma-back-to-U}
设 $\mathcal{C}$ 为范畴，且 $\mathcal{C}$ 的任意两个对象的余积都存在。
暂以
$\textit{FSSets}$ 表示各分量均有限非空的单纯集所成的范畴。
\begin{enumerate}
\item 规则 $(U, V) \mapsto U \times V$ 定义了函子
$\textit{FSSets} \times \text{Simp}(\mathcal{C})
\to \text{Simp}(\mathcal{C})$。
\item 对上面的每一对 $U$、$V$，都有单纯对象之间的典范映射
$$
U \times V \longrightarrow V
$$
它在 $(U \times V)_n = \coprod_u V_n$ 的每个分量上取恒等映射。
\end{enumerate}
\end{lemma}

\begin{proof}
略。
\end{proof}

\noindent
我们简要考察上述构造的一种特殊情形。设 $\mathcal{C}$ 为范畴，
$X$ 为 $\mathcal{C}$ 的对象，$k \geq 0$ 为整数。若所有余积
$X \amalg \ldots \amalg X$ 都存在，则按上述定义，积
$$
X \times \Delta[k]
$$
存在；这里把 $X$ 看作相应的常值单纯对象。

\begin{lemma}
\label{simplicial-lemma-morphism-from-coproduct}
沿用上面的 $X$ 和 $k$。对 $\mathcal{C}$ 的任意单纯对象 $V$，
有如下典范双射
$$
\Mor_{\text{Simp}(\mathcal{C})}(X \times \Delta[k], V)
\longrightarrow
\Mor_\mathcal{C}(X, V_k).
$$
它把 $\gamma$ 送到态射 $\gamma_k$ 在对应于
$\text{id}_{[k]}$ 的分量上的限制。类似地，对任意 $n \geq k$，
若 $W$ 是 $\mathcal{C}$ 的 $n$-截断单纯对象，则有
$$
\Mor_{\text{Simp}_n(\mathcal{C})}(\text{sk}_n(X \times \Delta[k]), W)
=
\Mor_\mathcal{C}(X, W_k).
$$
\end{lemma}

\begin{proof}
态射 $\gamma : X \times \Delta[k] \to V$ 由一族态射
$\gamma_\alpha : X \to V_n$ 给出，其中 $\alpha : [n] \to [k]$。
这些态射必须满足：对所有 $\varphi : [m] \to [n]$，图表
$$
\xymatrix{
X \ar[r]^{\gamma_\alpha} \ar[d]^{\text{id}_X} & V_n \ar[d]^{V(\varphi)} \\
X \ar[r]^{\gamma_{\alpha \circ \varphi}} & V_m
}
$$
可交换。取 $\alpha = \text{id}_{[k]}$，可见对任意
$\varphi : [m] \to [k]$，都有
$\gamma_\varphi = V(\varphi) \circ \gamma_{\text{id}_{[k]}}$。
因此，态射 $\gamma$ 由 $\gamma$ 在对应于 $\text{id}_{[k]}$ 的分量上的取值决定。
反之，给定这样的态射 $f : X \to V_k$，令
$\gamma_\alpha = V(\alpha) \circ f$，即可构造态射 $\gamma$。

\medskip\noindent
截断情形类似，留给读者验证。
\end{proof}

\noindent
一个特殊例子是 $k = 0$ 的情形。此时，引理中的公式只是在说
$$
\Mor_\mathcal{C}(X, V_0)
=
\Mor_{\text{Simp}(\mathcal{C})}(X, V)
$$
其中右边的 $X$ 表示取值为 $X$ 的常值单纯对象。下文将直接使用这个公式，
不再另行说明。




\section{从单纯集到余单纯对象的 Hom}
\label{simplicial-section-hom-from-simplicial-sets-into-cosimplicial}

\noindent
设 $\mathcal{C}$ 为范畴，$U$ 为 $\mathcal{C}$ 的单纯对象，
$V$ 为 $\mathcal{C}$ 的余单纯对象。于是得到如下余单纯集
$\Hom_\mathcal{C}(U, V)$：
\begin{enumerate}
\item 令 $\Hom_\mathcal{C}(U, V)_n = \Mor_\mathcal{C}(U_n, V_n)$；
\item 对 $\varphi : [m] \to [n]$，取映射
$\Hom_\mathcal{C}(U, V)_m \to \Hom_\mathcal{C}(U, V)_n$，
它由 $f \mapsto V(\varphi) \circ f \circ U(\varphi)$ 给出。
\end{enumerate}
这正是下述定义的动机。

\begin{definition}
\label{simplicial-definition-hom-deltak-cosimplicial}
设 $\mathcal{C}$ 为有有限积的范畴，$V$ 为 $\mathcal{C}$ 的余单纯对象，
$U$ 为单纯集，并且每个 $U_n$ 都是有限非空的。定义
{\it $\Hom(U, V)$} 为 $\mathcal{C}$ 的如下余单纯对象：
\begin{enumerate}
\item 令 $\Hom(U, V)_n = \prod_{u \in U_n} V_n$；换言之，
它是 $\mathcal{C}$ 中唯一满足下列 $X$-值点等式的对象：
$$
\Mor_\mathcal{C}(X, \Hom(U, V)_n)
=
\text{Map}(U_n, \Mor_\mathcal{C}(X, V_n))
$$
并且
\item 对 $\varphi : [m] \to [n]$，取映射
$\Hom(U, V)_m \to \Hom(U, V)_n$；在上述 $X$-值点上，
它由 $f \mapsto V(\varphi) \circ f \circ U(\varphi)$ 给出。
\end{enumerate}
\end{definition}

\noindent
把这个定义用积之间的映射完整写出，留给读者。还要指出，该构造关于 $U$
反变、关于 $V$ 协变，正如引理 \ref{simplicial-lemma-back-to-U} 中单纯集与单纯对象之积
的情形一样。




\section{从余单纯集到单纯对象的 Hom}
\label{simplicial-section-hom-from-cosimplicial-sets-into-simplicial}

\noindent
设 $\mathcal{C}$ 为范畴，$U$ 为 $\mathcal{C}$ 的余单纯对象，
$V$ 为 $\mathcal{C}$ 的单纯对象。于是得到如下单纯集
$\Hom_\mathcal{C}(U, V)$：
\begin{enumerate}
\item 令 $\Hom_\mathcal{C}(U, V)_n = \Mor_\mathcal{C}(U_n, V_n)$；
\item 对 $\varphi : [m] \to [n]$，取映射
$\Hom_\mathcal{C}(U, V)_n \to \Hom_\mathcal{C}(U, V)_m$，
它由 $f \mapsto V(\varphi) \circ f \circ U(\varphi)$ 给出。
\end{enumerate}
这正是下述定义的动机。

\begin{definition}
\label{simplicial-definition-hom-deltak-simplicial}
设 $\mathcal{C}$ 为有有限积的范畴，$V$ 为 $\mathcal{C}$ 的单纯对象，
$U$ 为余单纯集，并且每个 $U_n$ 都是有限非空的。定义
{\it $\Hom(U, V)$} 为 $\mathcal{C}$ 的如下单纯对象：
\begin{enumerate}
\item 令 $\Hom(U, V)_n = \prod_{u \in U_n} V_n$；换言之，
它是 $\mathcal{C}$ 中唯一满足下列 $X$-值点等式的对象：
$$
\Mor_\mathcal{C}(X, \Hom(U, V)_n)
=
\text{Map}(U_n, \Mor_\mathcal{C}(X, V_n))
$$
并且
\item 对 $\varphi : [m] \to [n]$，取映射
$\Hom(U, V)_n \to \Hom(U, V)_m$；在上述 $X$-值点上，
它由 $f \mapsto V(\varphi) \circ f \circ U(\varphi)$ 给出。
\end{enumerate}
\end{definition}

\noindent
把这个定义用积之间的映射完整写出，留给读者。还要指出，该构造关于 $U$
反变、关于 $V$ 协变，正如引理 \ref{simplicial-lemma-back-to-U} 中单纯集与单纯对象之积
的情形一样。

\medskip\noindent
我们在一种特殊情形下把上述构造明确写出。设 $X$ 为范畴
$\mathcal{C}$ 的对象，并假设自积 $X \times \ldots \times X$ 存在。
设 $k$ 为整数。考虑各项为
$$
U_n = \prod\nolimits_{\alpha \in \Mor([k], [n])} X
$$
的单纯对象 $U$；对 $\varphi : [m] \to [n]$，映射为
$$
U(\varphi) :
\prod\nolimits_{\alpha \in \Mor([k], [n])} X
\longrightarrow
\prod\nolimits_{\alpha' \in \Mor([k], [m])} X, \quad
(f_{\alpha})_{\alpha} \longmapsto
(f_{\varphi \circ \alpha'})_{\alpha'}
$$
用“坐标”表示，元素 $(x_\alpha)_\alpha$ 被映到元素
$(x_{\varphi \circ \alpha'})_{\alpha'}$。我们断言，这个对象等于
$\Hom(C[k], X)$；这里把 $X$ 看作取值为 $X$ 的常值单纯对象，而 $C[k]$ 是例
\ref{simplicial-example-simplex-cosimplicial-set} 中的余单纯集。

\begin{lemma}
\label{simplicial-lemma-morphism-into-product}
沿用上面的 $X$、$k$ 和 $U$。
\begin{enumerate}
\item 对 $\mathcal{C}$ 的任意单纯对象 $V$，有如下典范双射
$$
\Mor_{\text{Simp}(\mathcal{C})}(V, U)
\longrightarrow
\Mor_\mathcal{C}(V_k, X).
$$
它把 $\gamma$ 送到态射 $\gamma_k$ 与投影到对应于
$\text{id}_{[k]}$ 的因子的复合。
\item 类似地，若 $W$ 是 $\mathcal{C}$ 的 $k$-截断单纯对象，则有
$$
\Mor_{\text{Simp}_k(\mathcal{C})}(W, \text{sk}_k U)
=
\Mor_\mathcal{C}(W_k, X).
$$
\item 上面构造的对象 $U$ 是 $\Hom(C[k], X)$ 的一种实现；其中 $C[k]$
是例 \ref{simplicial-example-simplex-cosimplicial-set} 中的余单纯集。
\end{enumerate}
\end{lemma}

\begin{proof}
先证明 (1)。设 $\gamma : V \to U$ 为态射。它由一族态射
$\gamma_{\alpha} : V_n \to X$ 给出，其中 $\alpha : [k] \to [n]$。
这些态射必须满足：对所有 $\varphi : [m] \to [n]$，图表
$$
\xymatrix{
X \ar[d]^{\text{id}_X} &
V_n \ar[d]^{V(\varphi)}
\ar[l]^{\gamma_{\varphi \circ \alpha'}} \\
X &
V_m \ar[l]_{\gamma_{\alpha'}}
}
$$
对所有 $\alpha' : [k] \to [m]$ 都可交换。取
$\alpha' = \text{id}_{[k]}$，可见对任意 $\varphi : [k] \to [n]$，
都有 $\gamma_\varphi = \gamma_{\text{id}_{[k]}} \circ V(\varphi)$。
因此，态射 $\gamma$ 由 $\gamma_k$ 对应于 $\text{id}_{[k]}$ 的分量决定。
反之，给定这样的态射 $f : V_k \to X$，令
$\gamma_\alpha = f \circ V(\alpha)$，即可构造态射 $\gamma$。

\medskip\noindent
截断情形类似，留给读者验证。

\medskip\noindent
(3) 由 $U$ 的构造以及等式
$C[k]_n = \Mor([k], [n])$ 立即推出；后者正是构造 $U_n$ 时使用的指标集。
\end{proof}




\section{内部 Hom}
\label{simplicial-section-internal-hom}

\noindent
设 $\mathcal{C}$ 为有有限非空积的范畴，$U$、$V$ 为单纯对象
$\mathcal{C}$。
% P02-E0106: frozen source omits “of” before the ambient category.
在某些情形下，函子
$$
\text{Simp}(\mathcal{C})^{opp} \longrightarrow \textit{Sets}, \quad
W \longmapsto \Mor_{\text{Simp}(\mathcal{C})}(W \times V, U)
$$
是可表示的。在这种情形下，把所得的 $\mathcal{C}$ 的单纯对象记作
$\SheafHom(V, U)$，并称从 $V$ 到 $U$ 的{\it 内部 Hom 存在}。
此外，在这种情形下，给定 $\mathcal{C}$ 中的 $X$，有
\begin{align*}
\Mor_\mathcal{C}(X, \SheafHom(V, U)_n)
& =
\Mor_{\text{Simp}(\mathcal{C})}(X \times \Delta[n], \SheafHom(V, U)) \\
& =
\Mor_{\text{Simp}(\mathcal{C})}(X \times \Delta[n]\times V, U) \\
& =
\Mor_{\text{Simp}(\mathcal{C})}(X, \SheafHom(\Delta[n] \times V, U)) \\
& =
\Mor_\mathcal{C}(X, \SheafHom(\Delta[n] \times V, U)_0)
\end{align*}
只要 $\SheafHom(\Delta[n] \times V, U)$ 也存在即可。第一个和最后一个等式
由引理 \ref{simplicial-lemma-morphism-from-coproduct} 推出。

\medskip\noindent
由此得到的要点是：给定 $U$ 和 $V$，若要构造内部 Hom，就应尝试构造对象
$$
\SheafHom(\Delta[n] \times V, U)_0
$$
因为它们应当是 $\SheafHom(V, U)$ 的第 $n$ 项。下一节将研究单纯对象
“$\Hom(\Delta[n], U)$”的一种构造。


\section{从单纯集到单纯对象的 Hom}
\label{simplicial-section-hom-from-simplicial-sets}

\noindent
受上述内部 Hom 讨论的启发，我们定义一个单纯对象，它应当分类从单纯集到
范畴 $\mathcal{C}$ 的某个给定单纯对象的态射。

\begin{definition}
\label{simplicial-definition-hom-from-simplicial-set}
设 $\mathcal{C}$ 为任意两个对象的余积都存在的范畴。设 $U$ 为单纯集，
并且对所有 $n \geq 0$，$U_n$ 都有限非空。设 $V$ 为
$\mathcal{C}$ 的单纯对象。以 {\it $\Hom(U, V)$} 表示任意满足下式的
$\mathcal{C}$ 的单纯对象：
$$
\Mor_{\text{Simp}(\mathcal{C})}(W, \Hom(U, V))
=
\Mor_{\text{Simp}(\mathcal{C})}(W \times U, V)
$$
其中等式关于 $\mathcal{C}$ 的单纯对象 $W$ 函子性地成立。
\end{definition}

\noindent
当然，$\Hom(U, V)$ 未必存在。此外，由第 \ref{simplicial-section-internal-hom} 节
的讨论可预期，若它存在，则
$\Hom(U, V)_n = \Hom(U \times \Delta[n], V)_0$。
这里不对这些 Hom 对象使用斜体记号，因为 $\Hom(U, V)$ 不是内部 Hom。

\begin{lemma}
\label{simplicial-lemma-exists-hom-0-from-simplicial-set}
假设范畴 $\mathcal{C}$ 有任意两个对象的余积和可数极限。
设 $U$ 为单纯集，并且对所有 $n \geq 0$，$U_n$ 都有限非空。
设 $V$ 为 $\mathcal{C}$ 的单纯对象。则函子
\begin{eqnarray*}
\mathcal{C}^{opp} & \longrightarrow & \textit{Sets} \\
X
& \longmapsto &
\Mor_{\text{Simp}(\mathcal{C})}(X \times U, V)
\end{eqnarray*}
是可表示的。
\end{lemma}

\begin{proof}
从 $X \times U$ 到 $V$ 的态射由一族态射
$f_u : X \to V_n$ 给出，其中 $n \geq 0$ 且 $u \in U_n$。
而这样的一族态射恰好在下述条件下定义一个态射：对所有
$\varphi : [m] \to [n]$，图表
$$
\xymatrix{
X \ar[r]^{f_u} \ar[d]_{\text{id}_X} & V_n \ar[d]^{V(\varphi)} \\
X \ar[r]^{f_{U(\varphi)(u)}} & V_m
}
$$
都可交换。因此自然要引入范畴 $\mathcal{U}$ 和函子
$\mathcal{V} : \mathcal{U}^{opp} \to \mathcal{C}$，定义如下：
\begin{enumerate}
\item $\mathcal{U}$ 的对象集是 $\coprod_{n \geq 0} U_n$；
\item 从 $u' \in U_m$ 到 $u \in U_n$ 的态射，是满足
$U(\varphi)(u) = u'$ 的 $\varphi : [m] \to [n]$；
\item 对 $u \in U_n$，令 $\mathcal{V}(u) = V_n$；且
\item 对满足 $U(\varphi)(u) = u'$ 的 $\varphi : [m] \to [n]$，令
$\mathcal{V}(\varphi) = V(\varphi) : V_n \to V_m$。
\end{enumerate}
此时很清楚，我们的函子正是定义
$$
\lim_{\mathcal{U}^{opp}} \mathcal{V}
$$
的函子。因此，若 $\mathcal{C}$ 有可数极限，则此极限存在，从而表示引理中
函子的对象也存在。
\end{proof}

\begin{lemma}
\label{simplicial-lemma-exists-hom-0-from-simplicial-set-finite}
假设范畴 $\mathcal{C}$ 有任意两个对象的余积和有限极限。
设 $U$ 为单纯集，并且对所有 $n \geq 0$，$U_n$ 都有限非空。
假设对所有 $n \gg 0$，$U$ 的每个 $n$-单形都是退化的。
设 $V$ 为 $\mathcal{C}$ 的单纯对象。则函子
\begin{eqnarray*}
\mathcal{C}^{opp} & \longrightarrow & \textit{Sets} \\
X
& \longmapsto &
\Mor_{\text{Simp}(\mathcal{C})}(X \times U, V)
\end{eqnarray*}
是可表示的。
\end{lemma}

\begin{proof}
我们必须证明，引理 \ref{simplicial-lemma-exists-hom-0-from-simplicial-set} 的证明中
所述范畴 $\mathcal{U}$ 有一个有限子范畴 $\mathcal{U}'$，使得
$\mathcal{V}$ 在 $\mathcal{U}'$ 上的极限与 $\mathcal{V}$ 在
$\mathcal{U}$ 上的极限相同。我们将使用 Categories，引理
\ref{categories-lemma-initial}。对 $m > 0$，以
$\mathcal{U}_{\leq m}$ 表示对象为
$\coprod_{0 \leq n \leq m} U_m$ 的全子范畴。
% P02-E0107: frozen source uses U_m under an index n; likely U_n.
取整数 $m_0$，使得当 $n > m_0$ 时，单纯集 $U$ 的每个 $n$-单形都退化。
对任意充分大的 $m \geq m_0$，子范畴 $\mathcal{U}_{\leq m}$ 满足
Categories，定义 \ref{categories-definition-initial} 的性质 (1)。

\medskip\noindent
设 $u \in U_n$、$u' \in U_{n'}$，其中 $n, n' \leq m_0$；并设
$\varphi : [k] \to [n]$、$\varphi' : [k] \to [n']$ 为满足
$U(\varphi)(u) = U(\varphi')(u')$ 的态射。一个简单的组合论论证表明，
若 $k > 2m_0$，则存在指标 $0 \leq i \leq 2m_0$，使得
$\varphi(i) =\varphi(i + 1)$ 且
$\varphi'(i) = \varphi'(i + 1)$。（鸽巢原理会告诉你，当
$k > m_0^2$ 时这也成立；对下面的论证而言已经足够。）因此，若
$k > 2m_0$，则可写成
$\varphi = \psi \circ \sigma^{k - 1}_i$ 和
$\varphi' = \psi' \circ \sigma^{k - 1}_i$，其中
$\psi : [k - 1] \to [n]$、$\psi' : [k - 1] \to [n']$。
由于 $s^{k - 1}_i : U_{k - 1} \to U_k$ 是单射（见引理
\ref{simplicial-lemma-si-injective}），也可得
$U(\psi)(u) = U(\psi')(u')$。继续这一过程可知：给定
$\mathcal{U}$ 的态射 $u \to z$ 和 $u' \to z$，其中
$u, u' \in \mathcal{U}_{\leq m_0}$，存在可交换图表
$$
\xymatrix{
u \ar[rd] \ar[rrd] & & \\
& a \ar[r] & z \\
u' \ar[ru] \ar[rru]
}
$$
其中 $a \in \mathcal{U}_{\leq 2m_0}$。

\medskip\noindent
容易由此推出，有限子范畴 $\mathcal{U}_{\leq 2m_0}$ 符合要求。
事实上，设给定 $x' \in U_n$ 和 $x'' \in U_{n'}$，其中
$n, n' \leq 2m_0$；并给定 $\mathcal{U}$ 中具有同一目标的态射
$x' \to x$ 和 $x'' \to x$。由 $m_0$ 的选取，可以找到
$\mathcal{U}_{\leq m_0}$ 的对象 $u, u'$ 以及态射
$u \to x'$、$u' \to x''$。由上述结论，可以找到
$a \in \mathcal{U}_{\leq 2m_0}$ 以及态射 $u \to a$、$u' \to a$，
使得
$$
\xymatrix{
u \ar[rd] \ar[rrd] \ar[r] & x' \ar[rd] & \\
& a \ar[r] & x \\
u' \ar[ru] \ar[rru] \ar[r] & x'' \ar[ru] &
}
$$
可交换。把这个图顺时针旋转 90 度，就得到 Categories，定义
\ref{categories-definition-initial} 中 (2) 所要求的图表。
\end{proof}

\begin{lemma}
\label{simplicial-lemma-exists-hom-from-simplicial-set-finite}
假设范畴 $\mathcal{C}$ 有任意两个对象的余积和有限极限。
设 $U$ 为单纯集，并且对所有 $n \geq 0$，$U_n$ 都有限非空。
假设对所有 $n \gg 0$，$U$ 的每个 $n$-单形都是退化的。
设 $V$ 为 $\mathcal{C}$ 的单纯对象。则 $\Hom(U, V)$ 存在，而且有预期的等式
$$
\Hom(U, V)_n = \Hom(U \times \Delta[n], V)_0.
$$
\end{lemma}

\begin{proof}
如下构造这个单纯对象。对 $n \geq 0$，以 $\Hom(U, V)_n$ 表示
$\mathcal{C}$ 中表示函子
$$
X
\longmapsto
\Mor_{\text{Simp}(\mathcal{C})}(X \times U \times \Delta[n], V)
$$
的对象。由引理 \ref{simplicial-lemma-exists-hom-0-from-simplicial-set-finite}，
它存在，因为 $U \times \Delta[n]$ 是一个各单形集有限、且在足够高次数
没有非退化单形的单纯集；见引理 \ref{simplicial-lemma-product-degenerate}。
对 $\varphi : [m] \to [n]$，得到诱导的单纯集映射
$\varphi : \Delta[m] \to \Delta[n]$。因而得到关于 $X$ 函子性的态射
$X \times U \times \Delta[m] \to X \times U \times \Delta[n]$，
进而得到函子之间的变换，再进而给出
$$
\Hom(U, V)(\varphi) :
\Hom(U, V)_n
\longrightarrow
\Hom(U, V)_m.
$$
显然，这定义了从 $\Delta$ 到范畴 $\mathcal{C}$ 的反变函子
$\Hom(U, V)$。换言之，我们得到 $\mathcal{C}$ 的一个单纯对象。

\medskip\noindent
还需证明 $\Hom(U, V)$ 满足所需的泛性质
$$
\Mor_{\text{Simp}(\mathcal{C})}(W, \Hom(U, V))
=
\Mor_{\text{Simp}(\mathcal{C})}(W \times U, V)
$$
为此，给定 $f : W \to \Hom(U, V)$。我们要构造右端的元素
$f' : W \times U \to V$。由构造，每个
$f_n : W_n \to \Hom(U, V)_n$ 对应于态射
$f_n : W_n \times U \times \Delta[n] \to V$。此外，
对每个态射 $\varphi : [m] \to [n]$，图表
$$
\xymatrix{
W_n \times U \times \Delta[m]
\ar[rr]_{W(\varphi)\times \text{id} \times \text{id}}
\ar[d]_{\text{id} \times \text{id} \times \varphi} & &
W_m \times U \times \Delta[m] \ar[d]^{f_m} \\
W_n \times U \times \Delta[n] \ar[rr]^{f_n} & & V
}
$$
可交换。对 $(\Delta[n])_k$ 中的 $\psi : [n] \to [k]$，以
$(f_n)_{k, \psi} : W_n \times U_k \to V_k$ 表示 $(f_n)_k$
对应于元素 $\psi$ 的分量。定义
$f'_n : W_n \times U_n \to V_n$ 为
$f'_n = (f_n)_{n, \text{id}}$；换言之，它是
$(f_n)_n : W_n \times U_n \times (\Delta[n])_n \to V_n$
在 $W_n \times U_n \times \text{id}_{[n]}$ 上的限制。
为证明族 $(f'_n)$ 定义了单纯对象之间的态射，必须证明对任意
$\varphi : [m] \to [n]$，都有
$V(\varphi) \circ f'_n =
f'_m \circ W(\varphi) \times U(\varphi)$。
上面的可交换图表说明，
$(f_n)_{m, \varphi} : W_n \times U_m \to V_m$ 等于
$(f_m)_{m, \text{id}} \circ W(\varphi) :
W_n \times U_m \to V_m$。
另一方面，$f_n$ 是单纯对象之间的态射这一事实蕴含图表
$$
\xymatrix{
W_n \times U_n \times (\Delta[n])_n
\ar[r]_-{(f_n)_n}
\ar[d]_{\text{id} \times U(\varphi) \times \varphi}
& V_n \ar[d]^{V(\varphi)} \\
W_n \times U_m \times (\Delta[n])_m \ar[r]^-{(f_n)_m} & V_m
}
$$
可交换。这又蕴含 $(f_n)_{m, \varphi} \circ U(\varphi)$ 等于
$V(\varphi) \circ (f_n)_{n, \text{id}}$。综上，得到
$
V(\varphi) \circ (f_n)_{n, \text{id}}
=
(f_n)_{m, \varphi} \circ U(\varphi)
=
(f_m)_{m, \text{id}} \circ W(\varphi)\circ U(\varphi)
=
(f_m)_{m, \text{id}} \circ W(\varphi)\times U(\varphi)
$
如所欲证。

\medskip\noindent
反过来，给定态射 $f' : W \times U \to V$，我们如下定义态射
$f : W \to \Hom(U, V)$。由引理 \ref{simplicial-lemma-morphism-from-coproduct}，
态射 $\text{id} : W_n \to W_n$ 对应于唯一态射
$c_n : W_n \times \Delta[n] \to W$。因此可以考虑复合
$$
W_n \times \Delta[n] \times U
\xrightarrow{c_n}
W \times U
\xrightarrow{f'}
V.
$$
由构造，这对应于唯一态射 $f_n : W_n \to \Hom(U, V)_n$。
所需的这些态射定义了单纯集之间的态射这一点，留给读者验证。
% P02-E0108: frozen source uses plural “morphisms” with singular id.
% P02-E0109: frozen source says simplicial sets although these are C-valued simplicial objects.

\medskip\noindent
同样留给读者验证，$f \mapsto f'$ 和 $f' \mapsto f$ 是互逆操作。
\end{proof}

\begin{lemma}
\label{simplicial-lemma-hom-from-coprod}
假设范畴 $\mathcal{C}$ 有任意两个对象的余积和有限极限。
设 $a : U \to V$、$b : U \to W$ 为单纯集之间的态射。
假设对所有 $n \geq 0$，$U_n, V_n, W_n$ 都有限非空。
假设对所有 $n \gg 0$，$U, V, W$ 的每个 $n$-单形都是退化的。
设 $T$ 为 $\mathcal{C}$ 的单纯对象。则
$$
\Hom(V, T) \times_{\Hom(U, T)} \Hom(W, T)
=
\Hom(V \amalg_U W, T)
$$
换言之，左边的纤维积由右边的 Hom 对象表示。
\end{lemma}

\begin{proof}
由引理 \ref{simplicial-lemma-exists-hom-from-simplicial-set-finite}，所需的所有
$\Hom$ 对象都存在并满足正确的函子性。现在可以把左边的第 $n$ 项识别为
表示把 $X$ 送到如下函子性等式链中第一个集合的对象：
\begin{align*}
&
\Mor(X \times \Delta[n],
\Hom(V, T) \times_{\Hom(U, T)} \Hom(W, T)) \\
& =
\Mor(X \times \Delta[n], \Hom(V, T))
\times_{\Mor(X \times \Delta[n], \Hom(U, T))}
\Mor(X \times \Delta[n], \Hom(W, T)) \\
& =
\Mor(X \times \Delta[n] \times V, T)
\times_{\Mor(X \times \Delta[n] \times U, T)}
\Mor(X \times \Delta[n] \times W, T) \\
& =
\Mor(X \times \Delta[n] \times (V \amalg_U W), T))
\end{align*}
% P02-E0110: frozen source has an unmatched extra closing parenthesis in the final Hom term.
这里使用了等式
$$
(X \times \Delta[n] \times V)
\times_{X \times \Delta[n] \times U}
(X \times \Delta[n] \times W)
=
X \times \Delta[n] \times (V \amalg_U W)
$$
% P02-E0111: frozen source uses a fibre product on maps induced by U -> V and U -> W; the typed construction is a pushout.
逐项验证此等式很容易。引理的结论随即由上列等式的最后一项对应于
$\Hom(V \amalg_U W, T)_n$ 得出。
\end{proof}

\section{单纯对象的分裂}
\label{simplicial-section-splitting}

\noindent
范畴 $\mathcal{C}$ 的对象 $X$ 的子对象 $N$，是 $\mathcal{C}$ 的一个对象
$N$ 连同一个单态射 $N \to X$。当然，如果存在与到 $X$ 的态射相容的
同构 $N \to N'$，我们就（滥用记号地）称子对象 $N$、$N'$ 相等。
子对象全体构成偏序集。（这是由于我们对范畴的约定；例如，对差一个同伦
意义下的空间范畴并不成立。）

\begin{definition}
\label{simplicial-definition-split}
设 $\mathcal{C}$ 为有有限非空余积的范畴。称 $\mathcal{C}$ 的单纯对象
$U$ 是{\it 分裂的}，如果存在 $U_m$ 的子对象 $N(U_m)$（$m \geq 0$），
使得对所有 $n \geq 0$，映射
\begin{equation}
\label{simplicial-equation-splitting}
\coprod\nolimits_{\varphi : [n] \to [m]\text{ 满射}}
N(U_m)
\longrightarrow
U_n
\end{equation}
都是同构。若 $U$ 是 $\mathcal{C}$ 的 $r$-截断单纯对象，则称 $U$
是{\it 分裂的}，如果存在 $U_m$ 的子对象 $N(U_m)$（$r \geq m \geq 0$），
使得对 $r \geq n \geq 0$，映射 (\ref{simplicial-equation-splitting}) 是同构。
\end{definition}

\noindent
在这种情形下，$N(U_0) = U_0$。其次有
$U_1 = U_0 \amalg N(U_1)$。再者有
$$
U_2 = U_0 \amalg N(U_1) \amalg N(U_1) \amalg N(U_2).
$$
事实表明，在许多范畴 $\mathcal{C}$ 中，每个单纯对象都是分裂的。

\begin{lemma}
\label{simplicial-lemma-splitting-simplicial-sets}
设 $U$ 为单纯集。则 $U$ 有唯一的分裂，其中 $N(U_m)$ 等于非退化
$m$-单形的集合。
\end{lemma}

\begin{proof}
由定义立即可知，若存在分裂，则 $N(U_m)$ 必须是非退化单形的集合。
设 $x \in U_n$。假设存在满射 $\varphi : [n] \to [k]$ 和
$\psi : [n] \to [l]$，以及非退化单形 $y \in U_k$、$z \in U_l$，
使得 $x = U(\varphi)(y)$ 且 $x = U(\psi)(z)$。取 $\psi$ 的右逆
$\xi : [l] \to [n]$，即 $\psi \circ \xi = \text{id}_{[l]}$。
于是 $z = U(\xi)(x)$，从而
$z = U(\xi)(x) = U(\varphi \circ \xi)(y)$。由于 $z$ 非退化，
可知 $\varphi \circ \xi : [l] \to [k]$ 是满射，故 $l \geq k$。
类似地，$k \geq l$。因此，对 $\psi$ 的右逆 $\xi$ 的任意选取，
$\varphi \circ \xi : [l] \to [k]$ 都必须是恒等映射。
这很容易推出 $\psi = \varphi$。
\end{proof}

\noindent
当然，单纯集之间的映射可能把非退化 $n$-单形映到退化 $n$-单形。
因此，引理 \ref{simplicial-lemma-splitting-simplicial-sets} 的分裂不是函子性的。
下面给出一种它具有函子性的情形。

\begin{lemma}
\label{simplicial-lemma-injective-map-simplicial-sets}
设 $f : U \to V$ 为单纯集之间的态射。假设 (a) $U$ 的每个非退化单形
的像都是 $V$ 的非退化单形，并且 (b) $f$ 从 $U$ 的非退化单形集合
到 $V$ 的非退化单形集合的限制是单射。则对所有 $n$，$f_n$ 都是单射。
把“单射”替换成“满射”或“双射”时，结论同样成立。
\end{lemma}

\begin{proof}
在假设 (a) 下，映射 $f$ 保持引理
\ref{simplicial-lemma-splitting-simplicial-sets} 的分裂给出的不交并分解；
换言之，得到可交换图表
$$
\xymatrix{
\coprod\nolimits_{\varphi : [n] \to [m]\text{ 满射}}
N(U_m)
\ar[r] \ar[d] &
U_n \ar[d] \\
\coprod\nolimits_{\varphi : [n] \to [m]\text{ 满射}}
N(V_m)
\ar[r] &
V_n.
}
$$
于是 (b) 清楚地表明，左侧竖直箭头是单射（相应地，满射或双射）。
\end{proof}

\begin{lemma}
\label{simplicial-lemma-simplicial-set-n-skel-sub}
设 $U$ 为单纯集，$n \geq 0$ 为整数。规则
$$
U'_m = \bigcup\nolimits_{\varphi : [m] \to [i], \ i\leq n} \Im(U(\varphi))
$$
定义了子单纯集 $U' \subset U$，并且当 $i \leq n$ 时
$U'_i = U_i$。此外，对所有 $m > n$，$U'$ 的每个 $m$-单形都是退化的。
\end{lemma}

\begin{proof}
若 $x \in U_m$，且对某个 $y \in U_i$、$i \leq n$ 和某个
$\varphi : [m] \to [i]$ 有 $x = U(\varphi)(y)$，则对任意
$\psi : [m'] \to [m]$，像 $U(\psi)(x)$ 等于
$U(\varphi \circ \psi)(y)$，而且
$\varphi \circ \psi : [m'] \to [i]$。因此 $U'$ 是单纯集。
由构造，所有次数为 $n + 1$ 或更高的单形都是退化的。
\end{proof}

\begin{lemma}
\label{simplicial-lemma-splitting-simplicial-groups}
设 $U$ 为单纯阿贝尔群。则 $U$ 有如下分裂：取
$N(U_0) = U_0$，并对 $m \geq 1$ 取
$$
N(U_m) = \bigcap\nolimits_{i = 0}^{m - 1} \Ker(d^m_i).
$$
此外，这个分裂在单纯阿贝尔群范畴上是函子性的。
\end{lemma}

\begin{proof}
对 $n$ 归纳，我们将证明，引理中对 $N(U_m)$ 的选取保证
(\ref{simplicial-equation-splitting}) 对 $m \leq n$ 是同构。$n = 0$ 时显然。
在本证明的其余部分，为提高可读性，我们将省略映射 $d_i$ 和 $s_i$
的上标。还将反复使用备注 \ref{simplicial-remark-relations} 中的关系。

\medskip\noindent
先作一般性观察。对 $0 \leq i \leq m$ 和 $z \in U_m$，有
$d_i(s_i(z)) = z$。因此，通过取
$x' = (x - s_i(d_i(x)))$ 和 $x'' = s_i(d_i(x))$，可以把任意
$x \in U_{m + 1}$ 唯一写成 $x = x' + x''$，其中
$d_i(x') = 0$ 且 $x'' \in \Im(s_i)$。此外，满足
$x'' = s_i(z)$ 的元素 $z \in U_m$ 是唯一的，因为 $s_i$ 是单射。

\medskip\noindent
下面给出分解任意 $x \in U_{n + 1}$ 的步骤。先写成
$x = x_0 + s_0(z_0)$，其中 $d_0(x_0) = 0$。
再写成 $x_0 = x_1 + s_1(z_1)$，其中
$d_1(x_1) = 0$。如此继续，得到
\begin{eqnarray*}
x & = & x_0 + s_0(z_0), \\
x_0 & = & x_1 + s_1(z_1), \\
x_1 & = & x_2 + s_2(z_2), \\
\ldots & \ldots & \ldots \\
x_{n - 1} & = & x_n + s_n(z_n)
\end{eqnarray*}
其中对所有 $i = n, \ldots, 0$，都有 $d_i(x_i) = 0$。
由上面的一般性观察，所有 $x_i$ 和 $z_i$ 都由 $x$ 唯一确定。
我们断言
$x_i \in
\Ker(d_0) \cap
\Ker(d_1) \cap
\ldots \cap
\Ker(d_i)$
且
$z_i \in
\Ker(d_0) \cap
\ldots \cap
\Ker(d_{i - 1})$
对 $i = n, \ldots, 0$ 成立。在这里以及下文，核的空交表示整个空间；
也就是说，当 $i = 0$ 时，记号
$z_0 \in \Ker(d_0) \cap
\ldots \cap
\Ker(d_{i - 1})$
表示 $z_0 \in U_n$，没有任何限制。

\medskip\noindent
我们对 $i$ 作递增归纳来证明此断言。由 $x_0$ 和 $z_0$ 的构造，
$i = 0$ 时显然。假设对 $i - 1$ 结论成立，来证明
$0 < i \leq n$ 的情形。首先，由构造有 $d_i(x_i) = 0$。
取满足 $0 \leq j < i$ 的 $j$。由归纳假设，$d_j(x_{i - 1}) = 0$。
因此
$$
0 = d_j(x_{i - 1})
= d_j(x_i) + d_j(s_i(z_i))
= d_j(x_i) + s_{i - 1}(d_j(z_i)).
$$
最后一个等式由备注 \ref{simplicial-remark-relations} 的关系得到。
这些关系还蕴含
$d_{i - 1}(d_j(x_i)) = d_j(d_i(x_i)) = 0$，
因为由构造 $d_i(x_i)= 0$。于是，上面一般性观察中的唯一性表明，等式
$0 = x' + x'' = d_j(x_i) + s_{i - 1}(d_j(z_i))$
只能在两项都为零时成立。我们得到 $d_j(x_i) = 0$；又由
$s_{i - 1}$ 的单射性，得到 $d_j(z_i) = 0$。断言得证。

\medskip\noindent
该断言蕴含，可唯一写成
$$
x = s_0(z_0) + s_1(z_1) + \ldots + s_n(z_n) + x_n
$$
其中 $x_n \in N(U_{n + 1})$，且
$z_i \in \Ker(d_0) \cap \ldots \cap \Ker(d_{i - 1})$。
这也可以表述为我们找到了直和分解
$$
U_{n + 1}
=
N(U_{n + 1})
\oplus
\bigoplus\nolimits_{i = 0}^{i = n}
s_i\Big(\Ker(d_0) \cap \ldots \cap \Ker(d_{i - 1})\Big)
$$
并且它满足：对 $j = 0, \ldots, n$，有
$$
\Ker(d_0) \cap \ldots \cap \Ker(d_j)
=
N(U_{n + 1}) \oplus
\bigoplus\nolimits_{i = j + 1}^{i = n}
s_i\Big(\Ker(d_0) \cap \ldots \cap \Ker(d_{i - 1})\Big)
$$
结论可由此得到。$x$ 的表达式中的每个 $z_i$ 都可以唯一写成
$$
z_i = s_i(z'_{i, i}) + \ldots + s_{n - 1}(z'_{i, n - 1}) + z_{i, 0}
$$
其中 $z_{i, 0} \in N(U_n)$，且
$z'_{i, j} \in \Ker(d_0) \cap \ldots \cap \Ker(d_{j - 1})$。
因为 $z_i$ 已经属于 $d_0, \ldots, d_i$ 的核，分解 $z_i$ 时最初的几步
为零。这又唯一给出
$$
x = x_n + s_0(z_{0, 0}) + s_1(z_{1, 0}) + \ldots + s_n(z_{n, 0}) +
\sum\nolimits_{0 \leq i \leq j \leq n - 1} s_i(s_j(z'_{i, j})).
$$
继续这一过程，最终可见，$x$ 被分解成如下形式的项之和：
$$
s_{i_1} s_{i_2} \ldots s_{i_k} (z)
$$
其中 $0 \leq i_1 \leq i_2 \leq \ldots \leq i_k \leq n - k + 1$，
且 $z \in N(U_{n + 1 - k})$。这正是所需的分解，因为任意满射
$[n + 1] \to [n + 1 - k]$ 都可以唯一表示为
$$
\sigma^{n - k}_{i_k} \ldots \sigma^{n - 1}_{i_2} \sigma^n_{i_1}
$$
其中 $0 \leq i_1 \leq i_2 \leq \ldots \leq i_k \leq n - k + 1$。
\end{proof}

\begin{lemma}
\label{simplicial-lemma-splitting-abelian-category}
设 $\mathcal{A}$ 为阿贝尔范畴，$U$ 为 $\mathcal{A}$ 中的单纯对象。
则 $U$ 有如下分裂：取 $N(U_0) = U_0$，并对 $m \geq 1$ 取
$$
N(U_m) = \bigcap\nolimits_{i = 0}^{m - 1} \Ker(d^m_i).
$$
此外，这个分裂在 $\mathcal{A}$ 的单纯对象范畴上是函子性的。
\end{lemma}

\begin{proof}
对 $\mathcal{A}$ 的任意对象 $A$，得到单纯阿贝尔群
$\Mor_\mathcal{A}(A, U)$。其中每一个都由引理
\ref{simplicial-lemma-splitting-simplicial-groups} 典范地分裂。此外，
% P02-E0112: frozen source says “Each of these are” instead of “Each ... is”.
$$
N(\Mor_\mathcal{A}(A, U_m)) =
\bigcap\nolimits_{i = 0}^{m - 1} \Ker(d^m_i) =
\Mor_\mathcal{A}(A, N(U_m)).
$$
因此，对 $\mathcal{A}$ 的任意对象应用函子 $\Mor_\mathcal{A}(A, -)$ 后，
态射 (\ref{simplicial-equation-splitting}) 都成为同构。故由 Yoneda 引理，它本身是同构。
\end{proof}

\begin{lemma}
\label{simplicial-lemma-injective-map-simplicial-abelian}
\begin{slogan}
Dold--Kan 正规化函子反映单射性、满射性和同构性。
\end{slogan}
设 $\mathcal{A}$ 为阿贝尔范畴，$f : U \to V$ 为
$\mathcal{A}$ 的单纯对象之间的态射。若对所有 $i$，诱导态射
$N(f)_i : N(U)_i \to N(V)_i$ 都是单射，则对所有 $i$，$f_i$
都是单射。把“单射”替换成“满射”或“同构”时，结论同样成立。
\end{lemma}

\begin{proof}
这由引理 \ref{simplicial-lemma-splitting-abelian-category} 和分裂的定义立即推出。
\end{proof}


\begin{lemma}
\label{simplicial-lemma-N-d-in-N}
设 $\mathcal{A}$ 为阿贝尔范畴，$U$ 为 $\mathcal{A}$ 中的单纯对象。
令 $N(U_m)$ 如上面的引理 \ref{simplicial-lemma-splitting-abelian-category} 所述。
% P02-E0113: frozen source has “Let N(U_m) as in ...” without a finite verb.
则 $d^m_m(N(U_m)) \subset N(U_{m - 1})$。
\end{lemma}

\begin{proof}
对 $j = 0, \ldots, m - 2$，由备注 \ref{simplicial-remark-relations} 中的关系，有
$d^{m - 1}_j d^m_m = d^{m - 1}_{m - 1} d^m_j$。结论随即成立。
\end{proof}

\begin{lemma}
\label{simplicial-lemma-simplicial-abelian-n-skel-sub}
设 $\mathcal{A}$ 为阿贝尔范畴，$U$ 为 $\mathcal{A}$ 的单纯对象，
$n \geq 0$ 为整数。规则
$$
U'_m = \sum\nolimits_{\varphi : [m] \to [i], \ i\leq n} \Im(U(\varphi))
$$
定义了子单纯对象 $U' \subset U$，并且当 $i \leq n$ 时
$U'_i = U_i$。此外，对所有 $m > n$，都有 $N(U'_m) = 0$。
\end{lemma}

\begin{proof}
取 $m$、$i \leq n$ 和某个 $\varphi : [m] \to [i]$。
对任意 $\psi : [m'] \to [m]$，$\Im(U(\varphi))$ 在 $U(\psi)$
下的像等于 $U(\varphi \circ \psi)$ 的像，并且
$\varphi \circ \psi : [m'] \to [i]$。因此 $U'$ 是单纯对象。
取 $m > n$。必须证明 $N(U'_m) = 0$。由 $N(U_m)$ 和 $N(U'_m)$
的定义，有 $N(U'_m) = U'_m \cap N(U_m)$（子对象的交）。
由于引理 \ref{simplicial-lemma-splitting-abelian-category} 表明 $U$ 是分裂的，
只需证明 $U'_m$ 包含于和
$$
\sum\nolimits_{\varphi : [m] \to [m']\text{ 满射}, \ m' < m}
\Im(U(\varphi)|_{N(U_{m'})}).
$$
由分裂，每个 $U_{m'}$ 都是 $N(U_{m''})$ 在满射
$\psi : [m'] \to [m'']$ 所诱导的 $U(\psi)$ 下的像之和。
因此，上面显示的和等于
$$
\sum\nolimits_{\varphi : [m] \to [m']\text{ 满射}, \ m' < m}
\Im(U(\varphi)).
$$
显然，$U'_m$ 包含于其中：在 $U'_m$ 定义中出现的任意
$\varphi : [m] \to [i]$（$i \leq n$）都可分解为
$[m] \to [m'] \to [i]$，其中 $[m] \to [m']$ 是满射，且
$m' < m$，正如上面最后一个显示的和中那样。
\end{proof}

\section{余骨架函子}
\label{simplicial-section-coskeleton}

\noindent
设 $\mathcal{C}$ 为范畴。{\it 余骨架函子}（若存在）是函子
$$
\text{cosk}_n :
\text{Simp}_n(\mathcal{C}) \longrightarrow \text{Simp}(\mathcal{C})
$$
它是骨架函子的右伴随。用公式表示为
\begin{equation}
\label{simplicial-equation-cosk}
\Mor_{\text{Simp}(\mathcal{C})}(U, \text{cosk}_n V)
=
\Mor_{\text{Simp}_n(\mathcal{C})}(\text{sk}_n U, V)
\end{equation}
给定 $n$-截断单纯对象 $V$，如果存在
$\text{cosk}_nV \in \Ob(\text{Simp}(\mathcal{C}))$ 以及态射
$\text{sk}_n \text{cosk}_n V \to V$，使得上述公式成立，则称
{\it $\text{cosk}_nV$ 存在}；换言之，函子
$U \mapsto \Mor_{\text{Simp}_n(\mathcal{C})}(\text{sk}_n U, V)$
是可表示的。若它存在，则由 Yoneda 引理，它在唯一同构的意义下唯一。
见 Categories，第 \ref{categories-section-opposite} 节。

\begin{example}
\label{simplicial-example-cosk0}
假设范畴 $\mathcal{C}$ 有有限非空自积。$\mathcal{C}$ 的
$0$-截断单纯对象与 $\mathcal{C}$ 的对象 $X$ 是同一回事。
在这种情形下，我们断言 $\text{cosk}_0(X)$ 是单纯对象 $U$，其中
$U_n = X^{n + 1}$ 是 $X$ 的 $(n + 1)$ 重自积，而其单纯对象结构如例
\ref{simplicial-example-fibre-products-simplicial-object} 所述。事实上，态射
$V \to U$（其中 $V$ 是单纯对象）由态射 $V_n \to X^{n + 1}$ 给出，
并且所有图表
$$
\xymatrix{
V_n \ar[r] \ar[d]_{V([0] \to [n], 0 \mapsto i)} &
X^{n + 1} \ar[d]^{\text{pr}_i} \\
V_0 \ar[r] &
X
}
$$
都可交换。显然，这意味着该映射确定并且由唯一态射 $V_0 \to X$ 确定。
这就证明了公式 (\ref{simplicial-equation-cosk})。
\end{example}

\noindent
回忆范畴 $\Delta/[n]$，见例 \ref{simplicial-example-simplex-category}。
以 $(\Delta/[n])_{\leq m}$ 表示 $\Delta/[n]$ 的全子范畴，
其对象为 $\Delta/[n]$ 中满足 $k \leq m$ 的 $[k] \to [n]$。
换言之，有如下范畴与函子的可交换图表：
$$
\xymatrix{
(\Delta/[n])_{\leq m} \ar[r] \ar[d] &
\Delta/[n] \ar[d] \\
\Delta_{\leq m} \ar[r] &
\Delta
}
$$
给定 $\mathcal{C}$ 的 $m$-截断单纯对象 $U$，按下列规则定义函子
$$
U(n) : (\Delta/[n])_{\leq m}^{opp} \longrightarrow \mathcal{C}
$$
：
\begin{eqnarray*}
([k] \to [n]) & \longmapsto & U_k \\
\psi : ([k'] \to [n]) \to ([k] \to [n]) &
\longmapsto &
U(\psi) : U_k \to U_{k'}
\end{eqnarray*}
对 $\Delta$ 的给定态射 $\varphi : [n] \to [n']$，有相应函子
$$
\overline{\varphi} :
(\Delta/[n])_{\leq m}
\longrightarrow
(\Delta/[n'])_{\leq m}
$$
它把 $\alpha : [k] \to [n]$ 送到
$\varphi \circ \alpha : [k] \to [n']$。复合
$U(n') \circ \overline{\varphi}$ 等于函子 $U(n)$。

\begin{lemma}
\label{simplicial-lemma-existence-cosk}
若范畴 $\mathcal{C}$ 有有限极限，则对所有 $m$，函子
$\text{cosk}_m$ 都存在。此外，对任意 $m$-截断单纯对象 $U$，
单纯对象 $\text{cosk}_mU$ 由公式
$$
(\text{cosk}_mU)_n = \lim_{(\Delta/[n])_{\leq m}^{opp}} U(n)
$$
描述；而对 $\varphi : [n] \to [n']$，映射
$\text{cosk}_mU(\varphi)$ 由上面的等同
$U(n') \circ \overline{\varphi} = U(n)$，通过 Categories，引理
\ref{categories-lemma-functorial-limit} 得到。
\end{lemma}

\begin{proof}
在本引理的证明中，以 $\text{cosk}_mU$ 表示单纯对象，其中
$(\text{cosk}_mU)_n$ 等于
$\lim_{(\Delta/[n])_{\leq m}^{opp}} U(n)$。
证明最后将说明，它确实满足所需的映射性质。

\medskip\noindent
设 $V$ 为单纯对象。态射 $\gamma : V \to \text{cosk}_mU$
由一列态射 $\gamma_n : V_n \to (\text{cosk}_mU)_n$ 给出。
由极限的定义，这等价于给出一族态射
$\gamma(\alpha) : V_n \to U_k$，其中 $\alpha$ 取遍所有满足
$k \leq m$ 的 $\alpha : [k] \to [n]$。对任意
$0 \leq k, k' \leq m$、$0 \leq n, n'$ 以及 $\Delta$ 中满足
$\varphi \circ \alpha = \alpha' \circ \psi$ 的
$\psi : [k] \to [k']$、$\varphi : [n] \to [n']$、
$\alpha : [k] \to [n]$ 和 $\alpha' : [k'] \to [n']$，
这些态射还满足图表
$$
\xymatrix{
V_n \ar[r]_{\gamma(\alpha)} &  U_k \\
V_{n'} \ar[r]^{\gamma(\alpha')} \ar[u]^{V(\varphi)} & U_{k'} \ar[u]_{U(\psi)}
}
$$
可交换。取 $n = k = k'$、$\varphi = \alpha'$，以及
$\alpha = \psi = \text{id}_{[k]}$，得到
$\gamma(\alpha') = \gamma(\text{id}_{[k]}) \circ V(\alpha')$。
换言之，态射 $\gamma(\text{id}_{[k]})$（$k \leq m$）确定态射
$\gamma$。容易看出，这些态射组成一个态射
$\text{sk}_m V \to U$。

\medskip\noindent
反过来，给定态射 $\gamma : \text{sk}_m V \to U$，令
$\gamma(\alpha) = \gamma(\text{id}_{[k]}) \circ V(\alpha)$，
便得到一族态射 $\gamma(\alpha)$，其中 $\alpha$ 取遍所有满足
$k \leq m$ 的 $\alpha : [k] \to [n]$。这些态射满足上图所示的全部
可交换条件，因而给出态射 $V \to \text{cosk}_m U$。
\end{proof}

\begin{lemma}
\label{simplicial-lemma-trivial-cosk}
设 $\mathcal{C}$ 为范畴，$U$ 为 $\mathcal{C}$ 的 $m$-截断单纯对象。
对 $n \leq m$，极限
$\lim_{(\Delta/[n])_{\leq m}^{opp}} U(n)$ 存在，并典范同构于 $U_n$。
\end{lemma}

\begin{proof}
这是因为在这种情形下，范畴 $(\Delta/[n])_{\leq m}$ 有一个终对象，
即恒等映射 $[n] \to [n]$。
% P02-E0114: frozen source says “an final object”.
\end{proof}

\begin{lemma}
\label{simplicial-lemma-recover-cosk}
设 $\mathcal{C}$ 为有有限极限的范畴，$U$ 为 $\mathcal{C}$ 的
$n$-截断单纯对象。态射 $\text{sk}_n \text{cosk}_n U \to U$ 是同构。
\end{lemma}

\begin{proof}
合并引理 \ref{simplicial-lemma-existence-cosk} 与引理 \ref{simplicial-lemma-trivial-cosk}。
\end{proof}

\noindent
下面更详细地描述余骨架函子的一个特例。滥用记号，对任意
$n' \geq n$，也以 $\text{sk}_n$ 表示限制函子
$\text{Simp}_{n'}(\mathcal{C}) \to \text{Simp}_n(\mathcal{C})$。
我们将描述函子
$\text{sk}_n : \text{Simp}_{n + 1}(\mathcal{C})
\to \text{Simp}_n(\mathcal{C})$ 的一个右伴随。
对 $n \geq 1$、$0 \leq i < j \leq n + 1$，定义
$\delta^{n + 1}_{i, j} : [n - 1] \to [n + 1]$
为略去 $i$ 和 $j$ 的递增映射。注意
$\delta^{n + 1}_{i, j} =
\delta^{n + 1}_j \circ \delta^n_i =
\delta^{n + 1}_i \circ \delta^n_{j - 1}$；见引理
\ref{simplicial-lemma-relations-face-degeneracy}。这引出下述引理。

\begin{lemma}
\label{simplicial-lemma-formula-limit}
设整数 $n$，并要求其值 $\geq 1$。设 $U$ 为 $\mathcal{C}$ 的
$n$-截断单纯对象。考虑从 $\mathcal{C}$ 到 $\textit{Sets}$ 的反变函子，
它把对象 $T$ 送到集合
$$
\{ (f_0, \ldots, f_{n + 1}) \in \Mor_\mathcal{C}(T, U_n)
\mid
d^n_{j - 1} \circ f_i = d^n_i \circ f_j
\ \forall\ 0\leq i < j\leq n + 1\}
$$
若此函子由 $\mathcal{C}$ 的某个对象 $U_{n + 1}$ 表示，则
$$
U_{n + 1} = \lim_{(\Delta/[n + 1])_{\leq n}^{opp}} U(n)
$$
\end{lemma}

\begin{proof}
该极限若存在，就表示把对象 $T$ 送到集合
$$
\{
(f_\alpha)_{\alpha : [k] \to [n + 1], k \leq n}
\mid
f_{\alpha \circ \psi} = U(\psi) \circ f_\alpha\ \forall
\ \psi : [k'] \to [k], \alpha : [k] \to [n + 1]
\}.
$$
事实上，我们将证明此函子同构于引理中显示的函子。一个方向的映射由规则
$$
(f_\alpha)_{\alpha}
\longmapsto
(f_{\delta^{n + 1}_0}, \ldots, f_{\delta^{n + 1}_{n + 1}}).
$$
给出。由引理前回忆的关系，它满足引理中的条件，因为
$$
d^n_{j - 1} \circ f_{\delta^{n + 1}_i} =
f_{\delta^{n + 1}_i \circ \delta^n_{j - 1}} =
f_{\delta^{n + 1}_j \circ \delta^n_i} =
d^n_i \circ f_{\delta^{n + 1}_j}
$$
为构造反方向的映射，必须把引理中显示公式里的系统
$(f_0, \ldots, f_{n + 1})$ 送到映射系统 $f_\alpha$。
给定 $\alpha : [k] \to [n + 1]$。由于 $k \leq n$，映射 $\alpha$
不是满射。因此，可以写成
$\alpha = \delta^{n + 1}_i \circ \psi$，其中
$0 \leq i \leq n + 1$，且 $\psi : [k] \to [n]$。
我们只能定义
$$
f_\alpha = U(\psi) \circ f_i.
$$
当然，必须检查它与对偶 $(i, \psi)$ 的选择无关。首先注意，给定 $i$，
符合条件的 $\psi$ 唯一。其次，设 $(j, \phi)$ 为另一对。
则 $i \not = j$，并可假设 $i < j$。由于 $i, j$ 都不在
$\alpha$ 的像中，实际上可写
$\alpha = \delta^{n + 1}_{i, j} \circ \xi$；于是
$\psi = \delta^n_{j - 1} \circ \xi$，且
$\phi = \delta^n_i \circ \xi$。因此
\begin{eqnarray*}
U(\psi) \circ f_i & = & U(\delta^n_{j - 1} \circ \xi) \circ f_i \\
& = & U(\xi) \circ d^n_{j - 1} \circ f_i \\
& = & U(\xi) \circ d^n_i \circ f_j \\
& = & U(\delta^n_i \circ \xi) \circ f_j \\
& = & U(\phi) \circ f_j
\end{eqnarray*}
如所欲证。还需验证，如此定义的映射 $f_\alpha$ 满足映射系统
$(f_\alpha)_\alpha$ 的规则。为此，设
$\psi : [k'] \to [k]$、$\alpha : [k] \to [n + 1]$，其中
$k, k' \leq n$。令 $\alpha' = \alpha \circ \psi$。
取不在 $\alpha$ 的像中的 $i$，则 $i$ 显然也不在 $\alpha'$ 的像中。
写 $\alpha = \delta^{n + 1}_i \circ \phi$（这里不能使用字母 $\psi$，
因为它已经用过）。则显然
$\alpha' = \delta^{n + 1}_i \circ \phi \circ \psi$。
由上述构造，有
$$
U(\psi) \circ f_\alpha = U(\psi) \circ U(\phi) \circ f_i
= U(\phi \circ \psi) \circ f_i = f_{\alpha \circ \psi} = f_{\alpha'}
$$
如所欲证。留给读者一项愉快的工作：验证这些构造给出互逆双射，并且关于
$T$ 函子性地成立。
\end{proof}

\begin{lemma}
\label{simplicial-lemma-work-out}
设整数 $n$，并要求其值 $\geq 1$。设 $U$ 为 $\mathcal{C}$ 的
$n$-截断单纯对象。考虑从 $\mathcal{C}$ 到 $\textit{Sets}$ 的反变函子，
它把对象 $T$ 送到集合
$$
\{ (f_0, \ldots, f_{n + 1}) \in \Mor_\mathcal{C}(T, U_n)
\mid
d^n_{j - 1} \circ f_i = d^n_i \circ f_j
\ \forall\ 0\leq i < j\leq n + 1\}
$$
若此函子由 $\mathcal{C}$ 的某个对象 $U_{n + 1}$ 表示，则存在
$(n + 1)$-截断单纯对象 $\tilde U$，满足
$\text{sk}_n \tilde U = U$ 和 $\tilde U_{n + 1} = U_{n + 1}$，
并且有如下伴随关系：
$$
\Mor_{\text{Simp}_{n + 1}(\mathcal{C})}(V, \tilde U)
=
\Mor_{\text{Simp}_n(\mathcal{C})}(\text{sk}_nV, U)
$$
\end{lemma}

\begin{proof}
由引理 \ref{simplicial-lemma-trivial-cosk}，对 $0 \leq i \leq n$ 有等同
$$
U_i = \lim_{(\Delta/[i])_{\leq n}^{opp}} U(i)
$$
由引理 \ref{simplicial-lemma-formula-limit}，有
$$
U_{n + 1} = \lim_{(\Delta/[n + 1])_{\leq n}^{opp}} U(n).
$$
因此，对任意满足 $i, j \leq n + 1$ 的 $\varphi : [i] \to [j]$，
可以完全按照引理 \ref{simplicial-lemma-existence-cosk} 定义相应映射
$\tilde U(\varphi) : \tilde U_j \to \tilde U_i$。
这定义了满足 $\text{sk}_n \tilde U = U$ 的
$(n + 1)$-截断单纯对象 $\tilde U$。

\medskip\noindent
为证明伴随关系，论证如下。给定公式右端的任意元素
$\gamma : \text{sk}_n V \to U$，考虑态射
$f_i = \gamma_n \circ d^{n + 1}_i : V_{n + 1} \to V_n \to U_n$。
它们显然满足关系
$d^n_{j - 1} \circ f_i = d^n_i \circ f_j$，因而由 $U_{n + 1}$
的选取，定义唯一态射 $V_{n + 1} \to U_{n + 1}$。
反之，给定左端的态射 $\gamma' : V \to \tilde U$，只需限制到
$\Delta_{\leq n}$，便得到右端的一个元素。留给读者证明这两个构造互逆。
\end{proof}

\begin{remark}
\label{simplicial-remark-explicit-face-degeneracy}
设 $U$ 和 $U_{n + 1}$ 如引理 \ref{simplicial-lemma-work-out} 中所述。
在 $T$-值点上，可以容易地描述 $\tilde U$ 的面映射和退化映射。
具体地，映射 $d^{n + 1}_i : U_{n + 1} \to U_n$ 由
$$
(f_0, \ldots, f_{n + 1}) \longmapsto f_i.
$$
给出，而映射 $s^n_j : U_n \to U_{n + 1}$ 由
\begin{eqnarray*}
f & \longmapsto & (
s^{n - 1}_{j - 1} \circ d^{n - 1}_0 \circ f, \\
& &
s^{n - 1}_{j - 1} \circ d^{n - 1}_1 \circ f, \\
& &
\ldots\\
& &
s^{n - 1}_{j - 1} \circ d^{n - 1}_{j - 1} \circ f, \\
& &
f, \\
& &
f, \\
& &
s^{n - 1}_j \circ d^{n - 1}_{j + 1} \circ f, \\
& &
s^{n - 1}_j \circ d^{n - 1}_{j + 2} \circ f, \\
& &
\ldots\\
& &
s^{n - 1}_j \circ d^{n - 1}_n \circ f
)
\end{eqnarray*}
给出；留给读者验证，右端是引理 \ref{simplicial-lemma-work-out} 中显示集合的元素。
当 $n = 0$ 时，有一个映射，即 $f \mapsto (f, f)$。
当 $n = 1$ 时，有两个映射，即
$f \mapsto (f, f, s_0d_1f)$ 和
$f \mapsto (s_0d_0f, f, f)$。
当 $n = 2$ 时，有三个映射，即
$f \mapsto (f, f, s_0d_1f, s_0d_2f)$、
$f \mapsto (s_0d_0f, f, f, s_1d_2f)$ 和
$f \mapsto (s_1d_0f, s_1d_1f, f, f)$。依此类推。
\end{remark}


\begin{remark}
\label{simplicial-remark-cosk-simplicial-sets}
上面的引理 \ref{simplicial-lemma-work-out} 在单纯集情形下的构造如下。
给定 $n$-截断单纯集 $U$，如下作出典范的 $(n + 1)$-截断单纯集
$\tilde U$。利用引理中的公式添加一个 $(n + 1)$-单形集
$U_{n + 1}$。也就是说，$U_{n + 1}$ 的一个元素是一列编号的
$n$-单形 $(f_0, \ldots, f_{n + 1})$，它们具有与一个
$(n + 1)$-单形的各面相同的粘合性质。换言之，当 $i < j$ 时，
$f_j$ 的第 $i$ 个面是 $f_i$ 的第 $(j-1)$ 个面。
从几何上看，如何定义 $\tilde U$ 的面映射和退化映射是显然的。
若 $V$ 是 $(n + 1)$-截断单纯集，则它的 $(n + 1)$-单形给出相容的
$n$-单形族 $(f_0, \ldots, f_{n + 1})$，其中 $f_i \in V_n$。
因此有自然映射
$\Mor(\text{sk}_nV, U) \to \Mor(V, \tilde U)$，
它与反方向的典范限制映射互逆。

\medskip\noindent
此外，只需在截断单纯集的情形下完成该构造的组合论部分。
事实上，对范畴 $\mathcal{C}$ 的任意对象 $T$ 和 $\mathcal{C}$ 的任意
$n$-截断单纯对象 $U$，可以考虑 $n$-截断单纯集 $\Mor(T, U)$。
可以对此应用上述构造，取其 $(n + 1)$-单形集，并要求它可表示。
这是理解引理 \ref{simplicial-lemma-work-out} 结论的一种好方法。
\end{remark}

\begin{remark}
\label{simplicial-remark-inductive-coskeleton}
{\it 余骨架的归纳构造。}
假设 $\mathcal{C}$ 是有有限极限的范畴，$U$ 是 $\mathcal{C}$ 中的
$m$-截断单纯对象。可以如下归纳构造 $n$-截断对象 $U^n$：
\begin{enumerate}
\item 从 $U^m = U$ 开始。
\item 给定 $n \geq m$ 时的 $U^n$，令 $U^{n + 1} = \tilde U^n$，
其中 $\tilde U^n$ 按引理 \ref{simplicial-lemma-work-out} 由 $U^n$ 构造。
\end{enumerate}
由于引理 \ref{simplicial-lemma-work-out} 的构造保持 $U^n$ 的 $n$-骨架不变，
可以把 $\text{cosk}_m U$ 定义为满足
$(\text{cosk}_m U)_n = U^n_n = U^{n + 1}_n = \ldots$ 的单纯对象。
而且由引理 \ref{simplicial-lemma-work-out} 形式地推出，对所有 $n \geq m$，
$U^n$ 满足公式
$$
\Mor_{\text{Simp}_n(\mathcal{C})}(V, U^n)
=
\Mor_{\text{Simp}_m(\mathcal{C})}(\text{sk}_mV, U)
$$
进而又形式地推出
$$
\Mor_{\text{Simp}(\mathcal{C})}(V, \text{cosk}_mU)
=
\Mor_{\text{Simp}_m(\mathcal{C})}(\text{sk}_mV, U)
$$
其中 $\text{cosk}_mU$ 按上述方式选取。
\end{remark}

\begin{lemma}
\label{simplicial-lemma-cosk-up}
设 $\mathcal{C}$ 为有有限极限的范畴。
\begin{enumerate}
\item 对每个 $n$，函子 $\text{sk}_n : \text{Simp}(\mathcal{C})
\to \text{Simp}_n(\mathcal{C})$ 有右伴随 $\text{cosk}_n$。
\item 对每个 $n' \geq n$，函子
$\text{sk}_n : \text{Simp}_{n'}(\mathcal{C}) \to \text{Simp}_n(\mathcal{C})$
有右伴随，即 $\text{sk}_{n'}\text{cosk}_n$。
\item 对每个 $m \geq n \geq 0$ 以及 $\mathcal{C}$ 的每个
$n$-截断单纯对象 $U$，有
$\text{cosk}_m \text{sk}_m \text{cosk}_n U = \text{cosk}_n U$。
\item 若 $U$ 是 $\mathcal{C}$ 的单纯对象，并且对某个 $n \geq 0$，
典范映射 $U \to \text{cosk}_n \text{sk}_nU$ 是同构，则对所有
$m \geq n$，典范映射 $U \to \text{cosk}_m \text{sk}_mU$ 都是同构。
\end{enumerate}
\end{lemma}

\begin{proof}
(1) 中的存在性由上面的引理 \ref{simplicial-lemma-existence-cosk} 推出。
(2) 和 (3) 由备注 \ref{simplicial-remark-inductive-coskeleton} 中的讨论推出。
于是 (4) 显然。
\end{proof}

\begin{remark}
\label{simplicial-remark-existence-cosk}
定义余骨架函子并不需要所有有限极限。说明如下：
\begin{enumerate}
\item 在例 \ref{simplicial-example-cosk0} 中已经看到，若 $\mathcal{C}$ 有对象对的积，
则 $\text{cosk}_0$ 存在。
\item 对 $k > 0$，若 $\mathcal{C}$ 有有限连通极限，则函子
$\text{cosk}_k$ 存在。
\end{enumerate}
这由余骨架的归纳构造步骤（备注 \ref{simplicial-remark-cosk-simplicial-sets} 和
\ref{simplicial-remark-inductive-coskeleton}）清楚可见；也可以由下述事实推出：
引理 \ref{simplicial-lemma-existence-cosk} 中使用的范畴
$(\Delta/[n])_{\leq k}$ 在 $k \geq 1$ 且 $n \geq k + 1$ 时是连通的。
注意，由引理 \ref{simplicial-lemma-trivial-cosk} 或引理 \ref{simplicial-lemma-recover-cosk}，
我们不需要 $n \leq k$ 时的那些范畴。（随着 $k$ 增大，
在 $k \geq 1$ 且 $n \geq k + 1$ 时，范畴
$(\Delta/[n])_{\leq k}$ 在拓扑意义下越来越连通。）
\end{remark}

\begin{lemma}
\label{simplicial-lemma-cosk-product}
设 $U$、$V$ 为范畴 $\mathcal{C}$ 的 $n$-截断单纯对象。只要等式两端
都存在，就有
$$
\text{cosk}_n (U \times V) = \text{cosk}_nU \times \text{cosk}_nV
$$
\end{lemma}

\begin{proof}
设 $W$ 为单纯对象。有
\begin{eqnarray*}
\Mor(W, \text{cosk}_n (U \times V))
& = &
\Mor(\text{sk}_n W, U \times V) \\
& = &
\Mor(\text{sk}_n W, U)
\times
\Mor(\text{sk}_nW, V) \\
& = &
\Mor(W, \text{cosk}_n U)
\times
\Mor(W, \text{cosk}_n V) \\
& = &
\Mor(W, \text{cosk}_n U \times \text{cosk}_n V)
\end{eqnarray*}
引理得证。
\end{proof}

\begin{lemma}
\label{simplicial-lemma-cosk-fibre-product}
假设 $\mathcal{C}$ 有纤维积。设 $U \to V$ 和 $W \to V$ 为
范畴 $\mathcal{C}$ 的 $n$-截断单纯对象之间的态射。只要等式两端存在，
就有
$$
\text{cosk}_n (U \times_V W)
=
\text{cosk}_nU \times_{\text{cosk}_n V} \text{cosk}_nW
$$
\end{lemma}

\begin{proof}
略；证明与上面的引理 \ref{simplicial-lemma-cosk-product} 的证明非常相似。
\end{proof}

\begin{lemma}
\label{simplicial-lemma-cosk-above-object}
设 $\mathcal{C}$ 为有有限极限的范畴，并设
$X \in \Ob(\mathcal{C})$。对 $k \geq 1$，函子
$\mathcal{C}/X \to \mathcal{C}$ 与余骨架函子 $\text{cosk}_k$ 交换。
\end{lemma}

\begin{proof}
该断言的意思是：若 $U$ 是 $\mathcal{C}/X$ 的单纯对象，
可把它看成 $\mathcal{C}$ 的一个单纯对象，连同一个指向常值单纯对象
$X$ 的态射；那么在 $\mathcal{C}/X$ 中计算的 $\text{cosk}_k U$
与在 $\mathcal{C}$ 中计算的相同。例如，这由 Categories，引理
\ref{categories-lemma-connected-limit-over-X} 推出，因为引理
\ref{simplicial-lemma-existence-cosk} 中使用的范畴 $(\Delta/[n])_{\leq k}$
在 $k \geq 1$ 且 $n \geq k + 1$ 时是连通的。
注意，由引理 \ref{simplicial-lemma-trivial-cosk} 或引理 \ref{simplicial-lemma-recover-cosk}，
我们不需要 $n \leq k$ 时的那些范畴。
\end{proof}

\begin{lemma}
\label{simplicial-lemma-simplex-cosk}
典范映射 $\Delta[n] \to \text{cosk}_1 \text{sk}_1 \Delta[n]$ 是同构。
\end{lemma}

\begin{proof}
考虑单纯集 $U$ 和态射 $f : U \to \Delta[n]$。这是一条规则：
它把每个 $u \in U_i$ 关联到 $\Delta$ 中的映射
$f_u : [i] \to [n]$。此外，这些映射应满足：对任意
$\varphi : [j] \to [i]$，有
$f_u \circ \varphi = f_{U(\varphi)(u)}$。
以 $\epsilon^i_j : [0] \to [i]$ 表示把 $0$ 映到 $j$ 的映射。
以 $F : U_0 \to [n]$ 表示映射 $u \mapsto f_u(0)$。于是可见
$$
f_u(j) = F(\epsilon^i_j(u))
$$
对所有 $0 \leq j \leq i$ 和 $u \in U_i$ 成立。特别地，若知道函数
$F$，则对所有 $u\in U_i$ 和所有 $i$，都知道映射 $f_u$。
反之，给定映射 $F : U_0 \to [n]$，可以对任意 $i$、
任意 $u \in U_i$ 和任意 $0 \leq j \leq i$ 令
$$
f_u(j) = F(\epsilon^i_j(u))
$$
一般而言，这并不定义上面的单纯集态射 $f$。条件是所有映射
$f_u$ 都非降。这显然等价于：当 $0 \leq j \leq j' \leq i$
且 $u \in U_i$ 时，
$F(\epsilon^i_j(u)) \leq F(\epsilon^i_{j'}(u))$。
但在这种情形下，态射
$$
\epsilon^i_j, \epsilon^i_{j'} : [0] \to [i]
$$
都通过映射 $\epsilon^i_{j, j'} : [1] \to [i]$ 分解；后者由规则
$0 \mapsto j$、$1 \mapsto j'$ 定义。换言之，只需对 $i = 1$
和 $u \in X_1$ 检查这些不等式。
% P02-E0115: frozen source switches from the simplicial set U to undefined X_1.
也就是说，有
$$
\Mor(U, \Delta[n])
=
\Mor(\text{sk}_1 U, \text{sk}_1 \Delta[n])
$$
如所欲证。
\end{proof}



\section{增广}
\label{simplicial-section-augmentation}

\begin{definition}
\label{simplicial-definition-augmentation}
设 $\mathcal{C}$ 为范畴，$U$ 为 $\mathcal{C}$ 的单纯对象。
$U$ 向 $\mathcal{C}$ 的对象 $X$ 的一个{\it 增广
$\epsilon : U \to X$}，是从 $U$ 到常值单纯对象 $X$ 的态射。
\end{definition}

\begin{lemma}
\label{simplicial-lemma-augmentation-howto}
设 $\mathcal{C}$ 为范畴，$X \in \Ob(\mathcal{C})$，并设 $U$ 为
$\mathcal{C}$ 的单纯对象。给出 $U$ 向 $X$ 的一个增广，等价于给出
态射 $\epsilon_0 : U_0 \to X$，使得
$\epsilon_0 \circ d^1_0 = \epsilon_0 \circ d^1_1$。
\end{lemma}

\begin{proof}
给定态射 $\epsilon : U \to X$，当然得到引理中所述的
$\epsilon_0$。反之，给定引理中所述的 $\epsilon_0$，任选态射
$\alpha : [0] \to [n]$，并令
$\epsilon_n = \epsilon_0 \circ U(\alpha)$，从而定义
$\epsilon_n : U_n \to X$。事实上，若
$\beta : [0] \to [n]$ 是另一种选择，则存在态射
$\gamma : [1] \to [n]$，使得 $\alpha$ 和 $\beta$ 都分解为
$[0] \to [1] \to [n]$。因此，对 $\epsilon_0$ 的条件说明
$\epsilon_n$ 是良定义的。容易证明，
$(\epsilon_n) : U \to X$ 是单纯对象之间的态射。
\end{proof}

\begin{lemma}
\label{simplicial-lemma-cosk-minus-one}
设 $\mathcal{C}$ 为有纤维积的范畴。设 $f : Y\to X$ 为
$\mathcal{C}$ 的态射，并设 $U$ 为 $\mathcal{C}$ 的单纯对象，
其第 $n$ 项是 $(n + 1)$ 重纤维积
$Y \times_X Y \times_X \ldots \times_X Y$。
% P02-E0116: frozen source has “(n + 1)fold” without a hyphen.
见例 \ref{simplicial-example-fibre-products-simplicial-object}。
对 $\mathcal{C}$ 的任意单纯对象 $V$，有
\begin{align*}
\Mor_{\text{Simp}(\mathcal{C})}(V, U)
& =
\Mor_{\text{Simp}_1(\mathcal{C})}(\text{sk}_1 V, \text{sk}_1 U) \\
& =
\{g_0 : V_0 \to Y \mid f \circ g_0 \circ d^1_0 = f \circ g_0 \circ d^1_1\}
\end{align*}
特别地，有 $U = \text{cosk}_1 \text{sk}_1 U$。
\end{lemma}

\begin{proof}
设 $g : \text{sk}_1V \to \text{sk}_1U$ 为 $1$-截断单纯对象之间的态射。
则图表
$$
\xymatrix{
V_1 \ar@<1ex>[r]^{d^1_0} \ar@<-1ex>[r]_{d^1_1} \ar[d]_{g_1} &
V_0 \ar[d]^{g_0} \\
Y \times_X Y \ar@<1ex>[r]^{pr_1} \ar@<-1ex>[r]_{pr_0} &
Y \ar[r] & X
}
$$
可交换，这证明了引理中显示的关系。反过来，必须证明：给定满足关系
$f \circ g_0 \circ d^1_0 = f \circ g_0 \circ d^1_1$ 的态射
$g_0$，可得到唯一的单纯对象态射 $g : V \to U$。
作法如下。对任意 $n \geq 1$，令
$g_{n, i} = g_0 \circ V([0] \to [n], 0 \mapsto i) :
V_n \to Y$。由可交换图表
$$
\xymatrix{
[0] \ar[rd]_{0 \mapsto 0} \ar[rrrrrd]^{0 \mapsto i} \\
& [1] \ar[rrrr]^{0 \mapsto i, 1\mapsto i + 1} & & & & [n] \\
[0] \ar[ru]^{0 \mapsto 1} \ar[rrrrru]_{0 \mapsto i + 1}
}
$$
上述等式蕴含 $f \circ g_{n, i} = f \circ g_{n, i + 1}$。
因此得到
$(g_{n, 0}, \ldots, g_{n, n}) : V_n \to Y \times_X\ldots \times_X Y = U_n$。
留给读者验证，这是单纯对象之间的态射。引理最后一个断言等价于引理中
显示公式的第一个等式。
\end{proof}

\begin{remark}
\label{simplicial-remark-augmentation}
设 $\mathcal{C}$ 为有纤维积的范畴，$V$ 为单纯对象，
$\epsilon : V \to X$ 为增广。设 $U$ 为如下单纯对象：其第 $n$ 项
是 $V_0$ 在 $X$ 上的第 $(n + 1)$ 个纤维积。
% P02-E0117: frozen source says “(n + 1)st fibred product”; the construction is the (n + 1)-fold fibre product.
简单地结合引理 \ref{simplicial-lemma-augmentation-howto} 与引理
\ref{simplicial-lemma-cosk-minus-one}，得到典范态射 $V \to U$。
\end{remark}





\section{骨架函子的左伴随}
\label{simplicial-section-adjoint-left}

\noindent
本节在某些情形下构造骨架函子 $\text{sk}_m$ 的左伴随 $i_{m!}$。
伴随公式为
$$
\Mor_{\text{Simp}_m(\mathcal{C})}(U, \text{sk}_mV)
=
\Mor_{\text{Simp}(\mathcal{C})}(i_{m!}U, V).
$$
事实表明，当范畴 $\mathcal{C}$ 有有限余极限时，这个左伴随存在。

\medskip\noindent
我们使用与第 \ref{simplicial-section-skeleton} 节相似的构造。回忆 $[n]$ 下对象
所成的范畴 $[n]/\Delta$，见 Categories，例
\ref{categories-example-category-under-X}。其对象是态射
$\alpha : [n] \to [k]$，其态射是可交换三角形。
以 $([n]/\Delta)_{\leq m}$ 表示 $[n]/\Delta$ 的全子范畴，
由满足 $k \leq m$ 的对象 $[n] \to [k]$ 组成。
给定 $\mathcal{C}$ 的 $m$-截断单纯对象 $U$，按下列规则定义函子
$$
U(n) : ([n]/\Delta)_{\leq m}^{opp} \longrightarrow \mathcal{C}
$$
：
\begin{eqnarray*}
([n] \to [k]) & \longmapsto & U_k \\
\psi : ([n] \to [k']) \to ([n] \to [k])
& \longmapsto &
U(\psi) : U_k \to U_{k'}
\end{eqnarray*}
对 $\Delta$ 的给定态射 $\varphi : [n] \to [n']$，有相应函子
$$
\underline{\varphi} :
([n']/\Delta)_{\leq m}
\longrightarrow
([n]/\Delta)_{\leq m}
$$
它把 $\alpha : [n'] \to [k]$ 送到
$\varphi \circ \alpha : [n] \to [k]$。复合
$U(n) \circ \underline{\varphi}$ 等于函子 $U(n')$。

\begin{lemma}
\label{simplicial-lemma-left-adjoint-exists}
设 $\mathcal{C}$ 为有有限余极限的范畴。对所有 $m$，函子 $i_{m!}$
都存在。设 $U$ 为 $\mathcal{C}$ 的 $m$-截断单纯对象。
单纯对象 $i_{m!}U$ 由公式
$$
(i_{m!}U)_n = \colim_{([n]/\Delta)_{\leq m}^{opp}} U(n)
$$
描述；而对 $\varphi : [n] \to [n']$，映射 $i_{m!}U(\varphi)$
由上面的等同 $U(n) \circ \underline{\varphi} = U(n')$，通过
Categories，引理 \ref{categories-lemma-functorial-colimit} 得到。
\end{lemma}

\begin{proof}
在本证明中，以 $i_{m!}U$ 表示其第 $n$ 项由引理中显示公式给出的
单纯对象。我们将证明它满足伴随性质。

\medskip\noindent
设 $V$ 为 $\mathcal{C}$ 的单纯对象，并给定
$\gamma : U \to \text{sk}_mV$。态射
$$
\colim_{([n]/\Delta)_{\leq m}^{opp}} U(n) \to T
$$
由一个相容的态射系统 $f_\alpha : U_k \to T$ 给出，其中
$\alpha : [n] \to [k]$ 且 $k \leq m$。当然，取复合
$$
U_k \xrightarrow{\gamma_k} V_k \xrightarrow{V(\alpha)} V_n.
$$
便得到这样的态射系统。因此得到诱导态射
$(i_{m!}U)_n \to V_n$。留给读者验证，它们组成单纯对象之间的态射
$\gamma' : i_{m!}U \to V$。

\medskip\noindent
反过来，给定态射 $\gamma' : i_{m!}U \to V$，对
$0 \leq i \leq n$，令 $\gamma_i : U_i \to V_i$ 等于复合
$$
U_i
\xrightarrow{\text{id}_{[i]}}
\colim_{([i]/\Delta)_{\leq m}^{opp}} U(i)
\xrightarrow{\gamma'_i}
V_i
$$
% P02-E0118: frozen source bounds i by n here, but the reconstructed morphism has one component for every i <= m.
从而得到态射 $\gamma : U \to \text{sk}_m V$。留给读者验证，
这是上述构造的逆。
\end{proof}

\begin{lemma}
\label{simplicial-lemma-recovering-U}
设 $\mathcal{C}$ 为范畴，$U$ 为 $\mathcal{C}$ 的 $m$-截断单纯对象。
对任意 $n \leq m$，余极限
$$
\colim_{([n]/\Delta)_{\leq m}^{opp}} U(n)
$$
存在，并等于 $U_n$。
\end{lemma}

\begin{proof}
这是因为范畴 $([n]/\Delta)_{\leq m}$ 有初对象，即
$\text{id} : [n] \to [n]$。
\end{proof}

\begin{lemma}
\label{simplicial-lemma-recovering-U-for-real}
设 $\mathcal{C}$ 为有有限余极限的范畴，$U$ 为 $\mathcal{C}$ 的
$m$-截断单纯对象。映射 $U \to \text{sk}_m i_{m!}U$ 是同构。
\end{lemma}

\begin{proof}
合并引理 \ref{simplicial-lemma-left-adjoint-exists} 与引理 \ref{simplicial-lemma-recovering-U}。
\end{proof}

\begin{lemma}
\label{simplicial-lemma-imshriek-sets}
若 $U$ 是 $m$-截断单纯集且 $n > m$，则 $i_{m!}U$ 的所有
$n$-单形都是退化的。
\end{lemma}

\begin{proof}
这可以从引理 \ref{simplicial-lemma-left-adjoint-exists} 中 $i_{m!}U$ 的构造看出，
也可如下直接论证。令 $V = i_{m!}U$。设 $V' \subset V$ 为如下
单纯子集：当 $i \leq m$ 时 $V'_i = V_i$，而当 $i > m$ 时所有
$i$-单形都退化；见引理 \ref{simplicial-lemma-simplicial-set-n-skel-sub}。
由伴随公式，又因为 $\text{sk}_m V' = U$，单射 $V' \to V$ 有逆。
因此 $V' = V$。
\end{proof}

\begin{lemma}
\label{simplicial-lemma-n-skeleton-sets}
设 $U$ 为单纯集，$n \geq 0$ 为整数。态射
$i_{n!} \text{sk}_n U \to U$ 把 $i_{n!} \text{sk}_n U$ 与引理
\ref{simplicial-lemma-simplicial-set-n-skel-sub} 中定义的单纯集
$U' \subset U$ 等同起来。
\end{lemma}

\begin{proof}
由引理 \ref{simplicial-lemma-imshriek-sets}，$i_{n!} \text{sk}_n U$ 的非退化单形
只出现在次数 $\leq n$。映射 $i_{n!} \text{sk}_n U \to U$
在次数 $\leq n$ 上是同构。综合这两点，得到映射
$i_{n!} \text{sk}_n U \to U$ 把非退化单形映到非退化单形，
并且任意两个非退化单形的像都不相同。因此可应用引理
\ref{simplicial-lemma-injective-map-simplicial-sets}。于是
$i_{n!} \text{sk}_n U \to U$ 是单射。结论容易由此推出。
\end{proof}

\begin{remark}
\label{simplicial-remark-sk-literature}
在一些文献中，复合函子
$$
\text{Simp}(\mathcal{C})
\xrightarrow{\text{sk}_m}
\text{Simp}_m(\mathcal{C})
\xrightarrow{i_{m!}}
\text{Simp}(\mathcal{C})
$$
记作 $\text{sk}_m$。对单纯集而言这是合理的，因为此时引理
\ref{simplicial-lemma-n-skeleton-sets} 表明，$i_{m!} \text{sk}_m V$ 正是
$V$ 的如下子单纯集：它由 $V$ 的所有 $i$-单形（$i \leq m$）
及其退化组成。在这些文献中，通常还把复合
$$
\text{Simp}(\mathcal{C})
\xrightarrow{\text{sk}_m}
\text{Simp}_m(\mathcal{C})
\xrightarrow{\text{cosk}_m}
\text{Simp}(\mathcal{C})
$$
记作 $\text{cosk}_m$。
\end{remark}

\begin{lemma}
\label{simplicial-lemma-glue-simplex}
设 $U \subset V$ 为单纯集。假设 $n \geq 0$，且
$x \in V_n$、$x \not \in U_n$ 满足：
\begin{enumerate}
\item 当 $i < n$ 时，$V_i = U_i$；
\item $V_n = U_n \cup \{x\}$；
\item 当 $j > n$ 时，任意 $z \in V_j$、$z \not \in U_j$ 都退化。
\end{enumerate}
设 $\Delta[n] \to V$ 为把 $\Delta[n]$ 的非退化 $n$-单形映到
$x$ 的唯一态射。此时图表
$$
\xymatrix{
\Delta[n] \ar[r] & V \\
i_{(n - 1)!} \text{sk}_{n - 1} \Delta[n] \ar[r] \ar[u] & U \ar[u]
}
$$
是推出图表。
\end{lemma}

\begin{proof}
为方便起见，记
$\partial \Delta[n] = i_{(n - 1)!} \text{sk}_{n - 1} \Delta[n]$。
有自然映射 $U \amalg_{\partial \Delta[n]} \Delta[n] \to V$。
必须证明它在每个次数 $j$ 上都是双射，并且这对所有 $j$ 都成立。当 $j \leq n$ 时显然。
设 $j > n$。第三个条件意味着，任意 $z \in V_j$、
$z \not \in U_j$ 都是退化单形，例如
$z = s^{j - 1}_i(z')$。当然，$z' \not \in U_{j - 1}$。
由归纳可知，$z'$ 是 $x$ 的一个退化。因此，$V$ 的所有 $j$-单形
或者属于 $U$，或者是 $x$ 的退化。这蕴含映射
$U \amalg_{\partial \Delta[n]} \Delta[n] \to V$ 是满射。
注意，$U \amalg_{\partial \Delta[n]} \Delta[n]$ 的非退化单形，
或者是 $U$ 的某个非退化单形的像，或者是 $\Delta[n]$ 的（唯一）
非退化 $n$-单形的像。由于 $x$ 显然非退化，可推出
$U \amalg_{\partial \Delta[n]} \Delta[n] \to V$ 把非退化单形映到
非退化单形，并在非退化单形上是单射。因此，由引理
\ref{simplicial-lemma-injective-map-simplicial-sets}，它是单射。
\end{proof}

\begin{lemma}
\label{simplicial-lemma-add-simplices}
设 $U \subset V$ 为单纯集，并且对所有 $n$，$U_n, V_n$ 都有限非空。
假设 $U$ 和 $V$ 只有有限多个非退化单形。则存在子单纯集序列
$$
U = W^0 \subset W^1 \subset W^2 \subset \ldots W^r = V
$$
% P02-E0119: frozen source omits a subset relation between the ellipsis and W^r.
使得引理 \ref{simplicial-lemma-glue-simplex} 可应用于每个包含
$W^i \subset W^{i + 1}$。
\end{lemma}

\begin{proof}
设 $n$ 为使 $V$ 有一个不属于 $U$ 的非退化单形的最小整数。
取这样的非退化单形 $x \in V_n$、$x\not \in U_n$。
设 $W \subset V$ 为如下元素的集合：这些元素或者属于 $U$，
或者是 $x$ 的（反复）退化（换言之，具有形式 $V(\varphi)(x)$，
其中 $\varphi : [m] \to [n]$ 是满射）。容易看出，$W$ 是单纯集。
包含 $U \subset W$ 满足引理 \ref{simplicial-lemma-glue-simplex} 的条件。
此外，$V$ 中不属于 $W$ 的非退化单形数，恰比 $V$ 中不属于 $U$
的非退化单形数少一。因此，对这个数归纳即可。
\end{proof}

\begin{lemma}
\label{simplicial-lemma-imshriek-abelian}
设 $\mathcal{A}$ 为阿贝尔范畴
% P02-E0120: frozen source lacks sentence punctuation after “abelian category”.
设 $U$ 为 $\mathcal{A}$ 的 $m$-截断单纯对象。对 $n > m$，有
$N(i_{m!}U)_n = 0$。
\end{lemma}

\begin{proof}
令 $V = i_{m!}U$。设 $V' \subset V$ 为 $V$ 的如下单纯子对象：
当 $i \leq m$ 时 $V'_i = V_i$，而当 $i > m$ 时
$N(V'_i) = 0$；见引理 \ref{simplicial-lemma-simplicial-abelian-n-skel-sub}。
由伴随公式，又因为 $\text{sk}_m V' = U$，单射 $V' \to V$ 有逆。
因此 $V' = V$。
\end{proof}

\begin{lemma}
\label{simplicial-lemma-n-skeleton-abelian}
设 $\mathcal{A}$ 为阿贝尔范畴，$U$ 为 $\mathcal{A}$ 的单纯对象，
$n \geq 0$ 为整数。态射 $i_{n!} \text{sk}_n U \to U$
把 $i_{n!} \text{sk}_n U$ 与引理
\ref{simplicial-lemma-simplicial-abelian-n-skel-sub} 中定义的单纯子对象
$U' \subset U$ 等同起来。
\end{lemma}

\begin{proof}
由引理 \ref{simplicial-lemma-imshriek-abelian}，当 $i > n$ 时，有
$N(i_{n!} \text{sk}_n U)_i = 0$。映射
$i_{n!} \text{sk}_n U \to U$ 在次数 $\leq n$ 上是同构，
见引理 \ref{simplicial-lemma-recovering-U-for-real}。综合这两点，得到映射
$i_{n!} \text{sk}_n U \to U$ 对所有 $i$ 都诱导单射
$N(i_{n!} \text{sk}_n U)_i \to N(U)_i$。
因此可应用引理 \ref{simplicial-lemma-injective-map-simplicial-abelian}。
于是 $i_{n!} \text{sk}_n U \to U$ 是单射。结论容易由此推出。
\end{proof}

\noindent
下面利用上述材料给出理解余骨架函子的另一种方式。

\begin{lemma}
\label{simplicial-lemma-cosk-shriek}
设 $\mathcal{C}$ 为有有限余积和有限极限的范畴，$V$ 为
$\mathcal{C}$ 的单纯对象。此时
$$
(\text{cosk}_n \text{sk}_n V)_{n + 1}
=
\Hom(i_{n !}\text{sk}_n \Delta[n + 1], V)_0.
$$
\end{lemma}

\begin{proof}
由引理 \ref{simplicial-lemma-morphism-from-coproduct}，左边的对象表示把 $X$
送到下列等式中第一个集合的函子：
\begin{eqnarray*}
\Mor(X \times \Delta[n + 1], \text{cosk}_n \text{sk}_n V)
& = &
\Mor(X \times \text{sk}_n \Delta[n + 1], \text{sk}_n V) \\
& = &
\Mor(X \times i_{n !} \text{sk}_n \Delta[n + 1], V).
\end{eqnarray*}
引理公式右边的对象由把 $X$ 送到该等式链中最后一个集合的函子表示。
这证明了结论。

\medskip\noindent
在该等式链中，我们使用了
$\text{sk}_n (X \times \Delta[n + 1]) = X \times \text{sk}_n \Delta[n + 1]$
以及
$i_{n!}(X \times \text{sk}_n \Delta[n + 1]) =
X \times i_{n !} \text{sk}_n \Delta[n + 1]$。
第一个等式显然。对 $\mathcal{C}$ 的任意（可能截断的）单纯对象 $W$
和 $\mathcal{C}$ 的任意对象 $X$，暂以 $\Mor_\mathcal{C}(X, W)$ 表示（可能截断的）
单纯集 $[n] \mapsto \Mor_\mathcal{C}(X, W_n)$。
由定义，对任意（可能截断的）单纯集 $U$，有
$\Mor(U \times X, W) =
\Mor(U, \Mor_\mathcal{C}(X, W))$。因此
\begin{eqnarray*}
\Mor(X \times i_{n !} \text{sk}_n \Delta[n + 1], W)
& = &
\Mor(i_{n !} \text{sk}_n \Delta[n + 1], \Mor_\mathcal{C}(X, W)) \\
& = &
\Mor(\text{sk}_n \Delta[n + 1],
\text{sk}_n\Mor_\mathcal{C}(X, W)) \\
& = &
\Mor(X \times \text{sk}_n \Delta[n + 1], \text{sk}_nW) \\
& = &
\Mor(i_{n!}(X \times \text{sk}_n \Delta[n + 1]), W).
\end{eqnarray*}
这证明了所用的第二个等式，也结束了引理的证明。
\end{proof}

\section{阿贝尔范畴中的单纯对象}
\label{simplicial-section-abelian}

\noindent
回顾阿贝尔范畴的定义见《同调代数》第
\ref{homology-section-abelian-categories} 节。

\begin{lemma}
\label{simplicial-lemma-abelian}
设 $\mathcal{A}$ 为阿贝尔范畴。
\begin{enumerate}
\item 范畴 $\text{Simp}(\mathcal{A})$ 和
$\text{CoSimp}(\mathcal{A})$ 都是阿贝尔范畴。
\item （余）单纯对象的态射 $f : A \to B$ 为单射，当且仅当每个
$f_n : A_n \to B_n$ 都为单射。
\item （余）单纯对象的态射 $f : A \to B$ 为满射，当且仅当每个
$f_n : A_n \to B_n$ 都为满射。
\item （余）单纯对象的序列
$$
A \xrightarrow{f} B \xrightarrow{g} C
$$
在 $B$ 处正合，当且仅当每个序列
$$
A_i \xrightarrow{f_i} B_i \xrightarrow{g_i} C_i
$$
都在 $B_i$ 处正合。
\end{enumerate}
\end{lemma}

\begin{proof}
预加性很容易验证。终对象由所有次数中的 $U_n = 0$ 给出。
直积的存在性已见引理 \ref{simplicial-lemma-product} 和
\ref{simplicial-lemma-product-cosimplicial-objects}。
核与余核可逐项取核与余核得到。
\end{proof}

\noindent
对 $\mathcal{A}$ 的对象 $A$ 和整数 $k$，考虑满足下列条件的
$k$-截断单纯对象 $U$：
\begin{enumerate}
\item 当 $i < k$ 时，$U_i = 0$；
\item $U_k = A$；
\item 所有态射 $U(\varphi)$ 都等于零，唯有
$U(\text{id}_{[k]}) = \text{id}_A$。
\end{enumerate}
由于 $\mathcal{A}$ 既有有限极限也有有限余极限，
$\text{cosk}_k U$ 和 $i_{k!}U$ 都存在。
我们将描述二者以及典范映射 $i_{k!}U \to \text{cosk}_kU$。

\begin{lemma}
\label{simplicial-lemma-eilenberg-maclane-object}
设 $A$、$k$ 和 $U$ 如上，即当 $i < k$ 时 $U_i = 0$，且 $U_k = A$。
\begin{enumerate}
\item 给定 $k$-截断单纯对象 $V$，有
$$
\Mor(U, V)
=
\{ f : A \to V_k \mid d^k_i \circ f = 0, \ i = 0, \ldots, k \}
$$
以及
$$
\Mor(V, U)
=
\{ f : V_k \to A \mid f \circ s^{k - 1}_i = 0, \ i = 0, \ldots, k - 1 \}.
$$
\item 对象 $i_{k!} U$ 的第 $n$ 项等于
$\bigoplus_\alpha A$，其中 $\alpha$ 遍历所有满射
$\alpha : [n] \to [k]$。
\item 对任意 $\varphi : [m] \to [n]$，映射
$i_{k!} U(\varphi)$ 可描述为映射
$\bigoplus_\alpha A \to \bigoplus_{\alpha'} A$：
若 $\alpha \circ \varphi$ 不是满射，则它把与
$\alpha : [n] \to [k]$ 对应的分量映为零；若其为满射，
则以恒等映射把该分量映到与 $\alpha \circ \varphi$ 对应的分量。
\item 对象 $\text{cosk}_k U$ 的第 $n$ 项等于
$\bigoplus_\beta A$，其中 $\beta$ 遍历所有单射
$\beta : [k] \to [n]$。
\item 对任意 $\varphi : [m] \to [n]$，映射
$\text{cosk}_k U(\varphi)$ 可描述为映射
$\bigoplus_\beta A \to \bigoplus_{\beta'} A$：
若 $\beta$ 不经由 $\varphi$ 分解，则它把与
$\beta : [k] \to [n]$ 对应的分量映为零；若能够分解，
则以恒等映射把该分量映到每个满足
$\beta = \varphi \circ \beta'$ 的 $\beta'$ 所对应的分量。
\item 典范映射
$
c : i_{k !} U \to \text{cosk}_k U
$
在次数 $n$ 上的 $(\alpha, \beta)$ 系数 $A \to A$：
若 $\alpha \circ \beta$ 不是恒等映射，则等于零；
若其为恒等映射，则等于 $\text{id}_A$。
\item 典范映射
$
c : i_{k !} U \to \text{cosk}_k U
$
为单射。
\end{enumerate}
\end{lemma}

\begin{proof}
(1) 的证明留给读者。

\medskip\noindent
我们把 (2) 和 (3) 的规则作为一个单纯对象的定义，记为
$\tilde U$。下面证明它是 $i_{k!}U$ 的一种实现。
这将同时证明 (2) 和 (3)。必须证明：给定态射
$f : U \to \text{sk}_kV$，存在唯一态射
$\tilde f : \tilde U \to V$，使其取 $k$-骨架后恢复为 $f$。
由 (1)，$f$ 对应于态射 $f_k : A \to V_k$，且它对所有 $i$
都落在 $d^k_i$ 的核中。对任意满射
$\alpha : [n] \to [k]$，令
$\tilde f_\alpha : A \to V_n$ 等于复合
$\tilde f_\alpha = V(\alpha) \circ f_k : A \to V_n$。
定义 $\tilde f_n : \tilde U_n \to V_n$ 为所有满射
$\alpha : [n] \to [k]$ 所对应的 $\tilde f_\alpha$ 之和。
这样的一族 $\tilde f_\alpha$ 定义单纯对象的态射，当且仅当
对任意 $\varphi : [m] \to [n]$，图
$$
\xymatrix{
\bigoplus_{\alpha : [n] \to [k]\text{ 满射}} A
\ar[r]_-{\tilde f_n}
\ar[d]_{(3)} &
V_n \ar[d]^{V(\varphi)} \\
\bigoplus_{\alpha' : [m] \to [k]\text{ 满射}} A
\ar[r]^-{\tilde f_m} &
V_m
}
$$
交换。取 $\varphi = \alpha$ 可知，对 $\tilde f_\alpha$ 的选择由
$f_k$ 唯一决定。一般情形的交换性可以分别在左上角的每个直和分量上
检验。对于满足 $\alpha \circ \varphi$ 为满射的 $\alpha$ 所对应的
分量，这是显然的，因为这些分量通过 $\text{id}_A$ 映到
$\alpha' = \alpha \circ \varphi$ 所对应的分量，并且
$\tilde f_{\alpha'} = V(\alpha') \circ f_k =
V(\alpha \circ \varphi) \circ f_k = V(\varphi) \circ \tilde f_\alpha$。
对于 $\alpha \circ \varphi$ 不是满射的情形，必须证明
$V(\varphi) \circ \tilde f_\alpha = 0$。按定义，它等于
$V(\varphi) \circ V(\alpha) \circ f_k =
V(\alpha \circ \varphi) \circ f_k$。
由于 $\alpha \circ \varphi$ 不是满射，可以将其写成
$\delta^k_i \circ \psi$，于是由上可得
$V(\varphi) \circ V(\alpha) \circ f_k =
V(\psi) \circ d^k_i \circ f_k = 0$。

\medskip\noindent
我们把 (4) 和 (5) 的规则作为一个单纯对象的定义，也记为
$\tilde U$。下面证明它是 $\text{cosk}_k U$ 的一种实现。
这将同时证明 (4) 和 (5)。论证与上面 (2)、(3) 的证明完全对偶，
但我们仍将其写出。必须证明：给定态射
$f : \text{sk}_kV \to U$，存在唯一态射
$\tilde f : V \to \tilde U$，使其取 $k$-骨架后恢复为 $f$。
由 (1)，$f$ 对应于态射 $f_k : V_k \to A$，且对所有 $i$
都在 $s^{k - 1}_i$ 的像上为零。对任意单射 $\beta : [k] \to [n]$，
令 $\tilde f_\beta : V_n \to A$ 等于复合
$\tilde f_\beta = f_k \circ V(\beta) : V_n \to A$。
定义 $\tilde f_n : V_n \to \tilde U_n$ 为所有单射
$\beta : [k] \to [n]$ 所对应的 $\tilde f_\beta$ 之和。
这样的一族 $\tilde f_\beta$ 定义单纯对象的态射，当且仅当
对任意 $\varphi : [m] \to [n]$，图
$$
\xymatrix{
V_n
\ar[d]_{V(\varphi)}
\ar[r]_-{\tilde f_n}
&
\bigoplus_{\beta : [k] \to [n]\text{ 单射}} A
\ar[d]^{(5)}
\\
V_m
\ar[r]^-{\tilde f_m}
&
\bigoplus_{\beta' : [k] \to [m]\text{ 单射}} A
}
$$
交换。取 $\varphi = \beta$ 可知，对 $\tilde f_\beta$ 的选择由
$f_k$ 唯一决定。一般情形的交换性可以分别在右下角的每个直和分量上
检验。对于满足 $\varphi \circ \beta'$ 为单射的 $\beta'$ 所对应的
分量，这是显然的，因为恰有 $\beta = \varphi \circ \beta'$ 所对应的
分量映入其中，并且在这种情形下
$\tilde f_{\beta'} \circ V(\varphi) =
f_k \circ V(\beta') \circ V(\varphi) =
f_k \circ V(\beta) = \tilde f_\beta$。
对于 $\varphi \circ \beta'$ 不是单射的情形，必须证明
$\tilde f_{\beta'} \circ V(\varphi) = 0$。按定义，它等于
$f_k \circ V(\beta') \circ V(\varphi) =
f_k \circ V(\varphi \circ \beta')$。
由于 $\varphi \circ \beta'$ 不是单射，可以将其写成
$\psi \circ \sigma^{k - 1}_i$，于是由上可得
$f_k \circ V(\beta') \circ V(\varphi) =
f_k \circ s^{k - 1}_i \circ V(\psi) = 0$。

\medskip\noindent
复合 $i_{k!}U \to \text{cosk}_kU$ 是唯一的单纯对象态射，
并且在 $A = U_k = (i_{k!}U)_k = (\text{cosk}_kU)_k$ 上为恒等映射。
因此，只需检验所给规则定义了单纯对象的态射。为此，必须证明
对任意 $\varphi : [m] \to [n]$，图
$$
\xymatrix{
\bigoplus_{\alpha : [n] \to [k]\text{ 满射}} A
\ar[d]_{(3)}
\ar[r]_{(6)}
&
\bigoplus_{\beta : [k] \to [n]\text{ 单射}} A
\ar[d]^{(5)}
\\
\bigoplus_{\alpha' : [m] \to [k]\text{ 满射}} A
\ar[r]^{(6)}
&
\bigoplus_{\beta' : [k] \to [m]\text{ 单射}} A
}
$$
交换。现在可按如下方式把它理解为仅以 $0$ 和 $1$ 填充的矩阵：
(3) 的矩阵具有非零的 $(\alpha', \alpha)$ 项，当且仅当
$\alpha' = \alpha \circ \varphi$。同样，(5) 的矩阵具有非零的
$(\beta', \beta)$ 项，当且仅当
$\beta = \varphi \circ \beta'$。(6) 的上方矩阵具有非零的
$(\alpha, \beta)$ 项，当且仅当
$\alpha \circ \beta = \text{id}_{[k]}$；(6) 的下方矩阵也一样。
于是，该图的交换性归结为分别计算两个复合的
$(\alpha, \beta')$ 项并看出它们相等。这又归结为等式
$$
\# \left\{
\beta
\mid
\beta = \varphi \circ \beta'
\text{ 且 }
\alpha \circ \beta = \text{id}_{[k]}
\right\}
=
\# \left\{
\alpha'
\mid
\alpha' = \alpha \circ \varphi
\text{ 且 }
\alpha' \circ \beta' = \text{id}_{[k]}
\right\}
$$
其证明可以放心地留给读者。

\medskip\noindent
最后证明 (7)。它直接由引理
\ref{simplicial-lemma-injective-map-simplicial-abelian}、
\ref{simplicial-lemma-recover-cosk}、\ref{simplicial-lemma-recovering-U-for-real}
和 \ref{simplicial-lemma-imshriek-abelian} 得出。
\end{proof}

\begin{definition}
\label{simplicial-definition-eilenberg-maclane}
设 $\mathcal{A}$ 为阿贝尔范畴，$A$ 为 $\mathcal{A}$ 的对象，
且 $k$ 为整数并满足 $\geq 0$。{\it Eilenberg--Mac Lane 对象 $K(A, k)$}
是对象 $K(A, k) = i_{k!}U$，它已在上面的引理
\ref{simplicial-lemma-eilenberg-maclane-object} 中描述。
\end{definition}


\begin{lemma}
\label{simplicial-lemma-extension}
设 $\mathcal{A}$ 为阿贝尔范畴，$A$ 为 $\mathcal{A}$ 的对象，
且 $k$ 为整数并满足 $\geq 0$。考虑由下列规则定义的单纯对象 $E$：
\begin{enumerate}
\item $E_n = \bigoplus_\alpha A$，其中直和遍历所有
$\alpha : [n] \to [k + 1]$，其像为 $[k]$ 或 $[k + 1]$；
\item 给定 $\varphi : [m] \to [n]$，若
$\Im(\alpha \circ \varphi)$ 等于 $[k]$ 或 $[k + 1]$，
则映射 $E_n \to E_m$ 通过 $\text{id}_A$ 把
$\alpha$ 所对应的直和分量映到
$\alpha \circ \varphi$ 所对应的直和分量。
\end{enumerate}
则存在逐项分裂正合的短正合序列
$$
0 \to K(A, k) \to E \to K(A, k + 1) \to 0
$$
\end{lemma}

\begin{proof}
映射 $K(A, k)_n \to E_n$ 以及 $E_n \to K(A, k + 1)_n$
分别由直和的嵌入与投影给出，这从指标集的包含关系是显然的。
显然，它们是单纯对象的态射。
\end{proof}

\begin{lemma}
\label{simplicial-lemma-abelian-limit-skeleta}
设 $\mathcal{A}$ 为阿贝尔范畴。对 $\mathcal{A}$ 的任意单纯对象
$V$，有
$$
V = \colim_n i_{n!}\text{sk}_n V
$$
且所有转移映射都是单射。
\end{lemma}

\begin{proof}
这是因为一旦 $n \geq m$，每个 $V_m$ 就等于
$(i_{n!}\text{sk}_n V)_m$。关于转移映射，另见引理
\ref{simplicial-lemma-n-skeleton-abelian}。
\end{proof}

\section{单纯对象与链复形}
\label{simplicial-section-complexes}

\noindent
设 $\mathcal{A}$ 为阿贝尔范畴。关于链复形的约定和记号，
见《同调代数》第 \ref{homology-section-complexes} 节。
设 $U$ 为 $\mathcal{A}$ 的单纯对象。$U$ 的
{\it 伴随链复形 $s(U)$}，有时称为 {\it Moore 复形}，是链复形
$$
\ldots \to U_2 \to U_1 \to U_0 \to 0 \to 0 \to \ldots
$$
其边界映射 $d_n : U_n \to U_{n - 1}$ 由公式
$$
d_n = \sum\nolimits_{i = 0}^n (-1)^i d^n_i.
$$
给出。它确为复形，因为根据备注 \ref{simplicial-remark-relations} 中列出的关系，
有
\begin{eqnarray*}
d_n \circ d_{n + 1}
& = &
(\sum\nolimits_{i = 0}^n (-1)^i d^n_i) \circ
(\sum\nolimits_{j = 0}^{n + 1} (-1)^j d^{n + 1}_j) \\
& = &
\sum\nolimits_{0 \leq i < j \leq n + 1}
(-1)^{i + j} d^n_{j - 1} \circ d^{n + 1}_i
+
\sum\nolimits_{n \geq i \geq j \geq 0}
(-1)^{i + j} d^n_i \circ d^{n + 1}_j \\
& = & 0.
\end{eqnarray*}
各符号相互抵消！我们用 $s(U)$ 表示这个伴随链复形。
显然，该构造具有函子性，因此定义函子
$$
s : \text{Simp}(\mathcal{A}) \longrightarrow \text{Ch}_{\geq 0}(\mathcal{A}).
$$
于是得到这个看似令人困惑但正确的公式 $s(U)_n = U_n$。

\begin{lemma}
\label{simplicial-lemma-s-exact}
函子 $s$ 是正合的。
\end{lemma}

\begin{proof}
由引理 \ref{simplicial-lemma-abelian} 立即得出。
\end{proof}

\begin{lemma}
\label{simplicial-lemma-homology-extension}
设 $\mathcal{A}$ 为阿贝尔范畴，$A$ 为 $\mathcal{A}$ 的对象，
$k$ 为整数。设 $E$ 为引理 \ref{simplicial-lemma-extension} 中描述的对象。
则复形 $s(E)$ 是零调的。
\end{lemma}

\begin{proof}
对态射 $\alpha : [n] \to [k + 1]$，定义
$\alpha' : [n + 1] \to [k + 1]$ 为满足
$\alpha'|_{[n]} = \alpha$ 及 $\alpha'(n + 1) = k + 1$ 的映射。
注意，如果 $\alpha$ 的像为 $[k]$ 或 $[k + 1]$，
那么 $\alpha'$ 的像为 $[k + 1]$。考虑一族映射
$h_n : E_n \to E_{n + 1}$，它通过 $A$ 上的恒等映射，
把 $\alpha$ 所对应的分量映到 $\alpha'$ 所对应的分量。
我们来计算 $d_{n + 1} \circ h_n - h_{n - 1} \circ d_n$。
先在 $\mathcal{A}$ 为阿贝尔群范畴的情形下计算。
用记号 $x_\alpha$ 表示 $E_n$ 中由指标集内的映射 $\alpha$
所对应的直和分量里的元素 $x \in A$。我们还约定：
只要 $\Im(\alpha)$ 不是 $[k]$ 或 $[k + 1]$，
$x_\alpha$ 就表示 $E_n$ 的零元。由这些约定可见
$$
d_{n + 1}(h_n(x_\alpha)) =
\sum\nolimits_{i = 0}^{n + 1} (-1)^i x_{\alpha' \circ \delta^{n + 1}_i}
$$
以及
$$
h_{n - 1}(d_n(x_\alpha)) =
\sum\nolimits_{i = 0}^n (-1)^i x_{(\alpha \circ \delta_i^n)'}
$$
容易看出，当 $i = 0, \ldots, n$ 时，
$\alpha' \circ \delta^{n + 1}_i = (\alpha \circ \delta_i^n)'$；
也容易看出 $\alpha' \circ \delta^{n + 1}_{n + 1} = \alpha$。
因此
$$
(d_{n + 1} \circ h_n - h_{n - 1} \circ d_n)(x_\alpha)
=
(-1)^{n + 1} x_\alpha
$$
若 $\mathcal{A}$ 为任意阿贝尔范畴，这些恒等式仍然成立，
因为它们在单纯阿贝尔群 $[n] \mapsto \Hom(A, E_n)$ 中成立；
细节留给读者。由此，$E$ 上的恒等映射同伦于零，
同伦由一族映射
$h'_n = (-1)^{n + 1}h_n : E_n \to E_{n + 1}$ 给出。
因而 $E$ 是零调的；例如可由《同调代数》引理
\ref{homology-lemma-map-homology-homotopy} 得出。
\end{proof}

\begin{lemma}
\label{simplicial-lemma-homology-eilenberg-maclane}
设 $\mathcal{A}$ 为阿贝尔范畴，$A$ 为 $\mathcal{A}$ 的对象，
$k$ 为整数。当 $i = k$ 时 $H_i(s(K(A, k))) = A$，
其他情形为 $0$。
\end{lemma}

\begin{proof}
先证明 $k = 0$ 的情形。此时，对所有 $n$ 都有
$K(A, 0)_n = A$。此外，这个单纯阿贝尔群中的所有映射都是
$\text{id}_A$；换言之，$K(A, 0)$ 是取值为 $A$ 的常值单纯对象。
边界映射
$d_n = \sum_{i = 0}^n (-1)^i \text{id}_A = 0$
（当 $n$ 为奇数），而当 $n$ 为偶数时，则 $= \text{id}_A$。
因此 $s(K(A, 0))$ 形如
$$
\ldots \to A \xrightarrow{0} A \xrightarrow{1} A \xrightarrow{0} A \to 0
$$
结论显然。

\medskip\noindent
下面对 $k$ 归纳证明一般结论。假设 $k$ 的结论成立，考虑引理
\ref{simplicial-lemma-extension} 中的短正合序列
$$
0 \to K(A, k) \to E \to K(A, k + 1) \to 0
$$
由引理 \ref{simplicial-lemma-abelian}，相应的链复形序列正合。
由引理 \ref{simplicial-lemma-homology-extension}，$s(E)$ 是零调的。
于是 $k + 1$ 的结论由同调长正合序列得出；见《同调代数》引理
\ref{homology-lemma-long-exact-sequence-chain}。
\end{proof}

\noindent
还可以把第二个链复形关联给 $\mathcal{A}$ 的单纯对象。
回顾引理 \ref{simplicial-lemma-splitting-abelian-category}：
$\mathcal{A}$ 的任意单纯对象 $U$ 都典范分裂，并且
$N(U_m) = \bigcap_{i = 0}^{m - 1} \Ker(d^m_i)$。
定义 {\it 规范化链复形 $N(U)$} 为链复形
$$
\ldots \to N(U_2) \to N(U_1) \to N(U_0) \to 0 \to 0 \to \ldots
$$
其边界映射 $d_n : N(U_n) \to N(U_{n - 1})$
是 $(-1)^nd^n_n$ 在 $U_n$ 的直和分量 $N(U_n)$ 上的限制。
注意，引理 \ref{simplicial-lemma-N-d-in-N} 蕴含
$d^n_n(N(U_n)) \subset N(U_{n - 1})$。
它确为复形，因为
$d^n_n \circ d^{n + 1}_{n + 1} = d^n_n \circ d^{n + 1}_n$，
而按定义，$d^{n + 1}_n$ 在 $N(U_{n + 1})$ 上为零。
这样得到第二个函子
$$
N : \text{Simp}(\mathcal{A}) \longrightarrow \text{Ch}_{\geq 0}(\mathcal{A}).
$$
以下说明微分中符号的由来。

\begin{lemma}
\label{simplicial-lemma-map-associated-complexes}
设 $\mathcal{A}$ 为阿贝尔范畴，$U$ 为 $\mathcal{A}$ 的单纯对象。
典范映射 $N(U_n) \to U_n$ 给出复形态射 $N(U) \to s(U)$。
\end{lemma}

\begin{proof}
这是显然的，因为 $s(U)_n = U_n$ 上的微分是
$\sum (-1)^i d^n_i$，且当 $i < n$ 时，映射 $d^n_i$
在 $N(U_n)$ 上为零；而按定义，$(-1)^nd^n_n$ 的限制正是
$N(U)$ 的边界映射。
\end{proof}

\begin{lemma}
\label{simplicial-lemma-N-K}
设 $\mathcal{A}$ 为阿贝尔范畴，$A$ 为 $\mathcal{A}$ 的对象，
$k$ 为整数。当 $i = k$ 时，$N(K(A, k))_i = A$；
其他情形为 $0$。
\end{lemma}

\begin{proof}
当 $i < k$ 时，$N(K(A, k))_i = 0$ 是显然的，因为此时
$K(A, k)_i = 0$。又因 $K(A, k)_{k - 1} = 0$ 且
$K(A, k)_k = A$，显然有 $N(K(A, k))_k = A$。
当 $i > k$ 时，由引理 \ref{simplicial-lemma-imshriek-abelian} 及
$K(A, k)$ 的定义可得 $N(K(A, k))_i = 0$；见定义
\ref{simplicial-definition-eilenberg-maclane}。
\end{proof}

\begin{lemma}
\label{simplicial-lemma-decompose-associated-complexes}
设 $\mathcal{A}$ 为阿贝尔范畴，$U$ 为 $\mathcal{A}$ 的单纯对象。
链复形的典范态射 $N(U) \to s(U)$ 是分裂的。事实上，对某个复形
$D(U)$，有
$$
s(U) = N(U) \oplus D(U)
$$
构造 $U \mapsto D(U)$ 具有函子性。
\end{lemma}

\begin{proof}
定义 $D(U)_n$ 为映射
$$
\bigoplus\nolimits_{\varphi : [n] \to [m]\text{ 满射}, \ m < n} N(U_m)
\xrightarrow{\bigoplus U(\varphi)}
U_n
$$
的像。根据引理 \ref{simplicial-lemma-splitting-abelian-category} 和定义
\ref{simplicial-definition-split}，它是 $U_n$ 中与 $N(U_n)$ 互补的子对象。
下面证明 $D(U)$ 是子复形。取满射
$\varphi : [n] \to [m]$，其中 $m < n$，并考虑复合
$$
N(U_m) \xrightarrow{U(\varphi)} U_n \xrightarrow{d_n} U_{n - 1}.
$$
这个复合是映射
$$
N(U_m) \xrightarrow{U(\varphi \circ \delta^n_i)} U_{n - 1}
$$
之和，其符号为 $(-1)^i$，其中 $i = 0, \ldots, n$。

\medskip\noindent
先对 $m$ 作升序归纳，证明当 $0 \leq m < n - 1$ 时，
所有映射 $U(\varphi \circ \delta^n_i)$ 都把 $N(U_m)$ 映入
$D(U)_{n - 1}$。（情形 $m = n - 1$ 在下面处理。）
只要映射 $\varphi \circ \delta^n_i : [n - 1] \to [m]$
为满射，根据定义，$N(U_m)$ 在
$U(\varphi \circ \delta^n_i)$ 下的像就包含于 $D(U)_{n - 1}$。
若 $\varphi \circ \delta^n_i : [n - 1] \to [m]$ 不是满射，
令 $j = \varphi(i)$，并注意 $i$ 是在 $\varphi$ 下映为 $j$
的唯一指标。对某个 $\psi : [n - 1] \to [m - 1]$，可写成
$\varphi \circ \delta^n_i = \delta^m_j \circ \psi \circ \delta^n_i$。
% P02-E0121: the frozen factorization and its induced U-map contain an extra delta_i^n and are not type-correct.
因此
$U(\varphi \circ \delta^n_i) = U(\psi \circ \delta^n_i) \circ d^m_j$，
它在 $N(U_m)$ 上为零，除非 $j = m$。若 $j = m$，则
$d^m_m(N(U_m)) \subset N(U_{m - 1})$，从而
$U(\varphi \circ \delta^n_i)(N(U_m)) \subset
U(\psi \circ \delta^n_i)(N(U_{m - 1}))$，由归纳假设即得结论。

\medskip\noindent
为了完成 $D(U)$ 是子复形的证明，还须在 $m = n - 1$ 的情形下
处理复合
$$
N(U_m) \xrightarrow{U(\varphi)} U_n \xrightarrow{d_n} U_{n - 1}.
$$
。
% P02-E0122: frozen prose places a full stop between the displayed composition and the following “in case” clause.
此时，对某个 $0 \leq j \leq n - 1$，有
$\varphi = \sigma^{n - 1}_j$，且 $U(\varphi) = s^{n - 1}_j$。
因此，该复合由下列和给出：
$$
\sum (-1)^i d^n_i \circ s^{n - 1}_j
$$
回顾备注 \ref{simplicial-remark-relations}，
$d^n_j \circ s^{n - 1}_j = d^n_{j + 1} \circ s^{n - 1}_j = \text{id}$，
而相应两项符号相反，故相互抵消。若 $j < n - 1$，则映射
$d^n_n \circ s^{n - 1}_j$ 等于
$s^{n - 2}_j \circ d^{n - 1}_{n - 1}$。由于
$d^{n - 1}_{n - 1}$ 把 $N(U_{n - 1})$ 映入 $N(U_{n - 2})$，
可见像 $d^n_n ( s^{n - 1}_j (N(U_{n - 1}))$
% P02-E0123: frozen source leaves the outer parenthesis in this image expression unclosed.
包含于 $s^{n - 2}_j(N(U_{n - 2}))$，而根据定义，后者包含于
$D(U_{n - 1})$。
% P02-E0124: frozen source applies D to the component U_{n-1}; the established subcomplex term is D(U)_{n-1}.
对于所有其他的 $(i, j)$ 组合，或者
$d^n_i \circ s^{n - 1}_j = s^{n - 2}_{j - 1} \circ d^{n - 1}_i$
（若 $i < j$），或者
$d^n_i \circ s^{n - 1}_j = s^{n - 2}_j \circ d^{n - 1}_{i - 1}$
（若 $n > i > j + 1$）；在这些情形下，由 $N(U_{n - 1})$
的定义，该映射为零。
\end{proof}

\begin{remark}
\label{simplicial-remark-degenerate-subcomplex}
在引理 \ref{simplicial-lemma-decompose-associated-complexes} 的情形下，
子复形 $D(U) \subset s(U)$ 也可以定义为项
$$
D(U)_n = \Im\left(
\bigoplus\nolimits_{\varphi : [n] \to [m]\text{ 满射}, \ m < n} U_m
\xrightarrow{\bigoplus U(\varphi)}
U_n\right)
$$
所组成的子复形。事实上，$U_m$ 是子对象 $N(U_m)$
与所有满射 $[m] \to [k]$（其中 $k < m$）所给出的
$N(U_k)$ 之像的直和；因此这显然等同于引理
\ref{simplicial-lemma-decompose-associated-complexes} 的证明中给出的
$D(U)_n$ 的定义。由此可见，若 $U$ 是单纯阿贝尔群，
则 $D(U)_n$ 的元素恰为退化 $n$-单形之和。
\end{remark}

\begin{lemma}
\label{simplicial-lemma-N-exact}
函子 $N$ 是正合的。
\end{lemma}

\begin{proof}
由引理 \ref{simplicial-lemma-s-exact} 及引理
\ref{simplicial-lemma-decompose-associated-complexes} 的函子性分解即得。
\end{proof}

\begin{lemma}
\label{simplicial-lemma-quasi-isomorphism}
设 $\mathcal{A}$ 为阿贝尔范畴，$V$ 为 $\mathcal{A}$ 的单纯对象。
链复形的典范态射 $N(V) \to s(V)$ 是拟同构。换言之，引理
\ref{simplicial-lemma-decompose-associated-complexes} 中的复形 $D(V)$ 是零调的。
\end{lemma}

\begin{proof}
由引理 \ref{simplicial-lemma-homology-eilenberg-maclane} 和
\ref{simplicial-lemma-N-K}，对任意对象 $A$ 和任意 $k \geq 0$，
结论对 $K(A, k)$ 成立。考虑假设 $IH_{n, m}$：
对所有满足 $j \leq m$ 时 $V_j = 0$ 的 $V$，并对所有
$i \leq n$，映射 $N(V) \to s(V)$ 诱导同构
$H_i(N(V)) \to H_i(s(V))$。

\medskip\noindent
作为归纳的起点，$IH_{n, n}$ 显然成立，因为在这种情形下
$N(V)_n = 0$ 且 $s(V)_n = 0$。
% P02-E0125: frozen English says “To start of the induction”.

\medskip\noindent
假设 $IH_{n, m}$，其中 $m \leq n$。取满足当 $j < m$
时 $V_j = 0$ 的单纯对象 $V$。由引理
\ref{simplicial-lemma-eilenberg-maclane-object} 和定义
\ref{simplicial-definition-eilenberg-maclane}，有
$K(V_m, m) = i_{m!} \text{sk}_mV$。由引理
\ref{simplicial-lemma-n-skeleton-abelian}，自然态射
$$
K(V_m, m) = i_{m!} \text{sk}_mV \to V
$$
为单射。因此，对某个满足当 $i = 0, \ldots, m$ 时
$W_i = 0$ 的 $W$，得到短正合序列
$$
0 \to K(V_m, m) \to V \to W \to 0
$$
这个短正合序列诱导伴随复形的短正合序列之间的态射
% P02-E0126: frozen English has singular “a morphism of short exact sequence”.
$$
\xymatrix{
0 \ar[r] &
N(K(V_m, m)) \ar[r] \ar[d] &
N(V) \ar[r] \ar[d] &
N(W) \ar[r] \ar[d] &
0 \\
0 \ar[r] &
s(K(V_m, m)) \ar[r] &
s(V) \ar[r] &
s(W) \ar[r] &
0
}
$$
；见引理 \ref{simplicial-lemma-s-exact} 和 \ref{simplicial-lemma-N-exact}。
因此，由两端的结论可推出 $V$ 的结论。
\end{proof}






\section{Dold--Kan}
\label{simplicial-section-dold-kan}

\noindent
本节证明 Dold--Kan 定理，它把阿贝尔范畴中的单纯对象与链复形联系起来。

\begin{lemma}
\label{simplicial-lemma-N-faithful}
设 $\mathcal{A}$ 为阿贝尔范畴。函子 $N$ 是忠实的，并且反映同构、
单射和满射。
\end{lemma}

\begin{proof}
忠实性立即由引理 \ref{simplicial-lemma-splitting-abelian-category} 的典范分裂
得出。反映单射、满射和同构的断言由引理
\ref{simplicial-lemma-injective-map-simplicial-abelian} 得出。
\end{proof}

\begin{lemma}
\label{simplicial-lemma-S-N}
设 $\mathcal{A}$ 和 $\mathcal{B}$ 为阿贝尔范畴，并设
$N : \mathcal{A} \to \mathcal{B}$ 与
$S : \mathcal{B} \to \mathcal{A}$ 为函子。假设：
\begin{enumerate}
\item 函子 $S$ 和 $N$ 都正合；
\item 存在到 $\mathcal{B}$ 上恒等函子的同构
$g : N \circ S \to \text{id}_\mathcal{B}$；
\item $N$ 忠实；
\item $S$ 本质满。
\end{enumerate}
则 $S$ 和 $N$ 是互为拟逆的范畴等价。
\end{lemma}

\begin{proof}
只需构造函子性同构 $S(N(A)) \cong A$。为此，选取 $B$
和同构 $f : A  \to S(B)$。考虑映射
$$
f^{-1} \circ g_{S(B)} \circ S(N(f)) :
S(N(A)) \to S(N(S(B))) \to S(B) \to A.
$$
容易证明它不依赖于对 $f, B$ 的选择，并给出所需同构
$S \circ N \to \text{id}_\mathcal{A}$。
\end{proof}

\begin{theorem}
\label{simplicial-theorem-dold-kan}
设 $\mathcal{A}$ 为阿贝尔范畴。函子 $N$ 诱导范畴等价
$$
N :
\text{Simp}(\mathcal{A})
\longrightarrow
\text{Ch}_{\geq 0}(\mathcal{A})
$$
\end{theorem}

\begin{proof}
我们将描述一个反方向的函子，它受引理 \ref{simplicial-lemma-extension}
中的构造启发（但为了得到正确的边界映射，要加入一个符号）。
设 $A_\bullet$ 为边界映射
$d_{A, n} : A_n \to A_{n - 1}$ 的链复形。对每个 $n \geq 0$，记
$$
I_n =
\Big\{
\alpha : [n] \to \{0, 1, 2, \ldots\}
\mid
\Im(\alpha) = [k]\text{ 对某个 }k
\Big\}.
$$
对 $\alpha \in I_n$，以 $k(\alpha)$ 表示满足
$\Im(\alpha) = [k]$ 的唯一整数。如下定义单纯对象
$S(A_\bullet)$：
\begin{enumerate}
\item $S(A_\bullet)_n = \bigoplus_{\alpha \in I_n} A_{k(\alpha)}$，
并将其写成
$\bigoplus_{\alpha \in I_n} A_{k(\alpha)} \cdot \alpha$，
以提示把“$\alpha$”看作相应直和分量的基向量；
\item 给定 $\varphi : [m] \to [n]$，按如下规则，在
$S(A_\bullet)_n$ 的直和分量
$A_{k(\alpha)} \cdot \alpha$ 上定义 $S(A_\bullet)(\varphi)$：
\begin{enumerate}
\item 若 $\alpha \circ \varphi \not \in I_m$，则令它等于零；
% P02-E0127: frozen English omits “If” at the start of this case.
\item 若 $\alpha \circ \varphi \in I_m$，但
$k(\alpha \circ \varphi)$ 既不等于 $k(\alpha)$，也不等于
$k(\alpha) - 1$，则也令它等于零；
\item 若 $\alpha \circ \varphi \in I_m$ 且
$k(\alpha \circ \varphi) = k(\alpha)$，则通过恒等映射映到
$S(A_\bullet)_m$ 的分量
$A_{k(\alpha \circ \varphi)} \cdot (\alpha \circ \varphi)$；
\item 若 $\alpha \circ \varphi \in I_m$ 且
$k(\alpha \circ \varphi) = k(\alpha) - 1$，则使用
$(-1)^{k(\alpha)} d_{A, k(\alpha)}$ 映到
$S(A_\bullet)_m$ 的分量
$A_{k(\alpha \circ \varphi)}\cdot (\alpha \circ \varphi)$。
\end{enumerate}
\end{enumerate}
下面证明 $S(A_\bullet)$ 是 $\mathcal{A}$ 的单纯对象。
为此，设有映射 $\varphi : [m] \to [n]$ 和
$\psi : [n] \to [p]$。我们将证明
$S(A_\bullet)(\varphi) \circ S(A_\bullet)(\psi) =
S(A_\bullet)(\psi \circ \varphi)$。选取 $\beta \in I_p$，并令
$\alpha = \beta \circ \psi$ 和 $\gamma = \alpha \circ \varphi$，
将它们分别视为映射
$\alpha : [n] \to \{0, 1, 2, \ldots\}$ 和
$\gamma : [m] \to \{0, 1, 2, \ldots\}$。图示为
$$
\xymatrix{
[m] \ar[r]_\varphi \ar[d]_\gamma &
[n] \ar[r]_\psi \ar[d]_\alpha &
[p] \ar[d]^\beta \\
\Im(\gamma) \ar[r] &
\Im(\alpha) \ar[r] &
[k(\beta)]
}
$$
我们将证明映射
$S(A_\bullet)(\varphi) \circ S(A_\bullet)(\psi)$ 和
$S(A_\bullet)(\psi \circ \varphi)$
在分量 $A_{k(\beta)} \cdot \beta$ 上的限制相同。
% P02-E0128: frozen source breaks this sentence with a full stop before “to the summand”.
需要考虑若干情形：
\begin{enumerate}
\item 设 $\alpha \not \in I_n$，因而 $S(A_\bullet)(\psi)$
在 $A_{k(\beta)} \cdot \beta$ 上的限制为零。
此时，或者 $\gamma \not \in I_m$；或者
$[k(\gamma)] = \Im(\gamma) \subset \Im(\alpha) \subset [k(\beta)]$，
且 $[k(\beta)]$ 的子集 $\Im(\alpha)$ 有缺口，故
$k(\gamma) < k(\beta) - 1$。在两种情形下，
$S(A_\bullet)(\psi \circ \varphi)$ 在
$A_{k(\beta)} \cdot \beta$ 上的限制也为零。
\item 设 $\alpha \in I_n$ 且 $k(\alpha) < k(\beta) - 1$，
因而 $S(A_\bullet)(\psi)$ 在 $A_{k(\beta)} \cdot \beta$
上的限制为零。此时，或者 $\gamma \not \in I_m$；或者
$[k(\gamma)] \subset [k(\alpha)] \subset [k(\beta)]$，
从而 $k(\gamma) < k(\beta) - 1$。在两种情形下，
$S(A_\bullet)(\psi \circ \varphi)$ 在
$A_{k(\beta)} \cdot \beta$ 上的限制也为零。
\item 设 $\alpha \in I_n$ 且 $k(\alpha) = k(\beta)$，
因而 $S(A_\bullet)(\psi)$ 在 $A_{k(\beta)} \cdot \beta$
上的限制是从 $A_{k(\beta)} \cdot \beta$ 到
$A_{k(\alpha)} \cdot \alpha$ 的恒等映射。在这种情形下，
由于 $\Im(\alpha) =[k(\beta)]$，
$S(A_\bullet)(\psi \circ \varphi)$ 在分量
$A_{k(\beta)} \cdot \beta$ 上的限制规则，恰与
$S(A_\bullet)(\varphi)$ 在分量
$A_{k(\alpha)} \cdot \alpha$ 上的限制规则相同，故二者相合。
\item 设 $\alpha \in I_n$ 且 $k(\alpha) = k(\beta) - 1$，
因而 $S(A_\bullet)(\psi)$ 在 $A_{k(\beta)} \cdot \beta$
上的限制由到 $A_{k(\alpha)} \cdot \alpha$ 的
$(-1)^{k(\beta)}d_{A, k(\beta)}$ 给出。分以下子情形：
\begin{enumerate}
\item 若 $\gamma \not \in I_m$，则
$S(A_\bullet)(\psi \circ \varphi)$ 在分量
$A_{k(\beta)} \cdot \beta$ 上的限制，以及
$S(A_\bullet)(\varphi)$ 在分量
$A_{k(\alpha)} \cdot \alpha$ 上的限制都为零，故二者相合。
\item 若 $\gamma \in I_m$ 但 $k(\gamma) < k(\alpha) - 1$，
则两个限制同样都为零，故二者相合。
\item 若 $\gamma \in I_m$ 且 $k(\gamma) = k(\alpha)$，
则 $\Im(\gamma) = \Im(\alpha)$。此时，
$S(A_\bullet)(\psi \circ \varphi)$ 在分量
$A_{k(\beta)} \cdot \beta$ 上的限制由到
$A_{k(\gamma)} \cdot \gamma$ 的
$(-1)^{k(\beta)}d_{A, k(\beta)}$ 给出，而
$S(A_\bullet)(\varphi)$ 在分量
$A_{k(\alpha)} \cdot \alpha$ 上的限制是恒等映射
$A_{k(\alpha)} \cdot \alpha \to A_{k(\gamma)} \cdot \gamma$。
故二者相合。
\item 最后，若 $\gamma \in I_m$ 且
$k(\gamma) = k(\alpha) - 1$，则
$S(A_\bullet)(\varphi)$ 在分量
$A_{k(\alpha)} \cdot \alpha$ 上的限制由
$(-1)^{k(\alpha)} d_{A, k(\alpha)}$ 给出，作为映射
$A_{k(\alpha)} \cdot \alpha \to A_{k(\beta)} \cdot \beta$。
% P02-E0129: the frozen codomain here is the beta summand, but S(phi) should land in the gamma summand.
由于 $A_\bullet$ 是复形，可见复合
$A_{k(\beta)} \cdot \beta \to
A_{k(\alpha)} \cdot \alpha \to A_{k(\gamma)} \cdot \gamma$
为零；这与 $S(A_\bullet)(\psi \circ \varphi)$ 在分量
$A_{k(\beta)} \cdot \beta$ 上的限制所得结果相合，因为
$k(\gamma) = k(\beta) - 2 < k(\beta) - 1$。
\end{enumerate}
\end{enumerate}
所以 $S(A_\bullet)$ 是 $\mathcal{A}$ 的单纯对象。

\medskip\noindent
下面构造关于 $A_\bullet$ 具有函子性的同构
$A_\bullet \to N(S(A_\bullet))$。回顾引理
\ref{simplicial-lemma-decompose-associated-complexes} 给出的链复形分解
$$
S(A_\bullet) =  N(S(A_\bullet)) \oplus D(S(A_\bullet))
$$
另一方面，由备注 \ref{simplicial-remark-degenerate-subcomplex} 和
$S(A_\bullet)$ 的构造可得
$$
D(S(A_\bullet))_n =
\bigoplus\nolimits_{\alpha \in I_n,\ k(\alpha) < n} A_{k(\alpha)} \cdot \alpha
\subset
\bigoplus\nolimits_{\alpha \in I_n} A_{k(\alpha)} \cdot \alpha
$$
然而，若 $\alpha \in I_n$，则
$k(\alpha) \geq n \Leftrightarrow \alpha = \text{id}_{[n]} : [n] \to [n]$。
因此，$S(A_\bullet)_n$ 的分量
$A_n \cdot \text{id}_{[n]}$ 是分量 $D(S(A_\bullet))_n$ 的补。
所有映射 $d^n_i : S(A_\bullet)_n \to S(A_\bullet)_n$
% P02-E0130: frozen source gives the face maps codomain degree n instead of n-1.
在分量 $A_n \cdot \text{id}_{[n]}$ 上的限制都为零，
唯有 $d^n_n$ 从 $A_n \cdot \text{id}_{[n]}$ 到
$A_{n - 1} \cdot \text{id}_{[n - 1]}$ 给出
$(-1)^n d_{A, n}$。由此断定，$A_n \cdot \text{id}_{[n]}$
必定等于分量 $N(S(A_\bullet))_n$；而且微分
$d_n = \sum (-1)^id^n_i : S(A_\bullet)_n \to S(A_\bullet)_{n - 1}$
在分量 $A_n \cdot \text{id}_{[n]}$ 上的限制正给出所需映射！

\medskip\noindent
最后，必须证明 $S \circ N$ 同构于恒等函子。设 $U$ 为
$\mathcal{A}$ 的单纯对象。则可定义显然的映射
$$
S(N(U))_n = \bigoplus\nolimits_{\alpha \in I_n} N(U)_{k(\alpha)} \cdot \alpha
\longrightarrow
U_n
$$
：在 $\alpha$ 所对应的分量上使用
$U(\alpha) : N(U)_{k(\alpha)} \to U_n$。由定义
\ref{simplicial-definition-split}，这是同构。为了完成证明，还须证明它与单纯对象
中的映射相容。为此，设 $\varphi : [m] \to [n]$ 且
$\alpha \in I_n$，并令 $\beta = \alpha \circ \varphi$。图示为
$$
\xymatrix{
[m] \ar[r]_\varphi \ar[d]_\beta &
[n] \ar[d]_\alpha \\
\Im(\beta) \ar[r] &
[k(\alpha)]
}
$$
需要考虑若干情形：
\begin{enumerate}
\item 设 $\beta \not \in I_m$。则存在指标
$0 \leq j < k(\alpha)$，满足
$j \not \in \Im(\alpha \circ \varphi)$，从而可以选取分解
$\alpha \circ \varphi = \delta^{k(\alpha)}_j \circ \psi$，
其中 $\psi : [m] \to [k(\alpha) - 1]$。由此可知，
$U(\varphi)$ 在分量 $N(U)_{k(\alpha)} \cdot \alpha$ 的像上为零，
因为 $U(\varphi) \circ U(\alpha) =
U(\alpha \circ \varphi) = U(\psi) \circ d^{k(\alpha)}_j$
根据 $N$ 的构造在 $N(U)_{k(\alpha)}$ 上为零。这与上面对
$S(N(U))$ 所定的规则一致。
\item 设 $\beta \in I_m$ 且 $k(\beta) < k(\alpha) - 1$。
此时完全按情形 (1) 论证，取 $j = k(\alpha) - 1$。
\item 设 $\beta \in I_m$ 且 $k(\beta) = k(\alpha)$。
此时，分量 $N(U)_{k(\alpha)} \cdot \alpha$ 通过恒等映射
映到分量 $N(U)_{k(\beta)} \cdot \beta$。这与
$U(\varphi)$ 的作用相同，因为在这种情形下
$U(\varphi) \circ U(\alpha) = U(\beta)$。
\item 设 $\beta \in I_m$ 且 $k(\beta) = k(\alpha) - 1$。
此时，使用微分 $(-1)^{k(\alpha)} d_{N(U), k(\alpha)}$
把分量 $N(U)_{k(\alpha)} \cdot \alpha$ 映到分量
$N(U)_{k(\beta)} \cdot \beta$。另一方面，因为此时
$\Im(\beta) = [k(\beta)]$，有
$\alpha \circ \varphi = \delta^{k(\alpha)}_{k(\alpha)} \circ \beta$。
因此，$U(\varphi)$ 与 $U(\alpha)$ 在 $N(U)_{k(\alpha)}$
上的限制之复合，等于 $U(\beta)$ 前复合
$d^{k(\alpha)}_{k(\alpha)}$ 在 $N(U)_{k(\alpha)}$ 上的限制。
由于 $d_{N(U), k(\alpha)} = \sum (-1)^i d^{k(\alpha)}_i$，
并且当 $i < k(\alpha)$ 时，$d^{k(\alpha)}_i$ 在
$N(U)_{k(\alpha)}$ 上的限制为零，可见二者相等。
\end{enumerate}
定理得证。
\end{proof}







\section{余单纯对象的 Dold--Kan 对应}
\label{simplicial-section-dold-kan-cosimplicial}

\noindent
设 $\mathcal{A}$ 为阿贝尔范畴。根据《同调代数》引理
\ref{homology-lemma-abelian-opposite}，$\mathcal{A}^{opp}$
也是阿贝尔范畴。由定义形式地推出
$$
\text{CoSimp}(\mathcal{A}) = \text{Simp}(\mathcal{A}^{opp})^{opp}.
$$
于是 Dold--Kan 对应（定理 \ref{simplicial-theorem-dold-kan}）蕴含
$\text{CoSimp}(\mathcal{A})$ 等价于范畴
$\text{Ch}_{\geq 0}(\mathcal{A}^{opp})^{opp}$。
由定义还形式地推出
$$
\text{CoCh}_{\geq 0}(\mathcal{A}) =
\text{Ch}_{\geq 0}(\mathcal{A}^{opp})^{opp}.
$$
将这些箭头合在一起，得到等价
$$
Q :
\text{CoSimp}(\mathcal{A})
\longrightarrow
\text{CoCh}_{\geq 0}(\mathcal{A}).
$$
本节描述 $Q$。

\medskip\noindent
先定义余单纯对象 $U$ 的
{\it 伴随上链复形 $s(U)$}。它是在负次数中为零，并且当
$n \geq 0$ 时满足 $s(U)^n = U_n$ 的上链复形。
作为微分，使用映射
$d^n : s(U)^n \to s(U)^{n + 1}$，其定义为
$d^n = \sum_{i = 0}^{n + 1} (-1)^i \delta^{n + 1}_i$。
换言之，复形 $s(U)$ 形如
$$
\xymatrix{
0 \ar[r] &
U_0 \ar[rr]^{\delta^1_0 - \delta^1_1} & &
U_1 \ar[rr]^{\delta^2_0 - \delta^2_1 + \delta^2_2} & &
U_2 \ar[r] &
\ldots
}
$$
它有时也称为 $U$ 的 {\it Moore 复形}。

\medskip\noindent
另一方面，给定 $\mathcal{A}$ 的余单纯对象 $U$，令
$Q(U)^0 = U_0$，并令
$$
Q(U)^n = \Coker(
\xymatrix{
\bigoplus_{i = 0}^{n - 1} U_{n - 1} \ar[r]^-{\delta^n_i} &
U_n
}).
$$
微分 $d^n : Q(U)^n \to Q(U)^{n + 1}$ 由
$(-1)^{n + 1}\delta^{n + 1}_{n + 1}$ 诱导；也就是说，把态射
$(-1)^{n + 1}\delta^{n + 1}_{n + 1}$ 放入交换图
$$
\xymatrix{
U_n \ar[rr]_{(-1)^{n + 1}\delta^{n + 1}_{n + 1}} \ar[d] & &
U_{n + 1} \ar[d] \\
Q(U)^n \ar[rr]^{d_n} & &
Q(U)^{n + 1}.
}
$$
% P02-E0131: the frozen diagram labels the induced cochain differential d_n although the surrounding definition calls it d^n.
留给读者证明这个图是有意义的；即当
$i = 0, \ldots, n - 1$ 时，$\delta^n_i$ 的像映入右侧竖直箭头的核。
（这与引理 \ref{simplicial-lemma-N-d-in-N} 对偶。）
因此，上链复形 $Q(U)$ 形如
$$
0 \to Q(U)^0 \to Q(U)^1 \to Q(U)^2 \to \ldots
$$
它称为 $U$ 的 {\it 伴随规范化上链复形}。
Dold--Kan 定理 \ref{simplicial-theorem-dold-kan} 的对偶如下。

\begin{lemma}
\label{simplicial-lemma-dual-dold-kan}
设 $\mathcal{A}$ 为阿贝尔范畴。
\begin{enumerate}
\item 函子
$s : \text{CoSimp}(\mathcal{A}) \to \text{CoCh}_{\geq 0}(\mathcal{A})$
是正合的。
\item 映射 $s(U)^n \to Q(U)^n$ 定义上链复形的态射。
\item 在 $\text{CoCh}_{\geq 0}(\mathcal{A})$ 中存在函子性直和分解
$s(U) = D(U) \oplus Q(U)$。
\item 函子 $Q$ 是正合的。
\item 复形态射 $s(U) \to Q(U)$ 是拟同构。
\item 函子 $U \mapsto Q(U)^\bullet$ 定义范畴等价
$\text{CoSimp}(\mathcal{A}) \to \text{CoCh}_{\geq 0}(\mathcal{A})$。
\end{enumerate}
\end{lemma}

\begin{proof}
略。但这些结果正好分别是引理
\ref{simplicial-lemma-s-exact}、\ref{simplicial-lemma-map-associated-complexes}、
\ref{simplicial-lemma-decompose-associated-complexes}、
\ref{simplicial-lemma-N-exact}、\ref{simplicial-lemma-quasi-isomorphism}
以及定理 \ref{simplicial-theorem-dold-kan} 的对偶陈述。
\end{proof}

\section{同伦}
\label{simplicial-section-homotopy}

\noindent
考虑单纯集 $\Delta[0]$ 和 $\Delta[1]$。回顾有两个态射
$$
e_0, e_1 : \Delta[0] \longrightarrow \Delta[1],
$$
它们来自把 $0$ 映到 $[1] = \{0, 1\}$ 的一个元素的态射
$[0] \to [1]$。还回顾每个集合 $\Delta[1]_k$ 都是有限的。
因此，若范畴 $\mathcal{C}$ 有有限余积，则对 $\mathcal{C}$
的任意单纯对象 $U$，都可以构成积
$$
U \times \Delta[1]
$$
；见定义 \ref{simplicial-definition-product-with-simplicial-set}。
注意，$\Delta[0]$ 的性质是：对所有 $k \geq 0$，
$\Delta[0]_k = \{*\}$ 都是单点集。因此
$U \times \Delta[0] = U$。于是，上面的 $e_0, e_1$ 给出态射
% P02-E0132: frozen English uses singular “gives rise” with the compound subject e_0, e_1.
$$
e_0, e_1 : U \to U \times \Delta[1].
$$

\begin{definition}
\label{simplicial-definition-homotopy}
设 $\mathcal{C}$ 为有有限余积的范畴，$U$ 和 $V$ 为
$\mathcal{C}$ 的两个单纯对象，并设 $a, b : U \to V$ 为两个态射。
\begin{enumerate}
\item 如果态射
$$
h : U \times \Delta[1] \longrightarrow V
$$
满足 $a = h \circ e_0$ 且 $b = h \circ e_1$，则称它为
{\it 从 $a$ 到 $b$ 的同伦}。
\item 如果存在一列从 $U$ 到 $V$ 的态射
$a = a_0, a_1, \ldots, a_n = b$，使得对每个
$i = 1, \ldots, n$，或者存在从 $a_{i - 1}$ 到 $a_i$ 的同伦，
或者存在从 $a_i$ 到 $a_{i - 1}$ 的同伦，则称态射 $a$ 和 $b$
{\it 同伦}，或者称它们 {\it 属于同一个同伦类}。
\end{enumerate}
\end{definition}

\noindent
一般而言，“存在从 $a$ 到 $b$ 的同伦”这一关系既不传递也不对称；
例 \ref{simplicial-example-trivial-homotopy} 将说明它是自反的。当然，“同伦”
是集合 $\Mor(U, V)$ 上的等价关系，并且是由“存在从 $a$ 到 $b$
的同伦”这一关系生成的等价关系。事实上，可以在任意范畴中定义
单纯对象态射对之间的同伦。在详细分析存在态射
$h : U \times \Delta[1] \to V$ 的含义之后，我们将在备注
\ref{simplicial-remark-homotopy-better} 中完成这一点。

\medskip\noindent
设 $\mathcal{C}$ 为有有限余积的范畴，$U$、$V$ 为
$\mathcal{C}$ 的单纯对象，$a, b : U \to V$ 为态射。
进一步假设 $h : U \times \Delta[1] \to V$ 是从 $a$ 到 $b$
的同伦。对每个 $n \geq 0$，记
$$
\Delta[1]_n = \{\alpha^n_0, \ldots, \alpha^n_{n + 1}\}
$$
其中 $\alpha^n_i : [n] \to [1]$ 是满足
$$
\alpha^n_i(j)
=
\left\{
\begin{matrix}
0 & \text{若} & j < i\\
1 & \text{若} & j \geq i
\end{matrix}
\right.
$$
的映射。因此
$$
h_n : (U \times \Delta[1])_n = \coprod U_n \cdot \alpha^n_i
\longrightarrow
V_n
$$
具有分量 $h_{n, i} : U_n \to V_n$；对每个
$i = 0, \ldots, n + 1$，它是在 $\alpha^n_i$ 所对应的分量上的限制。

\begin{lemma}
\label{simplicial-lemma-relations-homotopy}
在上述情形下，有如下关系：
\begin{enumerate}
\item $h_{n, 0} = b_n$ 且 $h_{n, n + 1} = a_n$。
\item 当 $i > j$ 时，
$d^n_j \circ h_{n, i} = h_{n - 1, i - 1} \circ d^n_j$。
\item 当 $i \leq j$ 时，
$d^n_j \circ h_{n, i} = h_{n - 1, i} \circ d^n_j$。
\item 当 $i > j$ 时，
$s^n_j \circ h_{n, i} = h_{n + 1, i + 1} \circ s^n_j$。
\item 当 $i \leq j$ 时，
$s^n_j \circ h_{n, i} = h_{n + 1, i} \circ s^n_j$。
\end{enumerate}
反之，给定满足上述性质的一族映射 $h_{n, i}$，它们定义态射
$h$，并且这是从 $a$ 到 $b$ 的同伦。
\end{lemma}

\begin{proof}
略。最后一个断言可利用如下事实证明：根据引理
\ref{simplicial-lemma-face-degeneracy-category}，给出单纯对象的态射，
等价于给出与所有 $d^n_j$ 和 $s^n_j$ 交换的一列态射 $h_n$。
\end{proof}

\begin{example}
\label{simplicial-example-trivial-homotopy}
在上述情形下，假设 $a = b$。则存在从 $a$ 到 $b$ 的
{\it 平凡}同伦，即满足 $h_{n, i} = a_n = b_n$ 的同伦。
\end{example}

\begin{remark}
\label{simplicial-remark-homotopy-better}
设 $\mathcal{C}$ 为任意范畴（不作任何假设），$U$ 和 $V$
为 $\mathcal{C}$ 的单纯对象，$a, b : U \to V$ 为
$\mathcal{C}$ 中单纯对象的态射。一个 {\it 从 $a$ 到 $b$ 的同伦}
由态射\footnote{文献中经常使用映射
$h_{n + 1, i} \circ s_i : U_n \to V_{n + 1}$，而不是映射
$h_{n, i}$。当然，这些映射所满足的关系不同于引理
\ref{simplicial-lemma-relations-homotopy} 中的关系。}
$h_{n, i} : U_n \to V_n$ 给出，其中 $n \geq 0$ 且
$i = 0, \ldots, n + 1$，并要求它们满足引理
\ref{simplicial-lemma-relations-homotopy} 中的关系。与定义
\ref{simplicial-definition-homotopy} 一样，如果存在一列从 $U$ 到 $V$ 的态射
$a = a_0, a_1, \ldots, a_n = b$，使得对每个
$i = 1, \ldots, n$，或者存在从 $a_{i - 1}$ 到 $a_i$ 的同伦，
或者存在从 $a_i$ 到 $a_{i - 1}$ 的同伦，则称态射 $a$ 和 $b$
{\it 同伦}。显然，若 $F : \mathcal{C} \to \mathcal{C}'$
为任意函子，且 $\{h_{n, i}\}$ 是从 $a$ 到 $b$ 的同伦，
则 $\{F(h_{n, i})\}$ 是从 $F(a)$ 到 $F(b)$ 的同伦。
同样，若 $a$ 和 $b$ 同伦，则 $F(a)$ 和 $F(b)$ 同伦。
由于该引理表明，当有限余积存在时，新概念与旧概念相同，
% P02-E0133: frozen English says singular “finite coproduct exist”.
特别可知，只要两个范畴都有有限余积，函子就保持原来的同伦概念。
\end{remark}

\begin{remark}
\label{simplicial-remark-homotopy-pre-post-compose}
设 $\mathcal{C}$ 为任意范畴。假设单纯对象的两个态射
$a, a' : U \to V$ 同伦。则对任意态射 $b : V \to W$，
两个映射 $b \circ a, b \circ a' : U \to W$ 同伦。
同样，对任意态射 $c : X \to U$，两个映射
$a \circ c, a' \circ c : X \to V$ 同伦。事实上，映射
$b \circ a \circ c, b \circ a' \circ c : X \to W$ 同伦。
确切地说，如果映射 $h_{n, i} : U_n \to V_n$ 定义从
$a$ 到 $a'$ 的同伦，那么映射 $b \circ h_{n, i} \circ c$
定义从 $b \circ a \circ c$ 到 $b \circ a' \circ c$ 的同伦。
由此得到一个新范畴 $\text{hSimp}(\mathcal{C})$：它与
$\text{Simp}(\mathcal{C})$ 有相同的对象，而其态射是
$\text{Simp}(\mathcal{C})$ 中态射的同伦类。因此有典范函子
$$
\text{Simp}(\mathcal{C})
\longrightarrow
\text{hSimp}(\mathcal{C})
$$
它本质满，并且在态射集上满射。
\end{remark}

\begin{definition}
\label{simplicial-definition-homotopy-equivalent}
设 $U$ 和 $V$ 为范畴 $\mathcal{C}$ 的两个单纯对象。
如果存在态射 $b : V \to U$，使得 $a \circ b$ 同伦于
$\text{id}_V$，且 $b \circ a$ 同伦于 $\text{id}_U$，
则称态射 $a : U \to V$ 为 {\it 同伦等价}。
如果存在同伦等价 $a : U \to V$，则称 $U$ 和 $V$
{\it 同伦等价}。
\end{definition}

\begin{example}
\label{simplicial-example-simplex-contractible}
单纯集 $\Delta[m]$ 同伦等价于 $\Delta[0]$。事实上，考虑唯一态射
$f : \Delta[m] \to \Delta[0]$，以及由 $\Delta[m]$ 的最后一个
$0$-单形的包含给出的态射 $g : \Delta[0] \to \Delta[m]$。
有 $f \circ g = \text{id}$。下面给出从
$\text{id}_{\Delta[m]}$ 到 $g \circ f$ 的同伦
$h : \Delta[m] \times \Delta[1] \to \Delta[m]$。确切地说，
$h$ 由映射
$$
\Mor_\Delta([n], [m]) \times \Mor_\Delta([n], [1]) \to \Mor_\Delta([n], [m])
$$
给出，它们把 $(\varphi, \alpha)$ 映到
$$
k \mapsto
\left\{
\begin{matrix}
\varphi(k) & \text{若} & \alpha(k) = 0 \\
m & \text{若} & \alpha(k) = 1
\end{matrix}
\right.
$$
。注意，这只因我们取 $g$ 为最后一个 $0$-单形的包含而成立。
若取 $g$ 为第一个 $0$-单形的包含，则可找到从
$g \circ f$ 到 $\text{id}_{\Delta[m]}$ 的同伦。
这说明了单纯集范畴中同伦所固有的不对称性。
\end{example}

\noindent
下面的引理说明 $U \times \Delta[1]$ 同伦等价于 $U$。

\begin{lemma}
\label{simplicial-lemma-contractible}
设 $\mathcal{C}$ 为有有限余积的范畴，$U$ 为 $\mathcal{C}$
的单纯对象。考虑映射
$e_1, e_0 : U \to U \times \Delta[1]$ 和
$\pi : U \times \Delta[1] \to U$；见引理
\ref{simplicial-lemma-back-to-U}。
\begin{enumerate}
\item $\pi \circ e_1 = \pi \circ e_0 = \text{id}_U$；
\item 态射 $\text{id}_{U \times \Delta[1]}$ 与
$e_0 \circ \pi$ 同伦；
\item 态射 $\text{id}_{U \times \Delta[1]}$ 与
$e_1 \circ \pi$ 同伦。
\end{enumerate}
\end{lemma}

\begin{proof}
第一个断言平凡。考虑单纯集映射
$\Delta[1] \times \Delta[1] \longrightarrow \Delta[1]$，
它在次数 $n$ 上把一对 $(\beta_1, \beta_2)$，
其中 $\beta_i : [n] \to [1]$，映到由规则
$$
\beta(i) = \max\{\beta_1(i), \beta_2(i)\}.
$$
定义的态射 $\beta : [n] \to [1]$。
这是单纯集的态射，因为作用
$\Delta[1](\varphi) : \Delta[1]_n \to \Delta[1]_m$
（其中 $\varphi : [m] \to [n]$）是前复合。使用上面引入的记号
$\Delta[1]_n = \{\alpha^n_0, \ldots, \alpha^n_{n + 1}\}$，
若 $\beta_2 = \alpha^n_0$，则 $\beta = \alpha^n_0$；
若 $\beta_2 = \alpha^n_{n + 1}$，则 $\beta = \beta_1$。
由于 $\alpha^n_0$ 与 $\alpha^n_{n + 1}$ 分别是取值为
$1$ 与 $0$ 的常值映射，这意味着引理
\ref{simplicial-lemma-back-to-U} 的诱导态射
$$
U \times \Delta[1] \times \Delta[1]
\longrightarrow
U \times \Delta[1]
$$
是从 $\text{id}_{U \times \Delta[1]}$ 到
$e_1 \circ \pi$ 的同伦。$e_0 \circ \pi$ 的情形类似
（使用最小值而不是最大值）。
\end{proof}

\begin{lemma}
\label{simplicial-lemma-fibre-products-simplicial-object-w-section}
设 $f : Y \to X$ 为有纤维积的范畴 $\mathcal{C}$ 中的态射，
并假设 $f$ 有截面 $s$。考虑例
\ref{simplicial-example-fibre-products-simplicial-object} 中从 $f$
构造的单纯对象 $U$。态射 $U \to U$ 在每个次数上都是
$Y \times_X \ldots \times_X Y$ 的自映射
$(s \circ f)^{n + 1}$，即在每个因子上由 $s \circ f$ 给出；
这个态射同伦于 $U$ 上的恒等映射。特别地，$U$ 同伦等价于
常值单纯对象 $X$。
\end{lemma}

\begin{proof}
令 $g^0 = \text{id}_Y$ 且 $g^1 = s \circ f$。使用态射
\begin{eqnarray*}
Y \times_X \ldots \times_X Y \times \Mor([n], [1])
& \to &
Y \times_X \ldots \times_X Y \\
(y_0, \ldots, y_n) \times \alpha
& \mapsto &
(g^{\alpha(0)}(y_0), \ldots, g^{\alpha(n)}(y_n))
\end{eqnarray*}
这里采用函子点观点来定义这些映射。
% P02-E0134: frozen English repeats “point” in “functor of points point of view”.
换言之，
$h_{n, 0} = \text{id}$、
$h_{n, n + 1} = (s \circ f)^{n + 1}$，且
$h_{n, i} = \text{id}_Y^{i + 1} \times (s \circ f)^{n + 1 - i}$。
留给读者证明它们满足引理 \ref{simplicial-lemma-relations-homotopy} 的关系，
因而定义所需同伦。另见备注 \ref{simplicial-remark-homotopy-better}；
它说明不必对范畴 $\mathcal{C}$ 作任何其他假设。
\end{proof}

\begin{lemma}
\label{simplicial-lemma-products-homotopy}
设 $\mathcal{C}$ 为范畴，$T$ 为集合。对 $t \in T$，设
$X_t$、$Y_t$ 为 $\mathcal{C}$ 的单纯对象，并假设
$X = \prod_{t \in T} X_t$ 和 $Y = \prod_{t \in T} Y_t$ 存在。
\begin{enumerate}
\item 若对所有 $t \in T$，$X_t$ 与 $Y_t$ 同伦等价，且 $T$
有限，则 $X$ 与 $Y$ 同伦等价。
\end{enumerate}
对 $t \in T$，设 $a_t, b_t : X_t \to Y_t$ 为态射，并令
$a = \prod a_t : X \to Y$ 及 $b = \prod b_t : X \to Y$。
\begin{enumerate}
\item[(2)] 若对所有 $t \in T$ 都存在从 $a_t$ 到 $b_t$ 的同伦，
则存在从 $a$ 到 $b$ 的同伦。
\item[(3)] 若 $T$ 有限，且对 $t \in T$，
$a_t, b_t : X_t \to Y_t$ 同伦，则 $a$ 与 $b$ 同伦。
\end{enumerate}
\end{lemma}

\begin{proof}
若 $h_t = (h_{t, n , i})$ 是从 $a_t$ 到 $b_t$ 的同伦
（见备注 \ref{simplicial-remark-homotopy-better}），则
$h = (\prod_t h_{t, n, i})$ 是从 $\prod a_t$ 到
$\prod b_t$ 的同伦。这证明 (2)。

\medskip\noindent
证明 (3)。选取 $t \in T$。存在整数 $n \geq 0$ 和链
$a_t = a_{t, 0}, a_{t, 1}, \ldots, a_{t, n} = b_t$，
使得对每个 $1 \leq i \leq n$，或者存在从
$a_{t, i - 1}$ 到 $a_{t, i}$ 的同伦，或者存在从
$a_{t, i}$ 到 $a_{t, i - 1}$ 的同伦。若 $n = 0$，
就另取一个 $t$。（若对所有 $t \in T$ 都有 $a_t = b_t$，
则已经完成。）所以假设 $n > 0$。由例
\ref{simplicial-example-trivial-homotopy}，对所有
$t' \in T \setminus \{t\}$，存在从 $b_{t'}$ 到 $b_{t'}$
的同伦。
% P02-E0135: frozen English says “there are is a homotopy”.
因此由 (2)，或者存在从
$a_{t, n - 1} \times \prod_{t'} b_{t'}$ 到 $b$ 的同伦，
或者存在从 $b$ 到
$a_{t, n - 1} \times \prod_{t'} b_{t'}$ 的同伦。
如此可将 $n$ 减少 $1$。这证明 (3)。

\medskip\noindent
(1) 由 (3) 和定义得出。
\end{proof}




\section{阿贝尔范畴中的同伦}
\label{simplicial-section-homotopy-abelian}

\noindent
设 $\mathcal{A}$ 为加性范畴，$U$、$V$ 为 $\mathcal{A}$
的单纯对象，$a, b : U \to V$ 为态射。进一步假设
$h : U \times \Delta[1] \to V$ 是从 $a$ 到 $b$ 的同伦。
下面证明链复形的两个态射
$s(a), s(b) : s(U) \longrightarrow s(V)$
在《同调代数》第 \ref{homology-section-complexes} 节的意义下同伦。
使用第 \ref{simplicial-section-homotopy} 节引入的记号，按公式
\begin{equation}
\label{simplicial-equation-homotopy-to-homotopy}
s(h)_n = \sum\nolimits_{i = 0}^n (-1)^{i + 1} h_{n + 1, i + 1} \circ s^n_i.
\end{equation}
定义
$$
s(h)_n : U_n \longrightarrow V_{n + 1}
$$
。计算 $d_{n + 1} \circ s(h)_n + s(h)_{n - 1} \circ d_n$。
首先，
\begin{eqnarray*}
d_{n + 1} \circ s(h)_n & = &
\sum\nolimits_{j = 0}^{n + 1}
\sum\nolimits_{i = 0}^n
(-1)^{j + i + 1}
d^{n + 1}_j \circ h_{n + 1, i + 1} \circ s^n_i \\
& = &
\sum\nolimits_{1 \leq i + 1 \leq j \leq n + 1}
(-1)^{j + i + 1}
h_{n, i + 1} \circ d^{n + 1}_j \circ s^n_i \\
& &
+
\sum\nolimits_{n \geq i \geq j \geq 0}
(-1)^{i + j + 1}
h_{n, i} \circ d^{n + 1}_j \circ s^n_i \\
& = &
\sum\nolimits_{1 \leq i + 1 < j \leq n + 1}
(-1)^{j + i + 1}
h_{n, i + 1} \circ s^{n - 1}_i \circ d^n_{j - 1} \\
& &
+
\sum\nolimits_{1 \leq i + 1 = j \leq n + 1}
(-1)^{j + i + 1}
h_{n, i + 1} \\
& &
+
\sum\nolimits_{n \geq i = j \geq 0}
(-1)^{i + j + 1}
h_{n, i} \\
& &
+
\sum\nolimits_{n \geq i > j \geq 0}
(-1)^{i + j + 1}
h_{n, i} \circ s^{n - 1}_{i - 1} \circ d^n_j
\end{eqnarray*}
留给读者验证：四个和式中的第一个与最后一个，恰好同
$$
s(h)_{n - 1} \circ d_n
=
\sum_{i = 0}^{n - 1} \sum_{j = 0}^n
(-1)^{i + 1 + j} h_{n, i + 1} \circ s^{n - 1}_i \circ d^n_j.
$$
的所有项相消。因此得到
\begin{eqnarray*}
d_{n + 1} \circ s(h)_n + s(h)_{n - 1} \circ d_n
& = &
\sum_{j = 1}^{n + 1} (-1)^{2j} h_{n, j}
+
\sum_{i = 0}^n (-1)^{2i + 1} h_{n, i} \\
& = &
h_{n, n + 1} - h_{n , 0} \\
& = &
a_n - b_n
\end{eqnarray*}
，正是所需结论。

\begin{lemma}
\label{simplicial-lemma-homotopy-s-N}
设 $\mathcal{A}$ 为加性范畴，$a, b : U \to V$ 为
$\mathcal{A}$ 中单纯对象的态射。若 $a$、$b$ 同伦，则
$s(a), s(b) : s(U) \to s(V)$ 是链复形的同伦态射。
若 $\mathcal{A}$ 为阿贝尔范畴，则
$N(a), N(b) : N(U) \to N(V)$ 也是链复形的同伦态射。
\end{lemma}

\begin{proof}
可以选取一列从 $U$ 到 $V$ 的态射
$a = a_0, a_1, \ldots, a_n = b$，使得对每个
$i = 1, \ldots, n$，或者存在从 $a_i$ 到 $a_{i - 1}$
的同伦，或者存在从 $a_{i - 1}$ 到 $a_i$ 的同伦。
上面的计算说明，在此情形下，或者 $s(a_i)$ 作为链复形映射
同伦于 $s(a_{i - 1})$，或者 $s(a_{i - 1})$ 作为链复形映射
同伦于 $s(a_i)$。当然，这两个断言等价；而且同伦是链复形映射集
上的等价关系，见《同调代数》第 \ref{homology-section-complexes} 节。
这证明 $s(a)$ 和 $s(b)$ 作为链复形映射同伦。

\medskip\noindent
下面处理 $N(a)$ 和 $N(b)$。由引理
\ref{simplicial-lemma-decompose-associated-complexes}，
$N(a)$、$N(b)$ 是复合
$$
N(U) \to s(U) \to s(V) \to N(V)
$$
，其中中间使用 $s(a)$、$s(b)$。因此，断言由《同调代数》引理
\ref{homology-lemma-compose-homotopy} 得出。
\end{proof}

\begin{lemma}
\label{simplicial-lemma-homotopy-equivalence-s-N}
设 $\mathcal{A}$ 为加性范畴，$a : U \to V$ 为
$\mathcal{A}$ 中单纯对象的态射。若 $a$ 是同伦等价，则
$s(a) : s(U) \to s(V)$ 是链复形的同伦等价。
若进一步 $\mathcal{A}$ 为阿贝尔范畴，则
$N(a) : N(U) \to N(V)$ 也是链复形的同伦等价。
\end{lemma}

\begin{proof}
略。见上面的引理 \ref{simplicial-lemma-homotopy-s-N}。
\end{proof}




\section{同伦与余单纯对象}
\label{simplicial-section-homotopy-cosimplicial}

\noindent
设 $\mathcal{C}$ 为有有限积的范畴，$V$ 为余单纯对象，并考虑
$\Hom(\Delta[1], V)$；见第
\ref{simplicial-section-hom-from-simplicial-sets-into-cosimplicial} 节。
态射 $e_0, e_1 : \Delta[0] \to \Delta[1]$ 给出两个态射
$e_0, e_1 : \Hom(\Delta[1], V) \to V$。

\begin{definition}
\label{simplicial-definition-homotopy-cosimplicial}
设 $\mathcal{C}$ 为有有限积的范畴，$U$ 和 $V$ 为
$\mathcal{C}$ 的两个余单纯对象，$a, b : U \to V$ 为
$\mathcal{C}$ 中余单纯对象的两个态射。
\begin{enumerate}
\item 若态射
$$
h : U \longrightarrow \Hom(\Delta[1], V)
$$
满足 $a = e_0 \circ h$ 且 $b = e_1 \circ h$，则称它为
{\it 从 $a$ 到 $b$ 的同伦}。
\item 如果存在一列从 $U$ 到 $V$ 的态射
$a = a_0, a_1, \ldots, a_n = b$，使得对每个
$i = 1, \ldots, n$，或者存在从 $a_i$ 到 $a_{i - 1}$ 的同伦，
或者存在从 $a_{i - 1}$ 到 $a_i$ 的同伦，则称 $a$ 和 $b$
{\it 同伦}，或者称它们 {\it 属于同一个同伦类}。
\end{enumerate}
\end{definition}

\noindent
这与第 \ref{simplicial-section-homotopy} 节为单纯对象引入的概念对偶。
为作说明，考虑定义中的从 $a$ 到 $b$ 的同伦
$h : U \to \Hom(\Delta[1], V)$。回顾 $\Delta[1]_n$
是有限集。$h$ 的次数 $n$ 分量是态射
$$
h_n = (h_{n, \alpha}) :
U
\longrightarrow
\Hom(\Delta[1], V)_n =
\prod\nolimits_{\alpha \in \Delta[1]_n} V_n
$$
% P02-E0136: frozen degree-n component has source U rather than U_n, contradicted immediately by h_{n,alpha}: U_n -> V_n.
$\mathcal{C}$ 的态射 $h_{n, \alpha} : U_n \to V_n$ 具有如下性质：
对 $\Delta$ 的每个态射 $f : [n] \to [m]$，有
\begin{equation}
\label{simplicial-equation-property-homotopy-cosimplicial}
h_{m, \alpha} \circ U(f) = V(f) \circ h_{n, \alpha \circ f}
\end{equation}
此外，条件 $a = e_0 \circ h$ 意味着
$a_n = h_{n, 0 : [n] \to [1]}$，其中
$0 : [n] \to [1]$ 是取值为 $0$ 的常值映射。
同样，条件 $b = e_1 \circ h$ 意味着
$b_n = h_{n, 1 : [n] \to [1]}$，其中
$1 : [n] \to [1]$ 是取值为 $1$ 的常值映射。
反之，给定一族态射 $\{h_{n, \alpha}\}$，使得
(\ref{simplicial-equation-property-homotopy-cosimplicial})
对 $\Delta$ 的所有态射 $f$ 都成立，并且对所有 $n \geq 0$ 都有
$a_n = h_{n, 0 : [n] \to [1]}$ 和
$b_n = h_{n, 1 : [n] \to [1]}$，那么令
$h = \prod_{\alpha \in \Delta[1]_n} h_{n, \alpha}$，
便得到从 $a$ 到 $b$ 的同伦 $h$。

\begin{remark}
\label{simplicial-remark-homotopy-cosimplicial-better}
设 $\mathcal{C}$ 为任意范畴（不作任何假设），$U$ 和 $V$
为 $\mathcal{C}$ 的余单纯对象，$a, b : U \to V$ 为
$\mathcal{C}$ 中余单纯对象的态射。一个 {\it 从 $a$ 到 $b$ 的同伦}
由态射 $h_{n, \alpha} : U_n \to V_n$ 给出，其中
$n \geq 0$ 且 $\alpha \in \Delta[1]_n$；要求它们对
$\Delta$ 的所有态射 $f$ 都满足
(\ref{simplicial-equation-property-homotopy-cosimplicial})，
并且对所有 $n \geq 0$ 都有
$a_n = h_{n, 0 : [n] \to [1]}$ 与
$b_n = h_{n, 1 : [n] \to [1]}$。
与定义 \ref{simplicial-definition-homotopy-cosimplicial} 一样，
如果存在一列从 $U$ 到 $V$ 的态射
$a = a_0, a_1, \ldots, a_n = b$，使得对每个
$i = 1, \ldots, n$，或者存在从 $a_{i - 1}$ 到 $a_i$ 的同伦，
或者存在从 $a_i$ 到 $a_{i - 1}$ 的同伦，则称态射 $a$ 和 $b$
{\it 同伦}。显然，若 $F : \mathcal{C} \to \mathcal{C}'$
为任意函子，且 $\{h_{n, i}\}$ 是从 $a$ 到 $b$ 的同伦，
% P02-E0137: frozen cosimplicial remark switches the family index from alpha to i.
则 $\{F(h_{n, i})\}$ 是从 $F(a)$ 到 $F(b)$ 的同伦。
同样，若 $a$ 和 $b$ 同伦，则 $F(a)$ 和 $F(b)$ 同伦。
当有限积存在时，这个新概念与旧概念相同。因此特别可知，
只要两个范畴都有有限积，函子就保持原来的概念。
\end{remark}

\begin{lemma}
\label{simplicial-lemma-compare-homotopies}
设 $\mathcal{C}$ 为范畴，$U$ 和 $V$ 为 $\mathcal{C}$
的两个余单纯对象，$a, b : U \to V$ 为余单纯对象的态射。
回顾 $U$、$V$ 分别对应于 $\mathcal{C}^{opp}$ 的单纯对象
$U'$、$V'$，而 $a, b$ 分别对应于态射
$a', b' : V' \to U'$。下列条件等价：
\begin{enumerate}
\item 存在备注 \ref{simplicial-remark-homotopy-cosimplicial-better} 意义下，
从 $a$ 到 $b$ 的同伦 $h = \{h_{n, \alpha}\}$。
\item 存在备注 \ref{simplicial-remark-homotopy-better} 意义下，
从 $a'$ 到 $b'$ 的同伦 $h = \{h_{n, i}\}$。
\end{enumerate}
因此，$a$ 与 $b$ 在备注
\ref{simplicial-remark-homotopy-cosimplicial-better} 的意义下同伦，
当且仅当 $a'$ 与 $b'$ 在备注 \ref{simplicial-remark-homotopy-better}
的意义下同伦。
\end{lemma}

\begin{proof}
若 $\mathcal{C}$ 有有限积，则 $\mathcal{C}^{opp}$ 有有限余积，
并且可以用定义 \ref{simplicial-definition-homotopy-cosimplicial}
和 \ref{simplicial-definition-homotopy} 代替备注
\ref{simplicial-remark-homotopy-cosimplicial-better}
和 \ref{simplicial-remark-homotopy-better}。
此时，$h : U \to \Hom(\Delta[1], V)$ 等同于态射
$h' : \Hom(\Delta[1], V)' \to U'$。由于积与余积也相互交换，
立即有 $(\Hom(\Delta[1], V))' = V' \times \Delta[1]$。
此外，态射 $e_0, e_1 : \Hom(\Delta[1], V) \to V$
的“加撇”版本是态射
$e_0, e_1 : V' \to \Delta[1] \times V$。
因此，$e_0 \circ h = a$ 转化为 $h' \circ e_0 = a'$，
同样地，$e_1 \circ h = b$ 转化为 $h' \circ e_1 = b'$。
这证明了这种情形的引理。

\medskip\noindent
一般情形下，需要把
(\ref{simplicial-equation-property-homotopy-cosimplicial})
给出的关系转化为引理 \ref{simplicial-lemma-relations-homotopy} 中的关系。
细节从略。

\medskip\noindent
最后一个断言形式地由 (1) 与 (2) 的等价性得出。
\end{proof}

\begin{lemma}
\label{simplicial-lemma-functorial-homotopy}
\begin{slogan}
函子保持（余）单纯对象的同伦态射。
\end{slogan}
设 $\mathcal{C}, \mathcal{C}', \mathcal{D}, \mathcal{D}'$ 为范畴。
术语如备注 \ref{simplicial-remark-homotopy-cosimplicial-better} 和
\ref{simplicial-remark-homotopy-better} 所述。
% P02-E0138: frozen source presents “With terminology as in ...” as a sentence fragment.
\begin{enumerate}
\item 设 $a, b : U \to V$ 为 $\mathcal{D}$ 中单纯对象的态射，
$F : \mathcal{D} \to \mathcal{D}'$ 为协变函子。若 $a$ 与 $b$
同伦，则 $F(a)$、$F(b)$ 是单纯对象之间同伦的态射
$F(U) \to F(V)$。
\item 设 $a, b : U \to V$ 为 $\mathcal{C}$ 中余单纯对象的态射，
$F : \mathcal{C} \to \mathcal{C}'$ 为协变函子。若 $a$ 与 $b$
同伦，则 $F(a)$、$F(b)$ 是余单纯对象之间同伦的态射
$F(U) \to F(V)$。
\item 设 $a, b : U \to V$ 为 $\mathcal{D}$ 中单纯对象的态射，
$F : \mathcal{D} \to \mathcal{C}$ 为反变函子。若 $a$ 与 $b$
同伦，则 $F(a)$、$F(b)$ 是余单纯对象之间同伦的态射
$F(V) \to F(U)$。
\item 设 $a, b : U \to V$ 为 $\mathcal{C}$ 中余单纯对象的态射，
$F : \mathcal{C} \to \mathcal{D}$ 为反变函子。若 $a$ 与 $b$
同伦，则 $F(a)$、$F(b)$ 是单纯对象之间同伦的态射
$F(V) \to F(U)$。
\end{enumerate}
\end{lemma}

\begin{proof}
由上面的引理 \ref{simplicial-lemma-compare-homotopies}，可以把 $F$
变为一对范畴之间的协变函子，而只须证明函子保持同伦的映射对。
这已在备注 \ref{simplicial-remark-homotopy-better} 中说明。
\end{proof}

\begin{lemma}
\label{simplicial-lemma-push-outs-simplicial-object-w-section}
设 $f : X \to Y$ 为有推出的范畴 $\mathcal{C}$ 中的态射，
并假设存在态射 $s : Y \to X$，满足
$s \circ f = \text{id}_X$。考虑例
\ref{simplicial-example-push-outs-simplicial-object} 中从 $f$ 构造的
余单纯对象 $U$。态射 $U \to U$ 在每个次数上都是
$Y \amalg_X \ldots \amalg_X Y$ 的自映射，并且在每个因子上
由 $f \circ s$ 给出；它同伦于 $U$ 上的恒等映射。
特别地，$U$ 同伦等价于常值余单纯对象 $X$。
\end{lemma}

\begin{proof}
本引理与引理
\ref{simplicial-lemma-fibre-products-simplicial-object-w-section} 对偶。
因此，它由引理 \ref{simplicial-lemma-compare-homotopies} 得出。
\end{proof}

\begin{lemma}
\label{simplicial-lemma-homotopy-s-Q}
\begin{slogan}
（余单纯）Dold--Kan 函子把同伦映射映为同伦映射。
\end{slogan}
设 $\mathcal{A}$ 为加性范畴，$a, b : U \to V$ 为
$\mathcal{A}$ 中余单纯对象的态射。若 $a$、$b$ 同伦，则
$s(a), s(b) : s(U) \to s(V)$ 是上链复形的同伦映射。
若进一步 $\mathcal{A}$ 为阿贝尔范畴，则
$Q(a), Q(b) : Q(U) \to Q(V)$ 是上链复形的同伦映射。
\end{lemma}

\begin{proof}
设 $(-)' : \mathcal{A} \to \mathcal{A}^{opp}$ 为反变函子
$A \mapsto A$。根据引理
\ref{simplicial-lemma-push-outs-simplicial-object-w-section}，
映射 $a'$ 和 $b'$ 同伦。
% P02-E0139: frozen proof cites the pushout-with-section lemma here, but homotopy transfer to opposite categories is established by lemma-compare-homotopies.
由引理 \ref{simplicial-lemma-homotopy-s-N}，$s(a')$ 和 $s(b')$
是链复形的同伦映射。由于 $s(a') = (s(a))'$ 且
$s(b') = (s(b))'$，应用加性反变函子
$(-)'' : \mathcal{A}^{opp} \to \mathcal{A}$ 可知
$s(a)$ 和 $s(b)$ 也同伦。对 $Q$-复形的结论以相同方式得出，
这里使用 $Q(U)' = N(U')$。
\end{proof}

\begin{lemma}
\label{simplicial-lemma-homotopy-equivalence-s-Q}
设 $\mathcal{A}$ 为加性范畴，$a : U \to V$ 为
$\mathcal{A}$ 中余单纯对象的态射。若 $a$ 是同伦等价，则
$s(a) : s(U) \to s(V)$ 是链复形的同伦等价。
若进一步 $\mathcal{A}$ 为阿贝尔范畴，则
$Q(a) : Q(U) \to Q(V)$ 也是链复形的同伦等价。
% P02-E0140: frozen source calls both targets chain complexes although s(U) and Q(U) were defined here as cochain complexes.
\end{lemma}

\begin{proof}
略。见上面的引理 \ref{simplicial-lemma-homotopy-s-Q}。
\end{proof}




\section{阿贝尔范畴中的更多同伦}
\label{simplicial-section-backwards-homotopy}

\noindent
设 $\mathcal{A}$ 为阿贝尔范畴。本节将证明：
$\text{Ch}_{\geq 0}(\mathcal{A})$ 中态射之间的同伦，
总是来自单纯对象范畴中的态射 $U \times \Delta[1] \to V$。
在某种意义下，这将给出引理 \ref{simplicial-lemma-homotopy-s-N} 的逆命题。
先发展一些关于链复形态射之间同伦的材料。

\begin{lemma}
\label{simplicial-lemma-represent-homotopy}
设 $\mathcal{A}$ 为阿贝尔范畴，$A$ 为链复形。考虑协变函子
$$
B \longmapsto
\{
(a, b, h)
\mid
a, b : A \to B\text{ 且 }h\text{ 下列二者之间的同伦： }a, b
\}
$$
存在链复形 $\diamond A$，使得
$\Mor_{\text{Ch}(\mathcal{A})}(\diamond A, -)$
同构于所显示的函子。构造 $A \mapsto \diamond A$ 具有函子性。
\end{lemma}

\begin{proof}
令 $\diamond A_n = A_n \oplus A_n \oplus A_{n - 1}$，
并以矩阵
$$
d_{\diamond A, n}
=
\left(
\begin{matrix}
d_{A, n} & 0 & \text{id}_{A_{n - 1}} \\
0 & d_{A, n} & -\text{id}_{A_{n - 1}} \\
0 & 0 & -d_{A, n - 1}
\end{matrix}
\right) :
A_n \oplus A_n \oplus A_{n - 1} \to
A_{n - 1} \oplus A_{n - 1} \oplus A_{n - 2}
$$
定义 $d_{\diamond A, n}$。若 $\mathcal{A}$ 是阿贝尔群范畴，
且 $(x, y, z) \in A_n \oplus A_n \oplus A_{n - 1}$，则
$d_{\diamond A, n}(x, y, z) =
(d_n(x) + z, d_n(y) - z, -d_{n - 1}(z))$。
容易验证 $d^2 = 0$。显然有两个映射
$\diamond a, \diamond b : A \to \diamond A$
（第一与第二直和分量），以及映射
$\diamond A \to A[-1]$，它们给出逐项分裂的短正合序列
$$
0 \to
A \oplus A \to
\diamond A \to
A[-1] \to
0
$$
。此外，有一列映射
$\diamond h_n : A_n \to \diamond A_{n + 1}$，
即从 $A_n$ 到 $\diamond A_{n + 1}$ 的分量 $A_n$ 的恒等映射，
使得 $\diamond h$ 是 $\diamond a$ 与 $\diamond b$ 之间的同伦。

\medskip\noindent
由此，任意态射 $f : \diamond A \to B$ 都通过令
$a = f \circ \diamond a$、$b = f \circ \diamond b$ 以及
$h_n = f_{n + 1} \circ \diamond h_n$ 给出三元组 $(a, b, h)$。
反之，给定三元组 $(a, b, h)$，取
$$
f_n = (a_n, b_n, h_{n - 1}).
$$
即可得到态射 $f : \diamond A \to B$。
要看出这是链复形的态射，需要作一次计算。这里只在
$\mathcal{A}$ 为阿贝尔群范畴的情形下计算：
设 $(x, y, z) \in \diamond A_n = A_n \oplus A_n \oplus A_{n - 1}$。
则
\begin{eqnarray*}
f_{n - 1}(d_n(x, y, z)) & = &
f_{n - 1}(d_n(x) + z, d_n(y) - z, -d_{n - 1}(z)) \\
& = &
a_n(d_n(x)) + a_n(z) + b_n(d_n(y)) - b_n(z) - h_{n - 2}(d_{n - 1}(z))
\end{eqnarray*}
% P02-E0141: frozen calculation uses degree-n components a_n and b_n on degree-(n-1) inputs.
以及
\begin{eqnarray*}
d_n(f_n(x, y, z) & = &
% P02-E0142: frozen display leaves the parenthesis in d_n(f_n(x,y,z) unclosed.
d_n(a_n(x) + b_n(y) + h_{n - 1}(z)) \\
& = &
d_n(a_n(x)) + d_n(b_n(y)) + d_n(h_{n - 1}(z))
\end{eqnarray*}
根据同伦的定义，二者相同。
\end{proof}

\noindent
注意，扩张
$$
0 \to
A \oplus A \to
\diamond A \to
A[-1] \to
0
$$
配有态射 $\diamond A_n \to A[-1]_n$ 的截面，并且相应的态射
$\delta : A[-1] \to (A \oplus A)[-1]$（见《同调代数》引理
\ref{homology-lemma-ses-termwise-split}）等于态射
$(1, -1) : A[-1] \to A[-1] \oplus A[-1]$。

\begin{lemma}
\label{simplicial-lemma-map-into-diamond}
设 $\mathcal{A}$ 为阿贝尔范畴，并设
$$
0 \to A \oplus A \to B \to C \to 0
$$
为 $\mathcal{A}$ 中链复形的短正合序列。进一步假设给定态射
$s_n : C_n \to B_n$，它们在次数 $n$ 上分裂相应的短正合序列。
设 $\delta(s) : C \to (A \oplus A)[-1] = A[-1] \oplus A[-1]$
为相应的复形态射；见《同调代数》引理
\ref{homology-lemma-ses-termwise-split}。若 $\delta(s)$ 经态射
$(1, -1) : A[-1] \to A[-1] \oplus A[-1]$ 分解，则存在唯一态射
$B \to \diamond A$，使其嵌入交换图
$$
\xymatrix{
0 \ar[r] &
A \oplus A \ar[d] \ar[r] &
B \ar[r] \ar[d] &
C \ar[d] \ar[r] &
0 \\
0 \ar[r] &
A \oplus A \ar[r] &
\diamond A \ar[r] &
A[-1] \ar[r] &
0
}
$$
，其中竖直映射也分别与截面 $s_n$ 及
$\diamond A_n \to A[-1]_n$ 的截面相容。
\end{lemma}

\begin{proof}
以 $(p_n, q_n) : B_n \to A_n \oplus A_n$ 表示《同调代数》引理
\ref{homology-lemma-ses-termwise-split} 中的态射 $\pi_n$。
又将短正合序列中的映射写成
$(a, b) : A \oplus A \to B$ 和 $r : B \to C$。
把 $\delta(s)$ 的分解写成
$\delta(s) = (1, -1) \circ f$。这意味着
$p_{n - 1} \circ d_{B, n} \circ s_n = f_n$，并且
$q_{n - 1} \circ d_{B, n} \circ s_n = - f_n$，并且
% P02-E0143: frozen source leaves a dangling “and” before the imperative “Set”.
令 $B_n \to \diamond A_n = A_n \oplus A_n \oplus A_{n - 1}$
等于 $(p_n, q_n, f_n \circ r_n)$。

\medskip\noindent
现在必须检验这确实定义了复形态射。这里只在阿贝尔群情形下验证。
取 $x \in B_n$。则
$x = a_n(x_1) + b_n(x_2) + s_n(x_3)$，并且只须分别对每一项证明
我们的定义与微分交换。对于项 $a_n(x_1)$，有
$(p_n, q_n, f_n \circ r_n)(a_n(x_1)) = (x_1, 0, 0)$，
结论显然。项 $b_n(x_2)$ 的情形类似。对于项 $s_n(x_3)$，有
\begin{eqnarray*}
(p_n, q_n, f_n \circ r_n)(d_n(s_n(x_3))) & = &
(p_n, q_n, f_n \circ r_n)( \\
& & \ \ \ \ \ a_n(f_n(x_3)) - b_n(f_n(x_3)) + s_n(d_n(x_3))) \\
& = &
(f_n(x_3), -f_n(x_3), f_n(d_n(x_3)))
\end{eqnarray*}
% P02-E0144: frozen first calculation applies degree-n projections to d_n(s_n(x_3)), which lies in degree n-1.
，这是由 $f_n$ 的定义得到的。另一方面，
\begin{eqnarray*}
d_n(p_n, q_n, f_n \circ r_n)(s_n(x_3)) & = & d_n(0, 0, f_n(x_3)) \\
& = &
(f_n(x_3), - f_n(x_3), d_{A[-1], n}(f_n(x_3)))
\end{eqnarray*}
由于 $f$ 是复形态射，结论成立。
\end{proof}

\begin{lemma}
\label{simplicial-lemma-backwards-homotopy}
设 $\mathcal{A}$ 为阿贝尔范畴，$U$、$V$ 为 $\mathcal{A}$
的单纯对象，$a, b : U \to V$ 为一对态射。假设相应的链复形映射
$N(a), N(b) : N(U) \to N(V)$ 通过同伦
$\{N_n : N(U)_n \to N(V)_{n + 1}\}$ 同伦。
则存在定义 \ref{simplicial-definition-homotopy} 意义下从 $a$ 到 $b$
的同伦。此外，可以选取同伦
$h : U \times \Delta[1] \to V$，使得
$N_n = N(h)_n$，其中 $N(h)$ 是按第
\ref{simplicial-section-homotopy-abelian} 节从 $h$ 得到的同伦。
\end{lemma}

\begin{proof}
设 $(\diamond N(U), \diamond a, \diamond b, \diamond h)$
如引理 \ref{simplicial-lemma-represent-homotopy} 及其证明中所述。
由该引理，存在表示三元组 $(N(a), N(b), \{N_n\})$ 的态射
$\diamond N(U) \to N(V)$。下面证明存在态射
$\psi : N(U \times \Delta[1]) \to \diamond{N(U)}$，
使得 $\diamond a = \psi \circ N(e_0)$ 且
$\diamond b = \psi \circ N(e_1)$。还将证明：
由 (\ref{simplicial-equation-homotopy-to-homotopy}) 及引理
\ref{simplicial-lemma-homotopy-s-N}，取
$h = \text{id}_{U \times \Delta[1]}$ 得到的
$N(e_0), N(e_1) :
N(U) \to N(U \times \Delta[1])$ 之间的同伦，经 $\psi$
映到映射
$\diamond a, \diamond b : N(U) \to \diamond{N(U)}$
之间的典范同伦 $\diamond h$。这当然蕴含引理。

\medskip\noindent
注意，作为函子，
$N : \text{Simp}(\mathcal{A}) \to \text{Ch}_{\geq 0}(\mathcal{A})$
是函子
$s : \text{Simp}(\mathcal{A}) \to \text{Ch}_{\geq 0}(\mathcal{A})$
的直和分量。并且函子 $\diamond$ 与直和相容。
因此，只须改为构造具有相应性质的态射
$\Psi : s(U \times \Delta[1]) \to \diamond{s(U)}$。
下面正是这一构造。

\medskip\noindent
由定义 \ref{simplicial-definition-homotopy}，态射
$e_0 : U \to U \times \Delta[1]$ 与
$e_1 : U \to U \times \Delta[1]$
通过同伦 $\text{id}_{U \times \Delta[1]}$ 同伦。
由引理 \ref{simplicial-lemma-homotopy-s-N}，得到链复形态射
$s(e_0) : s(U) \to s(U \times \Delta[1])$ 与
$s(e_1) : s(U) \to s(U \times \Delta[1])$ 之间的显式同伦
$\{h_n : s(U)_n \to s(U \times \Delta[1])_{n + 1}\}$。
由上面的引理 \ref{simplicial-lemma-map-into-diamond}，得到相应态射
$$
\Phi : \diamond{s(U)} \to s(U \times \Delta[1])
$$
根据构造，$\Phi_n$ 在 $\diamond{s(U)}_n$ 的分量
$s(U)[-1]_n = s(U)_{n - 1}$ 上的限制等于 $h_{n - 1}$。并且
$$
h_{n - 1} = \sum\nolimits_{i = 0}^{n - 1}
(-1)^{i + 1} s^n_i \cdot \alpha^n_{i + 1} :
U_{n - 1} \to \bigoplus\nolimits_j U_n \cdot \alpha^n_j.
$$
% P02-E0145: frozen formula uses s_i^n on U_{n-1}; the degeneracy with this source degree is s_i^{n-1}.
这里采用显然的记号。

\medskip\noindent
另一方面，态射 $e_i : U \to U \times \Delta[1]$ 诱导态射
$(e_0, e_1) : U \oplus U \to U \times \Delta[1]$。
以 $W$ 表示余核。注意，若写成
$(U \times \Delta[1])_n =
\bigoplus_{\alpha : [n] \to [1]} U_n \cdot \alpha$，
则可以把
$W_n = \bigoplus_{i = 1}^n U_n \cdot \alpha^n_i$
作如第 \ref{simplicial-section-homotopy} 节中的 $\alpha^n_i$ 那样的识别。
有交换图
$$
\xymatrix{
0 \ar[r] &
U \oplus U \ar[rd]_{(1, 1)} \ar[r] &
U \times \Delta[1] \ar[d]^\pi \ar[r] &
W \ar[r] &
0 \\
& & U & &
}
$$
这意味着应用函子 $s$ 后有类似的交换图。然后选取分裂
$\sigma_n : s(W)_n \to s(U \times \Delta[1])_n$：
通过 $(-1, 1)$，把分量 $U_n \cdot \alpha^n_i \subset W_n$
映到分量
$U_n \cdot \alpha^n_0 \oplus U_n \cdot \alpha^n_i
\subset (U \times \Delta[1])_n$。
注意 $s(\pi)_n \circ \sigma_n = 0$，从而
$(1, 1) \circ \delta(\sigma)_n = 0$。因此
$\delta(\sigma)$ 如引理 \ref{simplicial-lemma-map-into-diamond} 中那样分解。
由该引理，得到典范态射
$\Psi : s(U \times \Delta[1]) \to \diamond{s(U)}$。

\medskip\noindent
为了计算 $\Psi$，先计算态射
$\delta(\sigma) : s(W) \to s(U)[-1] \oplus s(U)[-1]$。
根据《同调代数》引理
\ref{homology-lemma-ses-termwise-split} 及其证明，
为此必须计算
% P02-E0146: frozen source has “we have compute”, omitting “to”.
$$
d_{s(U \times \delta[1]), n} \circ \sigma_n
-
\sigma_{n - 1} \circ d_{s(W), n}
$$
% P02-E0147: frozen source writes lowercase \delta[1] in three places where the simplicial interval is \Delta[1].
并把它写成到
$U_{n - 1} \cdot \alpha^{n - 1}_0
\oplus U_{n - 1} \cdot \alpha^{n - 1}_n$ 的态射。
这里只在 $\mathcal{A}$ 是阿贝尔群范畴的情形下计算。
采用简写 $x_{\alpha}$：对 $x \in U_n$，以它表示
$(U \times \Delta[1])_n$ 的分量
$U_n \cdot \alpha$ 中的元素 $x$。回顾
$$
d_{s(U \times \delta[1]), n}
=
\sum\nolimits_{i = 0}^n (-1)^i d^n_i
$$
其中 $d^n_i$ 通过单纯对象 $U$ 的态射 $d^n_i$，
把分量 $U_n \cdot \alpha$ 映到分量
$U_{n - 1} \cdot (\alpha \circ \delta^n_i)$。
用上述记号，这意味着
$$
d_{s(U \times \delta[1]), n}(x_\alpha) =
\sum\nolimits_{i = 0}^n (-1)^i (d^n_i(x))_{\alpha \circ \delta^n_i}
$$
从 $x_\alpha \in W_n$ 出发，换言之，对某个
$j \in \{1, \ldots, n\}$ 有 $\alpha = \alpha^n_j$，
可见 $\sigma_n(x_\alpha) = x_\alpha - x_{\alpha^n_0}$，因而
$$
(d_{s(U \times \delta[1]), n} \circ \sigma_n)(x_\alpha)
=
\sum\nolimits_{i = 0}^n (-1)^i (d^n_i(x))_{\alpha \circ \delta^n_i}
-
\sum\nolimits_{i = 0}^n (-1)^i (d^n_i(x))_{\alpha^n_0 \circ \delta^n_i}
$$
为了计算 $d_{s(W), n}(x_\alpha)$，必须删去所有满足
$\alpha \circ \delta^n_i = \alpha^{n - 1}_0, \alpha^{n - 1}_n$
的项。因此得到
$$
\begin{matrix}
(\sigma_{n - 1} \circ d_{s(W), n})(x_\alpha) = \\
\sum\nolimits_{i = 0, \ldots, n\text{ 且 }
\alpha \circ \delta^n_i \not = \alpha^{n - 1}_0\text{ 或 }\alpha^{n - 1}_n}
\Big((-1)^i (d^n_i(x))_{\alpha \circ \delta^n_i}
-
(-1)^i (d^n_i(x))_{\alpha^{n - 1}_0}
\Big)
\end{matrix}
$$
% P02-E0148: frozen exclusion predicate uses “not equal A or B”, which is logically nonrestrictive; excluding both boundary values requires conjunction.
显然，两项之差是和
$$
\sum\nolimits_{i = 0, \ldots, n\text{ 且 }
\alpha \circ \delta^n_i = \alpha^{n - 1}_0\text{ 或 }\alpha^{n - 1}_n}
\Big((-1)^i (d^n_i(x))_{\alpha \circ \delta^n_i}
-
(-1)^i (d^n_i(x))_{\alpha^{n - 1}_0}
\Big)
$$
当然，若 $\alpha \circ \delta^n_i = \alpha^{n - 1}_0$，
则该项消失。回顾对某个 $j \in \{1, \ldots, n\}$，
有 $\alpha = \alpha^n_j$。等式
$\alpha^n_j \circ \delta^n_i = \alpha^{n - 1}_n$
成立的唯一可能是 $j = n$ 且 $i = n$。因此，除非
$j = n$，实际所得为 $0$；在该情形下得到
$(-1)^n(d^n_n(x))_{\alpha^{n - 1}_n}
- (-1)^n(d^n_n(x))_{\alpha^{n - 1}_0}$。
换言之，态射
$$
\delta(\sigma)_n :
W_n
\to
(s(U)[-1] \oplus s(U)[-1])_n = U_{n - 1} \oplus U_{n - 1}
$$
在除 $U_n \cdot \alpha^n_n$ 以外的所有分量上为零，
而在该分量上等于
$((-1)^nd^n_n, -(-1)^nd^n_n)$。
（也就是说，两个分量中的第一个对应于
$\alpha^{n - 1}_n$ 所在的因子，因为这是把所有元素映到 $0$
的映射 $[n - 1] \to [1]$，因而对应于 $e_0$。）

\medskip\noindent
得到典范图
$$
\xymatrix{
0 \ar[r] &
s(U) \oplus s(U) \ar[r] \ar[d] &
\diamond{s(U)} \ar[r] \ar[d]^{\Phi}&
s(U)[-1] \ar[r] \ar[d] &
0 \\
0 \ar[r] &
s(U) \oplus s(U) \ar[r] \ar[d] &
s(U \times \Delta[1]) \ar[r] \ar[d]^\Psi &
s(W) \ar[r] \ar[d] &
0 \\
0 \ar[r] &
s(U) \oplus s(U) \ar[r] &
\diamond{s(U)} \ar[r] &
s(U)[-1] \ar[r] &
0
}
$$
断言 $\Phi \circ \Psi$ 是恒等映射。为此，只须证明
$\Phi$ 与 $\delta(\sigma)$ 的复合，作为映射
$s(U)[-1] \to s(W) \to s(U)[-1] \oplus s(U)[-1]$，
在第一个因子上是恒等映射，在第二个因子上是负恒等映射。
根据上面的计算，它是
$((-1)^nd^n_0, -(-1)^nd^n_0) \circ (-1)^n s^n_n = (1, -1)$，
% P02-E0149: frozen final composition uses d_0^n after the preceding computation produced d_n^n; its displayed degrees are also not composable with s_n^n.
正如所需。
\end{proof}









\section{平凡 Kan 纤维化}
\label{simplicial-section-trivial-kan}

\noindent
回顾当 $n \geq 0$ 时，单纯集 $\Delta[n]$ 由规则
$[k] \mapsto \Mor_\Delta([k], [n])$ 给出；见例
\ref{simplicial-example-simplex-simplicial-set}。回顾 $\Delta[n]$ 有唯一的
非退化 $n$-单形，并且所有非退化单形都是这个 $n$-单形的面。
事实上，$\Delta[n]$ 的非退化单形恰好对应于单射
$[k] \to [n]$，可以把它们与 $[n]$ 的子集等同起来。
此外，回顾对任意单纯集 $X$，有
$\Mor(\Delta[n], X) = X_n$（引理 \ref{simplicial-lemma-simplex-map}）。
令
$$
\partial \Delta[n] = i_{(n - 1)!}\text{sk}_{n - 1}\Delta[n]
$$
并称它为 $\Delta[n]$ 的 {\it 边界}。
由引理 \ref{simplicial-lemma-n-skeleton-sets} 可知，
$\partial \Delta[n] \subset \Delta[n]$ 是这样的单纯子集：
它在次数 $\leq n - 1$ 中具有相同的非退化单形，
但不含那个非退化 $n$-单形。

\begin{definition}
\label{simplicial-definition-trivial-kan}
若单纯集映射 $X \to Y$ 满足 $X_0 \to Y_0$ 为满射，并且对所有
$n \geq 1$ 及任意交换实线图
$$
\xymatrix{
\partial \Delta[n] \ar[r] \ar[d] & X \ar[d] \\
\Delta[n] \ar[r] \ar@{-->}[ru] & Y
}
$$
都存在使图交换的虚线箭头，则称它为 {\it 平凡 Kan 纤维化}。
\end{definition}

\noindent
平凡 Kan 纤维化满足一种非常一般的提升性质。

\begin{lemma}
\label{simplicial-lemma-trivial-kan}
设 $f : X \to Y$ 为单纯集的平凡 Kan 纤维化。
对任意单纯集交换实线图
$$
\xymatrix{
Z \ar[r]_b \ar[d] & X \ar[d] \\
W \ar[r]^a \ar@{-->}[ru] & Y
}
$$
，若 $Z \to W$（逐项）为单射，则存在使图交换的虚线箭头。
\end{lemma}

\begin{proof}
假设 $Z \not = W$。设 $n$ 为满足 $Z_n \not = W_n$
的最小整数。取 $x \in W_n$，且 $x \not \in Z_n$。
以 $Z' \subset W$ 表示包含 $Z$、$x$ 以及 $x$ 的所有退化的
单纯子集。设 $\varphi : \Delta[n] \to Z'$ 为与 $x$
对应的态射（引理 \ref{simplicial-lemma-simplex-map}）。
由于 $\partial \Delta[n]$ 的所有非退化单形都映入 $Z$，
$\varphi|_{\partial \Delta[n]}$ 映入 $Z$。
由假设，可将 $b \circ \varphi|_{\partial \Delta[n]}$
延拓为 $\beta : \Delta[n] \to X$。
由引理 \ref{simplicial-lemma-glue-simplex}，单纯集 $Z'$ 是
$\Delta[n]$ 与 $Z$ 沿 $\partial \Delta[n]$ 的推出。
因此 $b$ 和 $\beta$ 定义态射 $b' : Z' \to X$。
换言之，我们已把态射 $b$ 延拓到 $Z$ 的一个更大单纯子集。
% P02-E0150: frozen source says the larger simplicial subset is “of Z”, although Z' strictly contains Z and is a subset of W.

\medskip\noindent
应用 Zorn 引理即可完成证明（从略）。
\end{proof}

\begin{lemma}
\label{simplicial-lemma-trivial-kan-base-change}
设 $f : X \to Y$ 为单纯集的平凡 Kan 纤维化，并设
$Y' \to Y$ 为单纯集态射。则
$X \times_Y Y' \to Y'$ 是平凡 Kan 纤维化。
\end{lemma}

\begin{proof}
这立即由纤维积的函子性（引理 \ref{simplicial-lemma-fibre-product}）和定义得出。
\end{proof}

\begin{lemma}
\label{simplicial-lemma-trivial-kan-composition}
两个平凡 Kan 纤维化的复合仍是平凡 Kan 纤维化。
\end{lemma}

\begin{proof}
略。
\end{proof}

\begin{lemma}
\label{simplicial-lemma-limit-trivial-kan}
设 $\ldots \to U^2 \to U^1 \to U^0$ 为一列平凡 Kan 纤维化。
令 $U = \lim U^t$，其定义为逐项取 $U_n = \lim U_n^t$。
则 $U \to U^0$ 是平凡 Kan 纤维化。
\end{lemma}

\begin{proof}
略。提示：使用集合的可数满射序列之逆极限非空这一事实。
\end{proof}

\begin{lemma}
\label{simplicial-lemma-product-trivial-kan}
\begin{slogan}
平凡 Kan 纤维化的积仍是平凡 Kan 纤维化。
\end{slogan}
设 $X_i \to Y_i$ 为一族平凡 Kan 纤维化。则
$\prod X_i \to \prod Y_i$ 是平凡 Kan 纤维化。
\end{lemma}

\begin{proof}
略。
\end{proof}

\begin{lemma}
\label{simplicial-lemma-filtered-colimit-trivial-kan}
平凡 Kan 纤维化的滤过余极限仍是平凡 Kan 纤维化。
\end{lemma}

\begin{proof}
略。提示：见《范畴论》第
\ref{categories-section-directed-colimits} 节中集合的滤过余极限之描述。
\end{proof}

\begin{lemma}
\label{simplicial-lemma-trivial-kan-homotopy}
设 $f : X \to Y$ 为单纯集的平凡 Kan 纤维化。
则 $f$ 是同伦等价。
\end{lemma}

\begin{proof}
由引理 \ref{simplicial-lemma-trivial-kan}，可以选取 $f$ 的右逆
$g : Y \to X$。
% P02-E0151: frozen source says “an right inverse”.
考虑图
$$
\xymatrix{
\partial \Delta[1] \times X \ar[d] \ar[r] & X \ar[d] \\
\Delta[1] \times X \ar[r] \ar@{-->}[ru] & Y
}
$$
这里，上方水平箭头由 $\text{id}_X$ 和 $g \circ f$ 给出；
使用了对所有 $n \geq 0$ 都有
$(\partial \Delta[1] \times X)_n = X_n \amalg X_n$。
下方水平箭头由映射 $\Delta[1] \to \Delta[0]$ 和
$f : X \to Y$ 给出。由于 $f \circ g \circ f = f$，该图交换。
由引理 \ref{simplicial-lemma-trivial-kan} 可以填入虚线箭头，结论即得。
\end{proof}





\section{Kan 纤维化}
\label{simplicial-section-kan}

\noindent
设 $n$、$k$ 为满足 $0 \leq k \leq n$ 且 $1 \leq n$ 的整数。
令 $\sigma_0, \ldots, \sigma_n$ 为 $\Delta[n]$ 的唯一非退化
$n$-单形 $\sigma$ 的 $n + 1$ 个面，即
$\sigma_i = d_i\sigma$。令
$$
\Lambda_k[n] \subset \Delta[n]
$$
为 $n$-单形 $\Delta[n]$ 的 {\it 第 $k$ 个角}。
它是由 $\sigma_0, \ldots, \hat \sigma_k, \ldots, \sigma_n$
生成的 $\Delta[n]$ 的单纯子集。换言之，所显示包含的像含有
$\Delta[n]$ 的所有非退化单形，唯独不含 $\sigma$ 和 $\sigma_k$。

\begin{definition}
\label{simplicial-definition-kan}
若单纯集映射 $X \to Y$ 满足：对所有整数 $k, n$，其中
$1 \leq n$、$0 \leq k \leq n$，以及任意交换实线图
$$
\xymatrix{
\Lambda_k[n] \ar[r] \ar[d] & X \ar[d] \\
\Delta[n] \ar[r] \ar@{-->}[ru] & Y
}
$$
都存在使图交换的虚线箭头，则称它为 {\it Kan 纤维化}。
若单纯集 $X$ 到单点集上的常值单纯集 $*$ 的映射
$X \to *$ 为 Kan 纤维化，则称它为 {\it Kan 复形}。
\end{definition}

\noindent
注意，$\Lambda_k[n]$ 总是非空的。因此，从空单纯集到任意单纯集
的态射总是 Kan 纤维化。由引理 \ref{simplicial-lemma-trivial-kan} 可知，
平凡 Kan 纤维化是 Kan 纤维化。

\begin{lemma}
\label{simplicial-lemma-kan-base-change}
设 $f : X \to Y$ 为单纯集的 Kan 纤维化，并设
$Y' \to Y$ 为单纯集态射。则
$X \times_Y Y' \to Y'$ 是 Kan 纤维化。
\end{lemma}

\begin{proof}
这立即由纤维积的函子性（引理 \ref{simplicial-lemma-fibre-product}）和定义得出。
\end{proof}

\begin{lemma}
\label{simplicial-lemma-kan-composition}
两个 Kan 纤维化的复合仍是 Kan 纤维化。
\end{lemma}

\begin{proof}
略。
\end{proof}

\begin{lemma}
\label{simplicial-lemma-limit-kan}
设 $\ldots \to U^2 \to U^1 \to U^0$ 为一列 Kan 纤维化。
令 $U = \lim U^t$，其定义为逐项取 $U_n = \lim U_n^t$。
则 $U \to U^0$ 是 Kan 纤维化。
\end{lemma}

\begin{proof}
略。提示：使用集合的可数满射序列之逆极限非空这一事实。
\end{proof}

\begin{lemma}
\label{simplicial-lemma-product-kan}
设 $X_i \to Y_i$ 为一族 Kan 纤维化。则
$\prod X_i \to \prod Y_i$ 是 Kan 纤维化。
\end{lemma}

\begin{proof}
略。
\end{proof}

\noindent
下面的引理由 J.C.\ Moore 给出，见 \cite{Moore-Cartan}。

\begin{lemma}
\label{simplicial-lemma-simplicial-group-kan}
设 $X$ 为单纯群。则 $X$ 是 Kan 复形。
\end{lemma}

\begin{proof}
下面的证明基本上只是把上述参考文献中的证明译成英语。
% The frozen English source explicitly characterizes this proof as an English translation.
使用本节引言中说明的术语，设
$f : \Lambda_k[n] \to X$ 为从一个角出发的态射。
对 $i = 0, \ldots, \hat k, \ldots, n$，令
$x_i = f(\sigma_i) \in X_{n - 1}$。这意味着，当 $i < j$
且 $i, j \not = k$ 时，有 $d_i x_j = d_{j - 1} x_i$。
必须找到 $x \in X_n$，使得对
$i = 0, \ldots, \hat k, \ldots, n$，有 $x_i = d_ix$。

\medskip\noindent
先证明存在 $u \in X_n$，使得对 $i < k$ 有 $d_i u = x_i$。
当 $k = 0$ 时这是平凡的。若 $k > 0$，则归纳定义
$u^r \in X_n$，使得对 $0 \leq i \leq r$ 有
$d_i u^r = x_i$。从 $u^0 = s_0x_0$ 开始。若 $r < k - 1$，令
$$
y^r = s_{r + 1}((d_{r + 1}u^r)^{-1}x_{r + 1}),\quad
u^{r + 1} = u^r y^r.
$$
简单计算表明，当 $i \leq r$ 时，$d_iy^r = 1$
（群 $X_{n - 1}$ 的单位元），且
$d_{r + 1}y^r = (d_{r + 1}u^r)^{-1}x_{r + 1}$。
因此，当 $i \leq r + 1$ 时，$d_iu^{r + 1} = x_i$。
最后取 $u = u^{k - 1}$，得到所需的 $u$。

\medskip\noindent
接着对整数 $r$ 归纳证明：当 $0 \leq r \leq n - k$ 时，
存在 $x^r \in X_n$，使得
% P02-E0152: frozen source says “there exists a x^r”.
$$
d_i x^r = x_i\quad\text{对 }i < k\text{ 且 }i > n - r.
$$
当 $r = 0$ 时，从 $x^0 = u$ 开始。若已对
$r \leq n - k - 1$ 定义了 $x^r$，令
$$
z^r = s_{n - r - 1}((d_{n - r}x^r)^{-1}x_{n - r}),\quad x^{r + 1} = x^rz^r
$$
利用给定关系作简单计算，可知当 $i < k$ 及 $i > n - r$
时 $d_iz^r = 1$，并且
$d_{n - r}(z^r) = (d_{n - r}x^r)^{-1}x_{n - r}$。
从而当 $i < k$ 及 $i > n - r - 1$ 时，
$d_ix^{r + 1} = x_i$。最后取 $x = x^{n - k}$，证明完成。
\end{proof}

\begin{lemma}
\label{simplicial-lemma-surjection-simplicial-abelian-groups-kan}
设 $f : X \to Y$ 为逐项满射的单纯阿贝尔群同态。
则 $f$ 是单纯集的 Kan 纤维化。
\end{lemma}

\begin{proof}
考虑定义 \ref{simplicial-definition-kan} 中的交换实线图
$$
\xymatrix{
\Lambda_k[n] \ar[r]_a \ar[d] & X \ar[d] \\
\Delta[n] \ar[r]^b \ar@{-->}[ru] & Y
}
$$
映射 $a$ 对应于满足
$d_i x_j = d_{j - 1} x_i$（当 $i < j$ 且 $i, j \not = k$）
的元素
$x_0, \ldots, \hat x_k, \ldots, x_n \in X_{n - 1}$。
映射 $b$ 对应于元素 $y \in Y_n$，使得当 $i \not = k$
时 $d_iy = f(x_i)$。任务是构造 $x \in X_n$，
使得当 $i \not = k$ 时 $d_ix = x_i$，且 $f(x) = y$。

\medskip\noindent
由于 $f$ 逐项满射，可取 $x \in X_n$ 使得 $f(x) = y$。
以 $0 = y - f(x)$ 代替 $y$，并对 $i \not = k$，
以 $x_i - d_ix$ 代替 $x_i$。于是可以假设 $y = 0$。
特别地，$f(x_i) = 0$。换言之，可以把 $X$ 替换为
$\Ker(f) \subset X$，并把 $Y$ 替换为 $0$。
在这种情形下，该陈述化为引理
\ref{simplicial-lemma-simplicial-group-kan}。
% P02-E0153: frozen source says “the statement become”.
\end{proof}

\begin{lemma}
\label{simplicial-lemma-qis-simplicial-abelian-groups}
设 $f : X \to Y$ 为逐项满射的单纯阿贝尔群同态，并且在伴随链复形上
诱导拟同构。则 $f$ 是单纯集的平凡 Kan 纤维化。
\end{lemma}

\begin{proof}
考虑定义 \ref{simplicial-definition-trivial-kan} 中的交换实线图
$$
\xymatrix{
\partial \Delta[n] \ar[r]_a \ar[d] & X \ar[d] \\
\Delta[n] \ar[r]^b \ar@{-->}[ru] & Y
}
$$
映射 $a$ 对应于满足当 $i < j$ 时
$d_i x_j = d_{j - 1} x_i$ 的元素
$x_0, \ldots, x_n \in X_{n - 1}$。
映射 $b$ 对应于元素 $y \in Y_n$，使得 $d_iy = f(x_i)$。
任务是构造 $x \in X_n$，使得 $d_ix = x_i$ 且 $f(x) = y$。

\medskip\noindent
由于 $f$ 逐项满射，可取 $x \in X_n$ 使得 $f(x) = y$。
以 $0 = y - f(x)$ 代替 $y$，并以 $x_i - d_ix$ 代替 $x_i$。
于是可以假设 $y = 0$。特别地，$f(x_i) = 0$。
换言之，可以把 $X$ 替换为 $\Ker(f) \subset X$，
并把 $Y$ 替换为 $0$。这是可行的，因为由《同调代数》引理
\ref{homology-lemma-long-exact-sequence-chain}，
$\Ker(f)$ 的伴随链复形之同调为零，因此
$\Ker(f) \to 0$ 在伴随链复形上诱导拟同构。

\medskip\noindent
由于 $X$ 是 Kan 复形（引理 \ref{simplicial-lemma-simplicial-group-kan}），
可取 $x \in X_n$，使得对 $i = 0, \ldots, n - 1$ 有
$d_i x = x_i$。对 $i = 0, \ldots, n$，以
$x_i - d_ix$ 代替 $x_i$ 后，可以假设
$x_0 = x_1 = \ldots = x_{n - 1} = 0$。
此时可见，当 $i = 0, \ldots, n - 1$ 时，$d_i x_n = 0$。
因此 $x_n \in N(X)_{n - 1}$，并属于微分
$N(X)_{n - 1} \to N(X)_{n - 2}$ 的核。
这里 $N(X)$ 是与 $X$ 伴随的规范化链复形，见第
\ref{simplicial-section-complexes} 节。由于 $N(X)$ 拟同构于 $s(X)$
（引理 \ref{simplicial-lemma-quasi-isomorphism}），因而是零调的，
可取 $x \in N(X_n)$，其微分为 $x_n$。
% P02-E0154: frozen source writes N(X_n) although the normalized complex term is N(X)_n (equivalently N(X_n) is not defined by this notation).
这个 $x$ 回答了引理提出的问题，证明完成。
\end{proof}

\begin{lemma}
\label{simplicial-lemma-homotopy-equivalence}
设 $f : X \to Y$ 为单纯阿贝尔群的映射。若 $f$ 是单纯集的同伦等价，
则 $f$ 在伴随链复形上诱导拟同构。
\end{lemma}

\begin{proof}
在本证明中，当 $Z$ 是单纯阿贝尔群时，记
$H_n(Z) = H_n(s(Z)) = H_n(N(Z))$，其中 $s$ 与 $N$
如第 \ref{simplicial-section-complexes} 节所述。
以 $\mathbf{Z}[X]$ 表示把 $X$ 视为单纯集后生成的自由阿贝尔群，
$\mathbf{Z}[Y]$ 亦然。考虑单纯阿贝尔群的交换图
$$
\xymatrix{
\mathbf{Z}[X] \ar[r]_g \ar[d] &
\mathbf{Z}[Y] \ar[d] \\
X \ar[r]^f & Y
}
$$
由于在集合上取自由阿贝尔群是函子，水平箭头是单纯阿贝尔群的
同伦等价；见引理 \ref{simplicial-lemma-functorial-homotopy}。
由引理 \ref{simplicial-lemma-homotopy-equivalence-s-N}，对所有
$n \geq 0$，映射
$H_n(g) : H_n(\mathbf{Z}[X]) \to H_n(\mathbf{Z}[Y])$ 都是双射。

\medskip\noindent
设 $\xi \in H_n(Y)$。根据 $N(Y)$ 的定义，可以用边界为零的元素
$y \in N(Y_n)$ 表示 $\xi$。这意味着 $y \in Y_n$，并且
$d^n_0(y) = \ldots = d^n_{n - 1}(y) = 0$，因为
$y \in N(Y_n)$；并且 $d^n_n(y) = 0$，因为 $y$ 的边界为零。
以 $0_n \in Y_n$ 表示零元。则
$$
\tilde y = [y] - [0_n] \in (\mathbf{Z}[Y])_n
$$
满足 $d^n_0(\tilde y) = \ldots = d^n_{n - 1}(\tilde y) = 0$
以及 $d^n_n(\tilde y) = 0$。因此，$\tilde y$ 属于
$N(\mathbf{Z}[Y])_n$，其边界为 $0$；也就是说，
$\tilde y$ 决定一个映到 $\xi$ 的类
$\tilde \xi \in H_n(\mathbf{Z}[Y])$。
% P02-E0155: frozen source joins “is in ... has boundary 0” without “and”.
因为 $H_n(\mathbf{Z}[X]) \to H_n(\mathbf{Z}[Y])$ 是双射，
可以把 $\tilde \xi$ 提升为 $H_n(\mathbf{Z}[X])$ 中的一个类。
由上面的交换图可知，$\xi$ 属于 $H_n(X) \to H_n(Y)$ 的像。

\medskip\noindent
设 $\xi \in H_n(X)$ 映为 $H_n(Y)$ 中的零。
与上一段完全一样，可以用边界为零的元素
$x \in N(X_n)$ 表示 $\xi$，即
$d^n_0(x) = \ldots = d^n_{n - 1}(x) = d^n_n(x) = 0$。
% P02-E0156: frozen source misspells “paragraph” as “parapgraph”.
特别地，$[x] - [0_n]$ 是 $N(\mathbf{Z}[X])_n$ 中边界为零的元素，
从而定义 $\xi$ 的提升
$\tilde \xi \in H_n(\mathbf{Z}[x])$。
% P02-E0157: frozen source uses lowercase x in \mathbf{Z}[x] instead of the simplicial set X.
$\xi$ 映为 $H_n(Y)$ 中的零意味着存在
$y \in N(Y_{n + 1})$，其边界为 $f_n(x)$。
这意味着 $d^{n + 1}_0(y) = \ldots = d^{n + 1}_n(y) = 0$ 且
$d^{n + 1}_{n + 1}(y) = f(x)$。
然而，这恰好意味着
$z = [y] - [0_{n + 1}]$ 属于
$N(\mathbf{Z}[y])_{n + 1}$，并且
% P02-E0158: frozen source uses \mathbf{Z}[y] where the ambient free simplicial abelian group is \mathbf{Z}[Y].
$$
g([x] - [0_n]) = [f(x)] - [0_n] = \text{下列对象的边界： }z
$$
这证明 $\tilde \xi$ 映为 $H_n(\mathbf{Z}[y])$ 中的零。
% P02-E0159: frozen source again uses \mathbf{Z}[y] instead of \mathbf{Z}[Y].
因为 $H_n(\mathbf{Z}[X]) \to H_n(\mathbf{Z}[Y])$ 是双射，
可得 $\tilde \xi = 0$，从而 $\xi = 0$。
\end{proof}






\section{一个同伦等价}
\label{simplicial-section-homotopy-equivalence}

\noindent
设 $A$、$B$ 为集合，$f : A \to B$ 为映射。考虑相应的单纯集映射
$$
\xymatrix{
\text{cosk}_0(A) \ar@{=}[r] &
\Big(\ldots A \times A \times A
\ar[d] \ar@<2ex>[r] \ar@<0ex>[r] \ar@<-2ex>[r] &
A \times A \ar[d] \ar@<1ex>[r] \ar@<-1ex>[r] \ar@<1ex>[l] \ar@<-1ex>[l] &
A \Big) \ar[d] \ar@<0ex>[l] \\
\text{cosk}_0(B) \ar@{=}[r] &
\Big( \ldots B \times B \times B \ar@<2ex>[r] \ar@<0ex>[r] \ar@<-2ex>[r] &
B \times B \ar@<1ex>[r] \ar@<-1ex>[r] \ar@<1ex>[l] \ar@<-1ex>[l] &
B \Big) \ar@<0ex>[l]
}
$$
见例 \ref{simplicial-example-cosk0}。下面引理在 $n = 0$ 的情形说明：
若 $f$ 为满射，则这个单纯集映射是平凡 Kan 纤维化。

\begin{lemma}
\label{simplicial-lemma-section}
设 $f : V \to U$ 为单纯集态射，$n \geq 0$ 为整数。假设：
\begin{enumerate}
\item 当 $i < n$ 时，映射 $f_i : V_i \to U_i$ 是双射；
\item 映射 $f_n : V_n \to U_n$ 是满射；
\item 典范态射 $U \to \text{cosk}_n \text{sk}_n U$ 是同构；
\item 典范态射 $V \to \text{cosk}_n \text{sk}_n V$ 是同构。
\end{enumerate}
则 $f$ 是平凡 Kan 纤维化。
\end{lemma}

\begin{proof}
考虑定义 \ref{simplicial-definition-trivial-kan} 中的实线图
$$
\xymatrix{
\partial \Delta[k] \ar[r] \ar[d] & V \ar[d] \\
\Delta[k] \ar[r] \ar@{-->}[ru] & U
}
$$
。设 $x \in U_k$ 为与下方水平箭头对应的 $k$-单形。
若 $k \leq n$，则虚线箭头与 $x$ 的提升
$y \in V_k$ 对应；该图交换，因为 $\Delta[k]$ 的其他非退化单形
都在次数 $< k$ 中，而 $f$ 在这些次数中为同构。
若 $k > n$，则由条件 (3) 和 (4)，利用骨架函子与余骨架函子的伴随性，
有
$$
\Mor(\Delta[k], U) =
\Mor(\text{sk}_n\Delta[k], \text{sk}_nU) =
\Mor(\text{sk}_n\partial\Delta[k], \text{sk}_nU) =
\Mor(\partial \Delta[k], U)
$$
，$V$ 也同样，因为当 $k > n$ 时，
$\text{sk}_n\Delta[k] = \text{sk}_n\partial\Delta[k]$。
因此，在这种情形下也得到唯一的虚线箭头嵌入该图。
\end{proof}

\noindent
设 $A, B$ 为集合，$f^0, f^1 : A \to B$ 为集合映射。
考虑同样记号所表示的诱导映射
$f^0, f^1 : \text{cosk}_0(A) \to \text{cosk}_0(B)$。
下面引理在 $n = 0$ 的情形说明 $f^0$ 同伦于 $f^1$。
事实上，存在从 $f^0$ 到 $f^1$ 的同伦
$h : \text{cosk}_0(A) \times \Delta[1] \to \text{cosk}_0(A)$，
% P02-E0160: frozen source gives the homotopy codomain cosk_0(A), but f^0 and f^1 land in cosk_0(B).
其分量为
\begin{eqnarray*}
h_m : A \times \ldots \times A \times \Mor_{\Delta}([m], [1])
& \longrightarrow &
B \times \ldots \times B, \\
(a_0, \ldots, a_m, \alpha) & \longmapsto &
(f^{\alpha(0)}(a_0), \ldots, f^{\alpha(m)}(a_m))
\end{eqnarray*}
为了检验其成立，注意对映射 $\varphi : [k] \to [m]$，
诱导映射为
$(a_0, \ldots, a_m) \mapsto
(a_{\varphi(0)}, \ldots, a_{\varphi(k)})$
以及 $\alpha \mapsto \alpha \circ \varphi$。
因此 $h = (h_m)_{m \geq 0}$ 显然是所需的单纯集映射。

\begin{lemma}
\label{simplicial-lemma-homotopy}
设 $f^0, f^1 : V \to U$ 为单纯集映射，$n \geq 0$ 为整数。假设：
\begin{enumerate}
\item 当 $i < n$ 时，映射
$f^j_i : V_i \to U_i$（$j = 0, 1$）相等；
\item 典范态射 $U \to \text{cosk}_n \text{sk}_n U$ 是同构；
\item 典范态射 $V \to \text{cosk}_n \text{sk}_n V$ 是同构。
\end{enumerate}
则 $f^0$ 同伦于 $f^1$。
\end{lemma}

\begin{proof}[第一个证明]
设 $W$ 为 $n$-截断单纯集，其中当 $i < n$ 时
$W_i = U_i$，而 $W_n = U_n / \sim$；这里 $\sim$
是由对 $y \in V_n$ 的关系 $f^0(y) \sim f^1(y)$ 生成的等价关系。
这是有意义的，因为与 $\varphi : [i] \to [n]$（$i < n$）
对应的态射 $U(\varphi) : U_n \to U_i$ 经商映射
$U_n \to W_n$ 分解：这是由于 $f^0$ 和 $f^1$
为单纯集态射，并且在次数 $< n$ 中相等。
接着，取 $\text{cosk}_n W$，把 $W$ 提升为单纯集。
由引理 \ref{simplicial-lemma-section}，态射 $g : U \to W$
是平凡 Kan 纤维化。根据构造，
$g \circ f^0 = g \circ f^1$；以 $f : V \to W$ 表示该态射。
考虑图
$$
\xymatrix{
\partial \Delta[1] \times V \ar[rr]_{f^0, f^1} \ar[d] & & U \ar[d] \\
\Delta[1] \times V \ar[rr]^f \ar@{-->}[rru] & & W
}
$$
由引理 \ref{simplicial-lemma-trivial-kan}，虚线箭头存在，证明完成。
\end{proof}

\begin{proof}[第二个证明]
必须构造单纯集态射
$h : V \times \Delta[1] \to U$，它与 $e_i$ 复合后恢复 $f^i$。
$n = 0$ 的情形已在引理之前处理。因此可假设 $n \geq 1$。
映射
$\Delta[1] \to \text{cosk}_1 \text{sk}_1 \Delta[1]$
是同构；见引理 \ref{simplicial-lemma-simplex-cosk}。
由引理 \ref{simplicial-lemma-cosk-up}，因 $n \geq 1$，
$\Delta[1] \to \text{cosk}_n \text{sk}_n \Delta[1]$
也是同构。因此，由引理 \ref{simplicial-lemma-cosk-product}，
$V \times \Delta[1] \to
\text{cosk}_n \text{sk}_n (V \times \Delta[1])$
也是同构。换言之，为了构造同伦，只须构造合适的
$n$-截断单纯集态射
$h : \text{sk}_n V \times \text{sk}_n \Delta[1] \to \text{sk}_n U$。

\medskip\noindent
对 $k = 0, \ldots, n - 1$，以公式
$h_k(v, \alpha) = f^0(v) = f^1(v)$ 定义 $h_k$。
映射 $h_n : V_n \times \Mor_{\Delta}([n], [1]) \to U_n$
如下定义。取 $v \in V_n$ 与 $\alpha : [n] \to [1]$：
\begin{enumerate}
\item 若 $\Im(\alpha) = \{0\}$，令 $h_n(v, \alpha) = f^0(v)$；
\item 若 $\Im(\alpha) = \{0, 1\}$，令 $h_n(v, \alpha) = f^0(v)$；
\item 若 $\Im(\alpha) = \{1\}$，令 $h_n(v, \alpha) = f^1(v)$。
\end{enumerate}
设 $\varphi : [k] \to [l]$ 为 $\Delta_{\leq n}$ 的态射。
我们将证明图
$$
\xymatrix{
V_{l} \times \Mor([l], [1]) \ar[r] \ar[d] & U_{l} \ar[d] \\
V_{k} \times \Mor([k], [1]) \ar[r] & U_{k}
}
$$
交换。取 $v \in V_{l}$ 和 $\alpha : [l] \to [1]$。
交换性意味着
$$
h_k(V(\varphi)(v), \alpha \circ \varphi)
=
U(\varphi)(h_l(v, \alpha)).
$$
几乎在所有情形下，这都由
$h_k(V(\varphi)(v), \alpha \circ \varphi) =
f^0(V(\varphi)(v))$ 和
$U(\varphi)(h_l(v, \alpha)) = U(\varphi)(f^0(v))$，
结合 $f^0$ 为单纯集态射这一事实得出。
唯一不满足这一描述的情形是：
或者 (A) $\Im(\alpha) = \{1\}$ 且 $l = n$，
或者 (B) $\Im(\alpha \circ \varphi) = \{1\}$ 且 $k = n$。
还注意，对 $V$ 的任意退化 $n$-单形，必有
$f^0(v) = f^1(v)$。因此，可把上述情形进一步缩小为：
(A) $\Im(\alpha) = \{1\}$、$l = n$ 且 $v$ 非退化；
以及 (B) $\Im(\alpha \circ \varphi) = \{1\}$、
$k = n$ 且 $V(\varphi)(v)$ 非退化。

\medskip\noindent
在情形 (A) 中，也有
$\Im(\alpha \circ \varphi) = \{1\}$。
因此，不仅 $h_l(v, \alpha) = f^1(v)$，而且
$h_k(V(\varphi)(v), \alpha \circ \varphi) =
f^1(V(\varphi)(v))$。由于 $f^1$ 是单纯集态射，该关系成立。

\medskip\noindent
在情形 (B) 中，断定 $l = k = n$ 且 $\varphi$ 为双射，
否则 $V(\varphi)(v)$ 是退化的。因此
$\varphi = \text{id}_{[n]}$，这是平凡情形。
\end{proof}

\begin{lemma}
\label{simplicial-lemma-cosk-minus-one-equivalence}
设 $A$、$B$ 为集合，并设 $f : A \to B$ 为映射。
% P02-E0161: frozen source says “Let A, B be sets, and that f ...”.
考虑第 $n$ 个单形为
$$
A \times_B A \times_B \ldots \times_B A\ (n + 1 \text{ 可分解）}.
$$
的单纯集 $U$；见例 \ref{simplicial-example-fibre-products-simplicial-object}。
若 $f$ 满射，则态射 $U \to B$ 是平凡 Kan 纤维化；
这里 $B$ 表示取值为 $B$ 的常值单纯集。
\end{lemma}

\begin{proof}
注意，$U$ 嵌入笛卡尔方块
$$
\xymatrix{
U \ar[d] \ar[r] & \text{cosk}_0(A) \ar[d] \\
B \ar[r] & \text{cosk}_0(B)
}
$$
由引理 \ref{simplicial-lemma-section}，右侧竖直箭头是平凡 Kan 纤维化；
再由引理 \ref{simplicial-lemma-trivial-kan-base-change}，左侧也如此。
\end{proof}










\section{标准预解式的准备}
\label{simplicial-section-standard-pre}

\noindent
本节材料可见 \cite[Appendix 1]{Godement}
% P02-E0162: frozen introductory sentence lacks terminal punctuation.

\begin{example}
\label{simplicial-example-godement}
设 $Y : \mathcal{C} \to \mathcal{C}$ 为一个范畴到自身的函子，并设给定
函子变换
$$
d : Y \longrightarrow \text{id}_\mathcal{C}
\quad\text{且}\quad
s : Y \longrightarrow Y \circ Y
$$
利用这些变换，我们可以构造一个形似单纯对象的东西。具体地，对
$n \geq 0$，定义
$$
X_n = Y \circ \ldots \circ Y \quad (n + 1\text{ 复合})
$$
注意，对 $n, m \geq 0$，有 $X_{n + m + 1} = X_n \circ X_m$。
接着，对 $n \geq 0$ 和 $0 \leq j \leq n$，采用《范畴论》第
\ref{categories-section-formal-cat-cat} 节中的记号，定义
$$
d^n_j = 1_{X_{j - 1}} \star d \star 1_{X_{n - j - 1}} :
X_n
\to
X_{n - 1}
\quad\text{且}\quad
s^n_j = 1_{X_{j - 1}} \star s \star 1_{X_{n - j - 1}} :
X_n
\to
X_{n + 1}
$$
因此，$d^n_j$（相应地，$s^n_j$）是在定义 $X_n$ 的复合中，
对从左数第 $j$ 个 $Y$ 使用 $d$（相应地，$s$）所得的自然变换。
\end{example}

\begin{lemma}
\label{simplicial-lemma-godement}
在例 \ref{simplicial-example-godement} 中，若
$$
1_Y = (d \star 1_Y) \circ s = (1_Y \star d) \circ s
\quad\text{且}\quad
(s \star 1) \circ s = (1 \star s) \circ s
$$
则 $X = (X_n, d^n_j, s^n_j)$ 是 $\mathcal{C}$ 的自函子范畴中的
单纯对象，且 $d : X_0 = Y \to \text{id}_\mathcal{C}$ 定义一个增广。
\end{lemma}

\begin{proof}
为证明所得对象是单纯对象，须验证引理
\ref{simplicial-lemma-characterize-simplicial-object} 的关系 (1)(a)--(e)。
我们将使用简写
$$
1_a = 1_{X_{a - 1}} = 1_Y \star \ldots \star 1_Y \quad (a\text{ 可分解})
$$
，其中 $a \geq 0$。采用此记号，有
$$
d^n_j = 1_j \star d \star 1_{n - j}
\quad\text{且}\quad
s^n_j = 1_j \star s \star 1_{n - j}
$$
% P02-E0163: frozen source misspells “functors” as “funtors”.
下文反复使用如下规则：对函子变换 $a, a', b, b'$，只要公式中的
$\star$ 与 $\circ$ 复合有定义，就有
$(a' \circ a) \star (b' \circ b) = (a' \star b') \circ (a \star b)$；
见《范畴论》引理 \ref{categories-lemma-properties-2-cat-cats}。

\medskip\noindent
条件 (1)(a) 总是成立（不需要对 $d$ 和 $s$ 作任何假设）。
事实上，设 $0 \leq i < j \leq n + 1$。我们须证明
$d^n_i \circ d^{n + 1}_j = d^n_{j - 1} \circ d^{n + 1}_i$，即
$$
(1_i \star d \star 1_{n - i}) \circ
(1_j \star d \star 1_{n + 1 - j}) =
(1_{j - 1} \star d \star 1_{n + 1 - j}) \circ
(1_i \star d \star 1_{n + 1 - i})
$$
左端可改写为
\begin{align*}
& (1_i \star d \star 1_{j - i - 1} \star 1_{n + 1 - j}) \circ
(1_i \star 1_1 \star 1_{j - i - 1} \star d \star 1_{n + 1 - j}) \\
& =
1_i \star
\left((d \star 1_{j - i - 1}) \circ
(1_1 \star 1_{j - i - 1} \star d)\right) \star 1_{n + 1 - j} \\
& =
1_i \star d \star 1_{j - i - 1} \star d \star 1_{n + 1 - j}
\end{align*}
第二个等号成立，是因为 $d \circ 1_1 = d$，且
$1_{j - i} \circ (1_{j - i - 1} \star d) = 1_{j - i - 1} \star d$。
对右端作相同计算也得到同一结果。

\medskip\noindent
验证条件 (1)(b)。设 $0 \leq i < j \leq n - 1$。须证明
$d^n_i \circ s^{n - 1}_j = s^{n - 2}_{j - 1} \circ d^{n - 1}_i$，即
$$
(1_i \star d \star 1_{n - i}) \circ
(1_j \star s \star 1_{n - 1 - j}) =
(1_{j - 1} \star s \star 1_{n - 1 - j}) \circ
(1_i \star d \star 1_{n - 1 - i})
$$
通过与情形 (1)(a) 相同的计算，两端都化为
$1_i \star d \star 1_{j - i - 1} \star s \star 1_{n - j - 1}$。

\medskip\noindent
验证条件 (1)(c)。设 $0 \leq j \leq n - 1$。须证明
$\text{id} = d^n_j \circ s^{n - 1}_j = d^n_{j + 1} \circ s^{n - 1}_j$，即
$$
1_n =
(1_j \star d \star 1_{n - j}) \circ
(1_j \star s \star 1_{n - 1 - j}) =
(1_{j + 1} \star d \star 1_{n - j - 1}) \circ
(1_j \star s \star 1_{n - 1 - j})
$$
容易看出，这由引理的第一个假设推出。

\medskip\noindent
验证条件 (1)(d)。设 $0 < j + 1 < i \leq n$。须证明
$d^n_i \circ s^{n - 1}_j = s^{n - 2}_j \circ d^{n - 1}_{i - 1}$，即
$$
(1_i \star d \star 1_{n - i}) \circ
(1_j \star s \star 1_{n - 1 - j}) =
(1_j \star s \star 1_{n - 2 - j}) \circ
(1_{i - 1} \star d \star 1_{n - i})
$$
通过与情形 (1)(a) 相同的计算，两端都化为
$1_j \star s \star 1_{i - j - 2} \star d \star 1_{n - i}$。

\medskip\noindent
验证条件 (1)(e)。设 $0 \leq i \leq j \leq n - 1$。须证明
$s^n_i \circ s^{n - 1}_j = s^n_{j + 1} \circ s^{n - 1}_i$，即
$$
(1_i \star s \star 1_{n - i}) \circ
(1_j \star s \star 1_{n - 1 - j}) =
(1_{j + 1} \star s \star 1_{n - 1 - j}) \circ
(1_i \star s \star 1_{n - 1 - i})
$$
通过与情形 (1)(a) 相同的计算，这化为
$$
(s \star 1_{j - i + 1}) \circ (1_{j - i} \star s) =
(1_{j - i + 1} \star s) \circ (s \star 1_{j - i})
$$
若 $j = i$，这恰为引理的两个假设之一。若 $j > i$，通过与
情形 (1)(a) 相似的计算，左端和右端都化为等式
$s \star 1_{j - i - 1} \star s$。

\medskip\noindent
最后，为证明 $d$ 定义一个增广，须证明
$d \circ (1_1 \star d) = d \circ (d \star 1_1)$；这是成立的，
因为两端都等于 $d \star d$。
\end{proof}

\begin{example}
\label{simplicial-example-godement-functorial}
设 $\mathcal{C}$、$Y$、$d$、$s$ 如例 \ref{simplicial-example-godement}，
并满足引理 \ref{simplicial-lemma-godement} 的等式。给定函子
$F : \mathcal{A} \to \mathcal{C}$ 和
$G : \mathcal{C} \to \mathcal{B}$，我们在从 $\mathcal{A}$
到 $\mathcal{B}$ 的函子范畴中得到单纯对象 $G \circ X \circ F$，
它带有到 $G \circ F$ 的增广。
\end{example}

\begin{lemma}
\label{simplicial-lemma-godement-section}
设 $\mathcal{A}$、$\mathcal{B}$、$\mathcal{C}$、$Y$、$d$、$s$、$F$、
$G$ 如例 \ref{simplicial-example-godement-functorial}。给定函子变换
$h_0 : G \circ F \to G \circ Y \circ F$，使得
$$
1_{G \circ F} = (1_G \star d \star 1_F) \circ h_0
$$
则存在单纯对象态射 $h : G \circ F \to G \circ X \circ F$，
满足 $\epsilon \circ h = \text{id}$；这里
$\epsilon : G \circ X \circ F \to G \circ F$ 是增广。
\end{lemma}

\begin{proof}
以 $u_n : Y = X_0 \to X_n$ 表示单纯对象 $X$ 中与
$\Delta$ 的唯一态射 $[n] \to [0]$ 对应的映射。令
$h_n : G \circ F \to G \circ X_n \circ F$ 等于
$(1_G \star u_n \star 1_F) \circ h_0$。

\medskip\noindent
对任意范畴中的任意单纯对象 $X = (X_n)$，
$u =(u_n) : X_0 \to X$ 是从取值为 $X_0$ 的常值单纯对象到 $X$
的态射。因此，$h$ 是单纯对象态射，因为它是
$1_G \star u \star 1_F$ 与 $h_0$ 的复合。

\medskip\noindent
验证 $\epsilon \circ h = \text{id}$。计算得
$$
\epsilon_n \circ (1_G \star u_n \star 1_F) \circ h_0 =
\epsilon_0 \circ h_0 = \text{id}
$$
第一个等号成立是因为 $\epsilon$ 是单纯对象态射；第二个等号成立是因为
$\epsilon_0 = (1_G \star d \star 1_F)$，故可应用引理陈述中的假设。
\end{proof}

\begin{lemma}
\label{simplicial-lemma-godement-two-maps}
设 $\mathcal{A}$、$\mathcal{B}$、$\mathcal{C}$、$Y$、$d$、$s$、$F$、
$G$ 如例 \ref{simplicial-example-godement-functorial}。设
$F' : \mathcal{A} \to \mathcal{C}$ 和
$G' : \mathcal{C} \to \mathcal{B}$ 为两个函子。
设 $(a_n) : G \circ X \to G' \circ X$ 为单纯对象态射，
且经增广与 $a : G \to G'$ 相容。设
$(b_n) : X \circ F \to X \circ F'$ 为单纯对象态射，
且经增广与 $b : F \to F'$ 相容。则两个映射
$$
a \star (b_n), (a_n) \star b : G \circ X \circ F \to G' \circ X \circ F'
$$
同伦。
\end{lemma}

\begin{proof}
为证明这些态射同伦，我们构造态射
$$
h_{n, i} : G \circ X_n \circ F \to G' \circ X_n \circ F'
$$
，其中 $n \geq 0$ 且 $0 \leq i \leq n + 1$，并使其满足引理
\ref{simplicial-lemma-relations-homotopy} 中所述关系；另见注
\ref{simplicial-remark-homotopy-better}。为满足引理
\ref{simplicial-lemma-relations-homotopy} 的条件 (1)，必须令
$h_{n, 0} = a \star b_n$ 和 $h_{n , n + 1} = a_n \star b$。
因此合理的选择是
$$
h_{n , i} = a_{i - 1} \star b_{n - i}
$$
，其中 $1 \leq i \leq n$。令 $a = a_{-1}$、$b = b_{-1}$，可见
% P02-E0164: frozen source misspells “formula” as “formular”.
上述公式对 $0 \leq i \leq n + 1$ 都成立。

\medskip\noindent
回顾
$$
d^n_j = 1_G \star 1_j \star d \star 1_{n - j} \star 1_F
$$
在 $G \circ X \circ F$ 上成立；这里采用引理
\ref{simplicial-lemma-godement} 证明中引入的记号
$1_a = 1_{Y \circ \ldots \circ Y}$。下文将使用其改写
\begin{align*}
d^n_j
& =
d^j_j \star 1_{n - j} =
d^{j + 1}_j \star 1_{n - j} = \ldots = d^{n - 1}_j \star 1_1 \\
& =
1_j \star d^{n - j}_0 = 1_{j - 1} \star d^{n - j + 1}_1 = \ldots =
1_1 \star d^{n - 1}_{j - 1}
\end{align*}
当然，在 $G' \circ X \circ F'$ 上的 $d^n_j$ 也有类似公式。

\medskip\noindent
验证引理 \ref{simplicial-lemma-relations-homotopy} 的条件 (2)。
设 $i > j$。须证明
$$
d^n_j \circ (a_{i - 1} \star b_{n - i}) =
(a_{i - 2} \star b_{n - i}) \circ d^n_j
$$
因 $i - 1 \geq j$，可用 $d^n_j$ 的一种描述把左端改写为
$$
(d^{i - 1}_j \star 1_{n - i + 1}) \circ (a_{i - 1} \star b_{n - i}) =
(d^{i - 1}_j \circ a_{i - 1}) \star b_{n - i} =
(a_{i - 2} \circ d^{i - 1}_j) \star b_{n - i}
$$
同样，右端变为
$$
(a_{i - 2} \star b_{n - i}) \circ (d^{i - 1}_j \star 1_{n - i + 1}) =
(a_{i - 2} \circ d^{i - 1}_j) \star b_{n - i}
$$
故两端结果相同，条件 (2) 得证。

\medskip\noindent
验证引理 \ref{simplicial-lemma-relations-homotopy} 的条件 (3)。
设 $i \leq j$。须证明
$$
d^n_j \circ (a_{i - 1} \star b_{n - i}) =
(a_{i - 1} \star b_{n - 1 - i}) \circ d^n_j
$$
因 $j \geq i$，可把左端改写为
$$
(1_i \star d^{n - i}_{j - i}) \circ (a_{i - 1} \star b_{n - i}) =
a_{i - 1} \star (b_{n - 1 - i} \circ d^{n - i}_{j - i})
$$
相同的变形表明它与右端相等。

\medskip\noindent
回顾
$$
s^n_j = 1_G \star 1_j \star s \star 1_{n - j} \star 1_F
$$
在 $G \circ X \circ F$ 上成立。下文将使用其改写
\begin{align*}
s^n_j
& =
s^j_j \star 1_{n - j} = s^{j + 1}_j \star 1_{n - j - 1} = \ldots =
s^{n - 1}_j \star 1_1 \\
& =
1_j \star s^{n - j}_0 = 1_{j - 1} \star s^{n - j + 1}_1 = \ldots =
1_1 \star s^{n - 1}_{j - 1}
\end{align*}
当然，在 $G' \circ X \circ F'$ 上的 $s^n_j$ 也有类似公式。

\medskip\noindent
验证引理 \ref{simplicial-lemma-relations-homotopy} 的条件 (4)。
设 $i > j$。须证明
$$
s^n_j \circ (a_{i - 1} \star b_{n - i}) =
(a_i \star b_{n - i}) \circ s^n_j
$$
因 $i - 1 \geq j$，可把左端改写为
$$
(s^{i - 1}_j \star 1_{n - i + 1}) \circ (a_{i - 1} \star b_{n - i}) =
(s^{i - 1}_j \circ a_{i - 1}) \star b_{n - i} =
(a_i \circ s^{i - 1}_j) \star b_{n - i}
$$
同样，右端变为
$$
(a_i \star b_{n - i}) \circ (s^{i - 1}_j \star 1_{n - i + 1}) =
(a_i \circ s^{i - 1}_j) \star b_{n - i}
$$
，正是所需等式。

\medskip\noindent
验证引理 \ref{simplicial-lemma-relations-homotopy} 的条件 (5)。
设 $i \leq j$。须证明
$$
s^n_j \circ (a_{i - 1} \star b_{n - i}) =
(a_{i - 1} \star b_{n + 1 - i}) \circ s^n_j
$$
与上面的论证完全相同，两端都化为
$a_{i - 1} \star (s^{n - i}_{j - i} \circ b_{n - i}) =
a_{i - 1} \star (b_{n + 1 - i} \circ s^{n - i}_{j - i})$，
故等式成立。
\end{proof}

\begin{lemma}
\label{simplicial-lemma-godement-before-after}
设 $\mathcal{C}$、$Y$、$d$、$s$ 如例 \ref{simplicial-example-godement}，
并满足引理 \ref{simplicial-lemma-godement} 的等式。设
$f : \text{id}_\mathcal{C} \to \text{id}_\mathcal{C}$ 为恒等函子的
自同态。则 $f \star 1_X, 1_X \star f : X \to X$ 是单纯对象映射，
且经增广 $\epsilon : X \to \text{id}_\mathcal{C}$ 与 $f$ 相容。
此外，$f \star 1_X$ 与 $1_X \star f$ 同伦。
\end{lemma}

\begin{proof}
映射 $f \star 1_X$ 的分量为
$$
X_n = \text{id}_\mathcal{C} \circ X_n
\xrightarrow{f \star 1_{X_n}}
\text{id}_\mathcal{C} \circ X_n = X_n
$$
。对 $\mathcal{C}$ 的自函子之间的变换 $a : F \to G$，有
$a \circ (f \star 1_F) = f \star a = (f \star 1_G) \circ a$。
因此 $f \star 1_X$ 确为单纯对象态射；$1_X \star f$ 的情形相同。

\medskip\noindent
为证明这些态射同伦，我们对 $n \geq 0$ 和
$0 \leq i \leq n + 1$ 构造态射 $h_{n, i} : X_n \to X_n$，
使其满足引理 \ref{simplicial-lemma-relations-homotopy} 中所述关系；另见注
\ref{simplicial-remark-homotopy-better}。可以取
$$
h_{n, i} = 1_i \star f \star 1_{n + 1 - i}
$$
，其中 $1_i$ 是引理 \ref{simplicial-lemma-godement} 证明中
$Y \circ \ldots \circ Y$ 上的恒等变换。我们有
$h_{n, 0} = f \star 1_{X_n}$ 和 $h_{n, n + 1} = 1_{X_n} \star f$，
这验证了第一个条件。验证其余条件时，我们使用引理
\ref{simplicial-lemma-godement-two-maps} 证明中关于映射 $d^n_j$ 与 $s^n_j$ 的说明。

\medskip\noindent
验证引理 \ref{simplicial-lemma-relations-homotopy} 的条件 (2)。
设 $i > j$。须证明
$$
d^n_j \circ (1_i \star f \star 1_{n + 1 - i}) =
(1_{i - 1} \star f \star 1_{n + 1 - i}) \circ d^n_j
$$
因 $i - 1 \geq j$，可用 $d^n_j$ 的一种描述把左端改写为
$$
(d^{i - 1}_j \star 1_{n - i + 1}) \circ (1_i \star f \star 1_{n + 1 - i}) =
d^{i - 1}_j \star f \star 1_{n + 1 - i}
$$
同样，右端变为
$$
(1_{i - 1} \star f \star 1_{n + 1 - i}) \circ
(d^{i - 1}_j \star 1_{n - i + 1}) =
d^{i - 1}_j \star f \star 1_{n + 1 - i}
$$
故两端结果相同，条件 (2) 得证。

\medskip\noindent
引理 \ref{simplicial-lemma-relations-homotopy} 的条件 (3)、(4)、(5)
可完全同样地验证，采用引理 \ref{simplicial-lemma-godement-two-maps}
证明中的策略即可。细节从略\footnote{当
$f$ 可逆时，只须证明 $(a_n) = 1_X$ 与
$(b_n) = f^{-1} \star 1_X \star f$ 同伦。
但这由引理 \ref{simplicial-lemma-godement-two-maps} 得出，
因为此时 $a = b = 1_{\text{id}_\mathcal{C}}$。}。
\end{proof}






\section{标准预解式}
\label{simplicial-section-standard}

\noindent
本节部分材料可见 \cite[Appendix 1]{Godement} 和
\cite[I 1.5]{cotangent}。

\begin{situation}
\label{simplicial-situation-adjoint-functors}
设 $\mathcal{A}$、$\mathcal{S}$ 为范畴，且
$V : \mathcal{A} \to \mathcal{S}$ 是具有左伴随
$U : \mathcal{S} \to \mathcal{A}$ 的函子。
\end{situation}

\noindent
在这一极一般的情形下，我们将在从 $\mathcal{A}$ 到
$\mathcal{A}$ 的函子范畴中构造一个单纯对象 $X$。
建议在阅读本节正文之前，先看后面给出的例子。

\medskip\noindent
构造中将使用《范畴论》第 \ref{categories-section-formal-cat-cat}
节所定义的水平复合。《范畴论》第
\ref{categories-section-adjoint} 节中伴随态射\footnote{不能以
$\epsilon$ 表示伴随的余单位，因为我们要用
$\epsilon$ 表示单纯对象的增广。}的定义为
$$
d : U \circ V \to \text{id}_\mathcal{A} \quad (\text{余单位})
\quad\text{且}\quad
\eta : \text{id}_\mathcal{S} \to V \circ U \quad (\text{单位})
$$
，并说明复合
\begin{equation}
\label{simplicial-equation-composition}
V \xrightarrow{\eta \star 1_V} V \circ U \circ V \xrightarrow{1_V \star d} V
\quad\text{且}\quad
U \xrightarrow{1_U \star \eta} U \circ V \circ U \xrightarrow{d \star 1_U} U
\end{equation}
都是恒等态射。这里，为定义态射 $\eta \star 1_V$，我们默认把
$V$ 与 $\text{id}_\mathcal{S} \circ V$ 识别，而
$1_V$ 表示 $\text{id}_V : V \to V$。下文将反复使用这些记号和关系。
对 $n \geq 0$，令
$$
X_n = (U \circ V)^{\circ (n + 1)} =
U \circ V \circ U \circ \ldots \circ U \circ V
$$
。换言之，$X_n$ 是 $U \circ V$ 与自身的 $(n + 1)$ 重复合。
还令 $X_{-1} = \text{id}_\mathcal{A}$。对所有 $n, m \geq -1$，
有 $X_{n + m + 1} = X_n \circ X_m$。我们将通过构造函子态射
$$
d^n_j : X_n \to X_{n - 1},\quad s^n_j : X_n \to X_{n + 1}
$$
，使其满足引理 \ref{simplicial-lemma-relations-face-degeneracy} 中所列关系，
从而赋予这个函子序列以 $\text{Fun}(\mathcal{A}, \mathcal{A})$
中单纯对象的结构。具体地，令
$$
d^n_j = 1_{X_{j - 1}} \star d \star 1_{X_{n - j - 1}}
\quad\text{且}\quad
s^n_j = 1_{X_{j - 1} \circ U} \star \eta \star 1_{V \circ X_{n - j - 1}}
$$
最后，记 $\epsilon_0 = d : X_0 \to X_{-1}$。

\begin{lemma}
\label{simplicial-lemma-standard-simplicial}
在情形 \ref{simplicial-situation-adjoint-functors} 中，
系统 $X = (X_n, d^n_j, s^n_j)$ 是
$\text{Fun}(\mathcal{A}, \mathcal{A})$ 中的单纯对象，
且 $\epsilon_0$ 定义从 $X$ 到取值为
$X_{-1} = \text{id}_\mathcal{A}$ 的常值单纯对象的增广 $\epsilon$。
\end{lemma}

\begin{proof}
考虑 $Y = U \circ V : \mathcal{A} \to \mathcal{A}$。
我们已有变换 $d : Y = U \circ V \to \text{id}_\mathcal{A}$。记
$$
s = 1_U \star \eta \star 1_V :
Y = U \circ \text{id}_\mathcal{S} \circ V
\longrightarrow
U \circ V \circ U \circ V = Y \circ Y
$$
这使我们处于例 \ref{simplicial-example-godement} 的情形。由公式立即可见，
上面构造的 $X, d^n_i, s^n_i$ 与从 $Y, d, s$ 按例
\ref{simplicial-example-godement} 构造的
% P02-E0165: frozen source repeats “X, s^n_i, s^n_i” where the first s-family should be d.
$X, s^n_i, s^n_i$ 一致。因此，由引理 \ref{simplicial-lemma-godement}，
只须证明
$$
1_Y = (d \star 1_Y) \circ s = (1_Y \star d) \circ s
\quad\text{且}\quad
(s \star 1) \circ s = (1 \star s) \circ s
$$
第一个等号转化为
$$
1_U \star 1_V = (d \star 1_U \star 1_V) \circ (1_U \star \eta \star 1_V)
$$
；只要有 $1_U = (d \star 1_U) \circ (1_U \star \eta)$，它便成立，
而后者由 (\ref{simplicial-equation-composition}) 成立。第二个等号同理。
对最后一个等式，须证明
$$
(1_U \star \eta \star 1_V \star 1_U \star 1_V) \circ
(1_U \star \eta \star 1_V) =
(1_U \star 1_V \star 1_U \star \eta \star 1_V) \circ
(1_U \star \eta \star 1_V)
$$
为此，只须证明
$(\eta \star 1_V \star 1_U) \circ \eta = (1_V \star 1_U \star \eta) \circ \eta$；
这是成立的，因为两端都与 $\eta \star \eta$ 相同。
\end{proof}

\noindent
阅读下一个引理的证明之前，建议先看例
\ref{simplicial-example-polynomial-algebra-homotopy} 中讨论的例子，以理解该引理的目的。

\begin{lemma}
\label{simplicial-lemma-standard-simplicial-homotopy}
在情形 \ref{simplicial-situation-adjoint-functors} 中，映射
$$
1_V \star \epsilon : V \circ X \to V,
\quad\text{且}\quad
\epsilon \star 1_U : X \circ U \to U
$$
都是同伦等价。
\end{lemma}

\begin{proof}
如引理 \ref{simplicial-lemma-standard-simplicial} 的证明，令
$Y = U \circ V$，于是处于例 \ref{simplicial-example-godement} 的情形。

\medskip\noindent
证明第一个同伦等价。由引理 \ref{simplicial-lemma-godement-section}，
为构造映射 $h : V \to V \circ X$，使它成为
$1_V \star \epsilon$ 的右逆，只须构造映射
$h_0 : V \to V \circ Y = V \circ U \circ V$，使
$1_V = (1_V \star d) \circ h_0$。当然，取
$h_0 = \eta \star 1_V$，等式由 (\ref{simplicial-equation-composition}) 成立。
为完成证明，须说明两个映射
$$
(1_V \star \epsilon) \circ h, 1_V \star \text{id}_X :
V \circ X \longrightarrow V \circ X
$$
同伦。这立即由引理 \ref{simplicial-lemma-godement-two-maps} 得出（取
$G = G' = V$ 且 $F = F' = \text{id}_\mathcal{S}$）。

\medskip\noindent
证明第二个同伦等价。由引理 \ref{simplicial-lemma-godement-section}，
为构造映射 $h : U \to X \circ U$，使它成为
$\epsilon \star 1_U$ 的右逆，只须构造映射
$h_0 : U \to Y \circ U = U \circ V \circ U$，使
$1_U = (d \star 1_U) \circ h_0$。当然，取
$h_0 = 1_U \star \eta$，等式由 (\ref{simplicial-equation-composition}) 成立。
为完成证明，须说明两个映射
$$
(\epsilon \star 1_U) \circ h, \text{id}_X \star 1_U :
X \circ U \longrightarrow X \circ U
$$
同伦。这立即由引理 \ref{simplicial-lemma-godement-two-maps} 得出（取
$G = G' = \text{id}_\mathcal{A}$ 且 $F = F' = U$）。
\end{proof}


\begin{example}
\label{simplicial-example-module}
设 $R$ 为环。作为上述构造的一个例子，可取遗忘函子
$i : \text{Mod}_R \to \textit{Sets}$，以及把集合 $E$ 送到
由 $E$ 生成的自由 $R$-模 $R[E]$ 的函子
$F : \textit{Sets} \to \text{Mod}_R$。对 $R$-模 $M$，
单纯 $R$-模 $X(M)$ 具有如下形状：
$$
X(M) =
\left( \ldots
\xymatrix{
R[R[R[M]]]
\ar@<2ex>[r]
\ar@<0ex>[r]
\ar@<-2ex>[r]
&
R[R[M]]
\ar@<1ex>[r]
\ar@<-1ex>[r]
\ar@<1ex>[l]
\ar@<-1ex>[l]
&
R[M]
\ar@<0ex>[l]
}
\right)
$$
，并带有到 $M$ 的增广。我们还将证明该增广在集合层面为同伦等价。
由引理 \ref{simplicial-lemma-trivial-kan-homotopy}、
\ref{simplicial-lemma-homotopy-equivalence} 和
\ref{simplicial-lemma-qis-simplicial-abelian-groups}，这等价于要求
% P02-E0166: frozen source says “cohomology group” for the homologically indexed chain complex here and again in Example polynomial-algebra.
$M$ 是与单纯模 $X(M)$ 伴随的链复形唯一的非零上同调群。
\end{example}

\begin{example}
\label{simplicial-example-polynomial-algebra}
设 $A$ 为环，$\textit{Alg}_A$ 为交换 $A$-代数的范畴。
作为上述构造的一个例子，可取遗忘函子
$i : \textit{Alg}_A \to \textit{Sets}$，以及把集合 $E$ 送到
$A$ 上由 $E$ 生成的多项式代数 $A[E]$ 的函子
$F : \textit{Sets} \to \textit{Alg}_A$。
（对于本例与例 \ref{simplicial-example-module} 之间的记号重叠，我们深表歉意。）
对 $A$-代数 $B$，单纯 $A$-代数 $X(B)$ 具有如下形状：
$$
X(B) = \left( \ldots
\xymatrix{
A[A[A[B]]]
\ar@<2ex>[r]
\ar@<0ex>[r]
\ar@<-2ex>[r]
&
A[A[B]]
\ar@<1ex>[r]
\ar@<-1ex>[r]
\ar@<1ex>[l]
\ar@<-1ex>[l]
&
A[B]
\ar@<0ex>[l]
}
\right)
$$
，并带有到 $B$ 的增广。我们还将证明该增广在集合层面为同伦等价。
由引理 \ref{simplicial-lemma-trivial-kan-homotopy}、
\ref{simplicial-lemma-homotopy-equivalence} 和
\ref{simplicial-lemma-qis-simplicial-abelian-groups}，这等价于要求
$B$ 是把 $X(B)$ 视为单纯 $A$-模后与之伴随的
$A$-模链复形唯一的非零上同调群。
\end{example}

\begin{example}
\label{simplicial-example-module-maps}
在例 \ref{simplicial-example-module} 中，
$X_n(M) = R[R[\ldots [M]\ldots]]$ 有 $n + 1$ 层方括号。
我们用 $X_1(M)$ 的一个典型元素
$$
\xi = \sum\nolimits_i r_i\left[\sum\nolimits_j r_{ij}[m_{ij}]\right]
$$
描述上面构造的映射。映射
$d_0, d_1 : R[R[M]] \to R[M]$ 由
$$
d_0(\xi) = \sum\nolimits_{i, j} r_ir_{ij}[m_{ij}]
\quad\text{且}\quad
d_1(\xi) = \sum\nolimits_i r_i\left[\sum\nolimits_j r_{ij}m_{ij}\right].
$$
给出。映射 $s_0, s_1 : R[R[M]] \to R[R[R[M]]]$ 由
$$
s_0(\xi) =
\sum\nolimits_i r_i\left[\left[\sum\nolimits_j r_{ij}[m_{ij}]\right]\right]
\quad\text{且}\quad
s_1(\xi) = \sum\nolimits_i r_i\left[\sum\nolimits_j r_{ij}[[m_{ij}]]\right].
$$
给出。
\end{example}

\begin{example}
\label{simplicial-example-polynomial-algebra-maps}
在例 \ref{simplicial-example-polynomial-algebra} 中，
$X_n(B) = A[A[\ldots [B]\ldots]]$ 有 $n + 1$ 层方括号。
我们用一个典型元素
$$
\xi = \sum\nolimits_i a_i [x_{i, 1}] \ldots [x_{i, m_i}] \in A[A[B]] = X_1(B)
$$
描述上面构造的映射，其中对每个 $i, j$，可写
$$
x_{i, j} = \sum a_{i, j, k} [b_{i, j, k, 1}] \ldots [b_{i, j, k, n_{i, j, k}}]
\in A[B]
$$
显然，这个记号极其可怕！为简化记号，也为了看清 $A$-代数映射
$d_0, d_1 : A[A[B]] \to A[B]$ 的作用，只须考察变量 $[x]$；
这里
$$
x = \sum a_k [b_{k, 1}] \ldots [b_{k, n_k}] \in A[B]
$$
是一般元素。对这些变量，有
$$
d_0([x]) = x = \sum a_k [b_{k, 1}] \ldots [b_{k, n_k}]
\quad\text{且}\quad
d_1([x]) = \left[\sum a_k b_{k, 1} \ldots b_{k, n_k}\right]
$$
映射 $s_0, s_1 : A[A[B]] \to A[A[A[B]]]$ 由
$$
s_0([x]) = \left[\left[\sum a_k [b_{k, 1}] \ldots [b_{k, n_k}]\right]\right]
\quad\text{且}\quad
s_1([x]) = \left[\sum a_k [[b_{k, 1}]] \ldots [[b_{k, n_k}]]\right]
$$
给出。
\end{example}

\begin{example}
\label{simplicial-example-polynomial-algebra-homotopy}
回到例 \ref{simplicial-example-polynomial-algebra} 中讨论的例子。
引理 \ref{simplicial-lemma-standard-simplicial-homotopy} 表明，对任意环映射
$A \to B$，单纯环映射
$$
\xymatrix{
A[A[A[B]]] \ar[d]
\ar@<2ex>[r]
\ar@<0ex>[r]
\ar@<-2ex>[r]
&
A[A[B]] \ar[d]
\ar@<1ex>[r]
\ar@<-1ex>[r]
\ar@<1ex>[l]
\ar@<-1ex>[l]
&
A[B] \ar[d]
\ar@<0ex>[l] \\
B
\ar@<2ex>[r]
\ar@<0ex>[r]
\ar@<-2ex>[r]
&
B
\ar@<1ex>[r]
\ar@<-1ex>[r]
\ar@<1ex>[l]
\ar@<-1ex>[l]
&
B
\ar@<0ex>[l]
}
$$
在底层单纯集上为同伦等价。此外，引理
\ref{simplicial-lemma-standard-simplicial-homotopy} 中构造的逆映射在次数
$n$ 上由
$$
b \longmapsto [\ldots[b]\ldots]
$$
给出，记号含义显然。反过来，该引理还说明，对每个集合 $E$，
存在环的同伦等价
$$
\xymatrix{
A[A[A[A[E]]]] \ar[d]
\ar@<2ex>[r]
\ar@<0ex>[r]
\ar@<-2ex>[r]
&
A[A[A[E]]] \ar[d]
\ar@<1ex>[r]
\ar@<-1ex>[r]
\ar@<1ex>[l]
\ar@<-1ex>[l]
&
A[A[E]] \ar[d]
\ar@<0ex>[l] \\
A[E]
\ar@<2ex>[r]
\ar@<0ex>[r]
\ar@<-2ex>[r]
&
A[E]
\ar@<1ex>[r]
\ar@<-1ex>[r]
\ar@<1ex>[l]
\ar@<-1ex>[l]
&
A[E]
\ar@<0ex>[l]
}
$$
。引理中构造的逆映射在次数 $n$ 上由环映射
$$
\sum a_{e_1,\ldots,e_p}[e_1][e_2] \ldots [e_p]
\longmapsto
\sum a_{e_1,\ldots,e_p}[\ldots[e_1]\ldots][\ldots[e_2]\ldots]
\ldots [\ldots[e_p]\ldots]
$$
给出（记号含义显然）。
\end{example}
